\documentclass[11pt,a4paper]{article}
\usepackage[UTF8]{ctex}
\usepackage{amsmath,amssymb,amsthm}
\usepackage{geometry}
\usepackage{booktabs}
\usepackage{array}
\newcolumntype{L}[1]{>{\raggedright\arraybackslash}p{#1}}  % 定宽列，左对齐不拉伸（中文窄列表格用）
\usepackage{enumitem}
\usepackage{hyperref}
\usepackage{titlesec}
\usepackage{fancyhdr}
\usepackage{tikz}
\usepackage{mathtools}

\geometry{margin=2.5cm}
\hypersetup{colorlinks=true,linkcolor=blue,urlcolor=blue}

\pagestyle{fancy}
\fancyhead[L]{现代数学物理与计算机图形学}
\fancyhead[R]{\thepage}
\fancyfoot[C]{}

\title{\textbf{现代数学物理与计算机图形学\\完整学习框架与笔记}}
\author{}
\date{}

\begin{document}
\maketitle
\tableofcontents
\newpage

\section{A. 分析学（Analysis）}

\subsection{A1. 从黎曼积分到勒贝格积分}

\subsubsection{黎曼积分的本质与局限}

黎曼积分的核心思想：将定义域 $[a,b]$ 分割为 $n$ 个小区间，在每个小区间上取函数值乘以区间长度求和，取极限：
\[
\int_a^b f(x)\,dx = \lim_{\max \Delta x_i \to 0} \sum_{i=1}^n f(\xi_i)\,\Delta x_i
\]

\textbf{局限性：}
\begin{itemize}
  \item 黎曼可积的充要条件：$f$ 有界且不连续点集合的勒贝格测度为零（即“几乎处处连续”）
  \item 狄利克雷函数（有理数处为 1，无理数处为 0）黎曼不可积
  \item 黎曼积分对极限运算不友好：$f_n \to f$ 逐点收敛时，$\int f_n$ 不一定收敛到 $\int f$
\end{itemize}

\subsubsection{勒贝格积分的构造}

核心思想转变：\textbf{不再分割定义域，而是分割值域。}

对于函数 $f: X \to \mathbb{R}$，考虑值域中的水平集：
\[
E_y = \{x : f(x) > y\}
\]

勒贝格积分定义为：
\[
\int f\,d\mu = \int_0^\infty \mu(\{x: f(x) > t\})\,dt
\]

构造步骤：
\begin{enumerate}
  \item \textbf{简单函数：} $s(x) = \sum_{i=1}^n a_i \mathbf{1}_{A_i}(x)$，其中 $A_i$ 是可测集
  \item 定义 $\int s\,d\mu = \sum a_i\,\mu(A_i)$
  \item 对非负可测函数 $f$：$\int f\,d\mu = \sup\left\{\int s\,d\mu : 0 \leq s \leq f,\; s \text{ 简单}\right\}$
  \item 一般函数：$f = f^+ - f^-$，分别积分
\end{enumerate}

\subsubsection{三大收敛定理}

\begin{center}
\begin{tabular}{L{3.5cm}L{5cm}L{4cm}}
\toprule
\textbf{定理} & \textbf{条件} & \textbf{结论} \\
\midrule
单调收敛定理 (MCT) & $0 \leq f_1 \leq f_2 \leq \cdots$，$f_n \to f$ & $\int f_n \to \int f$ \\
控制收敛定理 (DCT) & $f_n \to f$ a.e.，$\exists$ 可积 $g$ 使 $|f_n| \leq g$ & $\int f_n \to \int f$ \\
Fatou 引理 & $f_n \geq 0$ & $\int \liminf f_n \leq \liminf \int f_n$ \\
\bottomrule
\end{tabular}
\end{center}

在 CG 中：渲染方程的蒙特卡洛估计收敛性证明依赖 DCT。

\subsubsection{测度论基础}

\textbf{$\sigma$-代数：} 集合 $X$ 上的 $\sigma$-代数 $\mathcal{F}$ 是 $X$ 的子集族，满足：
\begin{itemize}
  \item $X \in \mathcal{F}$
  \item 对补集封闭：$A \in \mathcal{F} \Rightarrow A^c \in \mathcal{F}$
  \item 对可数并封闭：$A_1, A_2, \ldots \in \mathcal{F} \Rightarrow \bigcup A_i \in \mathcal{F}$
\end{itemize}

\textbf{为什么是“可数”并，而不是“任意”并？} 这是为了与下面测度的\textbf{可数可加}相匹配。
$\mathbb{R}$ 是可分的（存在可数拓扑基），每个开集都是可数个基开集之并，因此 Borel $\sigma$-代数
完全由\textbf{可数}层级的操作生成（超限层级 $\Sigma^0_\alpha$）。若改成对\emph{任意}并封闭，
$\mathcal{F}$ 立刻退化为幂集 $2^X$——而 Vitali 集证明：在幂集上无法同时拥有可数可加与平移不变的测度。
反过来说，$\sigma$-代数的“可数”二字，正是为可数可加腾出的逻辑空间。

\textbf{测度：} $\mu: \mathcal{F} \to [0, +\infty]$，满足：
\begin{itemize}
  \item $\mu(\emptyset) = 0$
  \item \textbf{可数可加性}：若 $\{A_n\}_{n=1}^\infty \subset \mathcal{F}$ 两两不交，则
  \[
  \mu\Big(\bigcup_{n=1}^\infty A_n\Big) = \sum_{n=1}^\infty \mu(A_n)
  \]
\end{itemize}

\textbf{可数可加的必要性——为什么恰好是“可数”？}\\
可数可加是\textbf{唯一同时满足两个要求的强度}：强到足以支撑分析，弱到不排除连续测度。

\textit{(1) 有限可加太弱：丢失全部极限工具。} 有限可加\emph{不}蕴含可数可加，有三个经典反例：

\begin{center}
\begin{tabular}{L{2.8cm}L{5.6cm}L{4.4cm}}
\toprule
\textbf{集函数} & \textbf{定义} & \textbf{失效方式} \\
\midrule
自然密度 & $\mu(A)=\lim\limits_{n}\frac{|A\cap\{1,\dots,n\}|}{n}$（极限存在时） &
$\mu(\{k\})=0$ 但 $\mu(\mathbb{N})=1$ \\
\addlinespace
退化测度 & $\mu(A)=0$（$A$ 有限），$\mu(A)=\infty$（$A$ 无限） &
$\mathbb{N}=\bigsqcup_k\{k\}$ 却 $\mu(\mathbb{N})=\infty \neq \sum_k 0$ \\
\addlinespace
Banach 测度 & $\mathbb{R}$ 上 Lebesgue 测度的有限可加扩张 &
三个条件全满足，但无任何收敛定理 \\
\bottomrule
\end{tabular}
\end{center}

Banach 测度说明：\emph{只}要求有限可加时，可以做得“漂亮”——但代价是\textbf{一整套定理同时失效}：

\begin{center}
\begin{tabular}{L{5.2cm}L{7.8cm}}
\toprule
\textbf{失去的东西} & \textbf{后果} \\
\midrule
MCT / Fatou / DCT & 积分与极限不能交换 $\to$ 无法做分析 \\
$L^p$ 完备性（Riesz--Fischer） & 泛函分析地基消失 \\
Fubini--Tonelli & 累次积分不成立 \\
Radon--Nikodym、Lebesgue 分解 & 无“密度函数”，条件期望无法定义（$\to$ C4） \\
Carathéodory 扩张的\textbf{唯一性} & 由区间长度构造 Lebesgue 测度时失去确定性 \\
Kolmogorov 相容性定理 & 无穷维概率空间（布朗运动）的存在性无法保证 \\
Borel--Cantelli、强大数定律、0--1 律 & 概率渐近理论全部失效 \\
\bottomrule
\end{tabular}
\end{center}

\textbf{对本笔记的直接影响：} H4 蒙特卡洛积分的收敛性证明依赖 DCT，而 DCT 经 MCT 依赖可数可加。
\textbf{有限可加下 MC 渲染没有理论保证}——这是 CG 侧最实际的理由。

\textit{(2) 不可数可加太强：只剩原子测度。} 若对\emph{任意}指标集可加，则对任意 $X$：
\[
\mu(X) = \sum_{x\in X}\mu(\{x\})
\]
而非负实数族的无序求和收敛\textbf{当且仅当至多可数项非零}。于是 $\mu(\{x\})=0\ (\forall x) \Rightarrow \mu \equiv 0$：
\begin{itemize}
  \item Lebesgue 测度满足 $\lambda(\{x\})=0$，将被迫 $\lambda(\mathbb{R})=0$——荒谬
  \item 正态分布等一切连续型分布被迫退化（C2 整节失效）
  \item A2.5 的分布理论、狄拉克 $\delta$ 建模（点光源）全部无法表述
\end{itemize}
即：不可数可加\textbf{只允许集中在至多可数个原子上的测度}，把连续世界排除在外。

\textit{(3) 可数可加的等价形式：对集合列的连续性。}
\[
A_n \uparrow A \implies \mu(A_n)\to\mu(A), \qquad
A_n \downarrow \varnothing \implies \mu(A_n)\to 0
\]
这是理解其必要性的钥匙——\textbf{分析学的一切建立在极限上，而有限可加对极限一无所知}。

\textit{(4) 代价：Vitali 三难。} 在选择公理下，
\[
\text{可数可加} \;+\; \text{平移不变} \;+\; \text{定义在 }\mathbb{R}\text{ 的全部子集上}
\]
三者\textbf{不可兼得}（Banach 测度靠放弃可数可加才同时做到）。因此我们放弃“全部子集”，
把定义域限制为 $\sigma$-代数，并做完备化。

\[
\boxed{\text{$\sigma$-代数这个看似繁琐的概念，其存在理由正是为可数可加腾出逻辑空间。}}
\]

\textbf{关键测度：}
\begin{itemize}
  \item 勒贝格测度：$\mathbb{R}^n$ 上的“体积”
  \item 计数测度：$\mu(A) = |A|$
  \item 狄拉克测度：$\delta_x(A) = \mathbf{1}_A(x)$（在 CG 中用于点光源建模）
  \item 豪斯多夫测度：用于分形维数（CG 中的分形地形）
\end{itemize}

\subsubsection{$L^p$ 空间}

\[
L^p(X, \mu) = \left\{f : \int |f|^p\,d\mu < \infty\right\} / \sim
\]

其中 $f \sim g$ 当且仅当 $f = g$ 几乎处处。

\begin{itemize}
  \item $L^1$：绝对可积函数（概率密度函数所在空间）
  \item $L^2$：平方可积函数（\textbf{量子力学的状态空间}，信号处理的核心）
  \item $L^\infty$：本质有界函数
\end{itemize}

\textbf{$L^2$ 是希尔伯特空间}，内积为 $\langle f, g \rangle = \int f\bar{g}\,d\mu$。这是连接分析学与量子力学的桥梁。

%--------------------------------------------------------------
\subsection{A2. 泛函分析}

\subsubsection{赋范空间与巴拿赫空间}

\begin{itemize}
  \item \textbf{赋范空间：} 向量空间 + 范数 $\|\cdot\|$（满足正定性、齐次性、三角不等式）
  \item \textbf{巴拿赫空间：} 完备的赋范空间（每个柯西序列都收敛）
  \item 例子：$\mathbb{R}^n$（有限维），$L^p$（无限维），$C[0,1]$（连续函数空间，sup 范数）
\end{itemize}

\subsubsection{希尔伯特空间}

\textbf{定义：} 完备的内积空间。

\textbf{核心性质：}
\begin{itemize}
  \item 内积诱导范数：$\|x\| = \sqrt{\langle x, x \rangle}$
  \item 柯西--施瓦茨不等式：$|\langle x, y \rangle| \leq \|x\|\,\|y\|$
  \item 平行四边形法则：$\|x+y\|^2 + \|x-y\|^2 = 2\|x\|^2 + 2\|y\|^2$
\end{itemize}

\textbf{正交分解定理：} 若 $M$ 是希尔伯特空间 $H$ 的闭子空间，则 $H = M \oplus M^\perp$。

这就是正交投影的严格基础——在 CG 中，最小二乘拟合、法线估计、球谐函数投影都是这个定理的应用。

\subsubsection{正交基与傅里叶分析}

可分希尔伯特空间存在可数正交基 $\{e_n\}$，任何元素可展开：
\[
f = \sum_{n} \langle f, e_n \rangle\, e_n
\]

\textbf{经典正交基：}
\begin{center}
\begin{tabular}{lll}
\toprule
\textbf{空间} & \textbf{正交基} & \textbf{应用} \\
\midrule
$L^2[0, 2\pi]$ & $\{e^{inx}\}$ & 傅里叶级数，信号处理 \\
$L^2(S^2)$ & 球谐函数 $Y_l^m$ & 环境光照、PRT \\
$L^2(\mathbb{R})$ & Hermite 函数 & 量子谐振子 \\
$L^2[0,1]$ & Legendre 多项式 & 数值积分 \\
$L^2(\mathbb{R}^3)$ & 平面波 $e^{i\mathbf{k}\cdot\mathbf{x}}$ & 衍射、散射 \\
\bottomrule
\end{tabular}
\end{center}

\textbf{Parseval 定理：} $\|f\|^2 = \sum |\langle f, e_n \rangle|^2$（能量守恒在频域的表述）

\subsubsection{有界线性算子与谱理论}
\label{subsec:spectral}

\textbf{算子：} $T: H \to H$，线性且有界（$\|Tx\| \leq C\|x\|$）。

\textbf{自伴算子（Hermitian）：} $\langle Tx, y \rangle = \langle x, Ty \rangle$

\textbf{谱定理（有限维）：} 自伴算子可以对角化，特征值全为实数，特征向量构成正交基。

\textbf{谱定理（无限维）：} 自伴算子的谱可能是连续的。例如位置算子 $\hat{x}$ 在 $L^2(\mathbb{R})$ 上的谱是整个实数轴。

\textbf{在 CG 中的应用：}
\begin{itemize}
  \item 拉普拉斯算子 $\Delta$ 的谱 $\to$ 网格上的振动模式 $\to$ 形状分析
  \item 协方差矩阵的特征分解 $\to$ PCA $\to$ 降维、法线估计
\end{itemize}

\textbf{微分算子的无穷维矩阵表示、矩阵指数与 Taylor 公式、若尔当标准型的推广，见~\ref{subsec:infdim-matrix}~节。}

\subsubsection{分布（广义函数）}

\textbf{动机：} 狄拉克 $\delta$ 函数不是普通函数，但在物理中无处不在。

\textbf{定义：} 分布是测试函数空间 $\mathcal{D} = C_c^\infty$ 上的连续线性泛函。
\[
\langle \delta, \varphi \rangle = \varphi(0)
\]

\textbf{关键分布：}
\begin{itemize}
  \item $\delta(x)$：点源（点电荷、点光源）
  \item $\delta'(x)$：偶极子
  \item $\text{p.v.}\frac{1}{x}$：主值分布
\end{itemize}

在 CG 中：点光源的辐射度用 $\delta$ 函数建模，BRDF 中的镜面反射分量包含 $\delta$ 函数。

%--------------------------------------------------------------
\subsection{A3. 微分学的高阶推广：微分形式}

\textbf{0-形式：} 标量函数 $f$ \\
\textbf{1-形式：} $\omega = P\,dx + Q\,dy + R\,dz$（对应向量场） \\
\textbf{2-形式：} $\omega = P\,dy\wedge dz + Q\,dz\wedge dx + R\,dx\wedge dy$（对应通量） \\
\textbf{3-形式：} $\omega = f\,dx\wedge dy\wedge dz$（对应体密度）

\textbf{外微分 $d$：}
\begin{itemize}
  \item $df = \frac{\partial f}{\partial x}dx + \frac{\partial f}{\partial y}dy + \frac{\partial f}{\partial z}dz$（梯度）
  \item $d(\text{1-形式})$ 对应旋度
  \item $d(\text{2-形式})$ 对应散度
  \item $d^2 = 0$（旋度的散度为零，梯度的旋度为零）
\end{itemize}

\textbf{广义 Stokes 定理：}
\[
\boxed{\int_{\partial M} \omega = \int_M d\omega}
\]

这一个公式统一了：
\begin{itemize}
  \item 微积分基本定理（$n=1$）
  \item 格林公式（$n=2$）
  \item 经典 Stokes 定理（$n=2$，三维中）
  \item 散度定理/高斯定理（$n=3$）
\end{itemize}

在 CG 中：离散外微分（DEC）用于网格上的物理模拟、守恒律的离散保持、网格参数化。

%--------------------------------------------------------------
\subsection{A4. Sobolev 空间与弱解}
\label{subsec:sobolev}

\textbf{为什么单列一节：} A1 的 $L^p$、A2 的 Hilbert 空间、P9 的能量法，三条线在这里汇合，
产出的东西是\textbf{整个椭圆 PDE 存在性理论的骨架}。没有它，「NS 方程有解吗」这句话根本没法说清。

\subsubsection{经典导数的不足}

$a$ 处的导数要求 $u$ 在 $a$ 的邻域内有定义且极限存在。可是：

\begin{itemize}
  \item $u(x)=|x|$ 在 $0$ 处无经典导数，但它是很好的函数，其「导数」应当是符号函数；
  \item 激波、间断初值、Leray--Hopf 弱解——这些\textbf{物理上有意义的解根本不光滑}；
  \item 极限过程会破坏光滑性：$C^1$ 中的序列可以在 $L^p$ 中收敛到一个不 $C^1$ 的函数。
\end{itemize}

\textbf{出路：} 用分部积分把导数「搬」到测试函数上——这正是 P9「移导数」那一手，
只不过现在把它\textbf{当作定义}。

\subsubsection{弱导数}

\textbf{定义：} 设 $u\in L^1_{\text{loc}}(\Omega)$。若存在 $v\in L^1_{\text{loc}}(\Omega)$，使对一切
$\varphi\in C_c^\infty(\Omega)$（紧支光滑测试函数）成立
\[
\int_\Omega u\,\partial_i\varphi\,dx=-\int_\Omega v\,\varphi\,dx,
\]
则称 $v$ 是 $u$ 关于 $x_i$ 的\textbf{弱导数}，记 $v=\partial_i u$。

\textbf{直觉：} 经典分部积分 $\int u\varphi'=-\int u'\varphi$（边界项因 $\varphi$ 紧支而消失）本来是
\emph{由 $u'$ 推出}的恒等式；现在反过来，\textbf{用它来定义 $u'$}。$u$ 只需要可积，不需要连续。

弱导数若存在则\textbf{几乎处处唯一}，且与经典导数在后者存在时一致。

\subsubsection{Sobolev 空间}

\[
W^{k,p}(\Omega)=\Big\{u\in L^p(\Omega):\ \partial^\alpha u\in L^p(\Omega)\ \text{对一切}\ |\alpha|\le k\Big\}
\]
（$\alpha$ 是多重指标，$\partial^\alpha$ 按弱导数理解）配备范数
\[
\|u\|_{W^{k,p}}=\Big(\sum_{|\alpha|\le k}\|\partial^\alpha u\|_{L^p}^p\Big)^{1/p}
\qquad(p<\infty)
\]

记 $H^k(\Omega):=W^{k,2}(\Omega)$。$H_0^k(\Omega)$ 是 $C_c^\infty(\Omega)$ 在 $H^k$ 范数下的闭包——
\textbf{「零边界条件」的精确含义}。

\begin{center}
\begin{tabular}{L{2.4cm}L{5.6cm}L{6.2cm}}
\toprule
\textbf{空间} & \textbf{性质} & \textbf{要点} \\
\midrule
$W^{k,p}$ & 完备（Banach） & 所以能取极限、能套压缩映射、能谈存在唯一 \\
$H^k=W^{k,2}$ & Hilbert（内积 $\sum_{|\alpha|\le k}\langle\partial^\alpha u,\partial^\alpha v\rangle$） &
可直接用 A2 的正交分解与投影，$H_0^1$ 更是 \\
磨光（mollification） & $C_c^\infty(\Omega)$ 在 $W^{k,p}(\Omega)$ 中稠密（$p<\infty$） &
\textbf{光滑函数够用}：先在光滑里证，再取极限 \\
\bottomrule
\end{tabular}
\end{center}

\subsubsection{嵌入定理}

\textbf{Sobolev 嵌入：} 设 $\Omega\subset\mathbb R^n$ 有界、边界「足够好」，$1\le p<n$，则
\[
W^{1,p}(\Omega)\hookrightarrow L^{p^*}(\Omega),\qquad
\frac1{p^*}=\frac1p-\frac1n
\]
（$p>n$ 时 $W^{1,p}\hookrightarrow C^{0,\alpha}$，$\alpha=1-n/p$——变成了连续函数。）

\textbf{怎么读这个公式：} 多「一个」导数的可积性，换一个「更好」的空间；
\textbf{损失的额度由维数 $n$ 决定}。$n=1$ 时 $p^*=\infty$：一维情形 $W^{1,p}\hookrightarrow C$，
所以一维「光滑性」几乎是免费的——这也是为什么一维 ODE 理论比 PDE 简单得多。

\textbf{紧嵌入（Rellich--Kondrachov）：} 有界 $\Omega$ 上，$p<n$ 时
\[
W^{1,p}(\Omega)\hookrightarrow\hookrightarrow L^q(\Omega)\qquad(\textstyle 1\le q<p^*)
\]
\textbf{「紧」是这里的全部重点：} 有界集映成相对紧集 $\Longrightarrow$ 可取收敛子列
（见 Q5：这就是变分法直接方法要的那一步）。$q<p^*$ 的严格不等号不能改成等号——
临界指数处嵌入恰好不紧，这正是\textbf{临界指数问题}难的原因。

\textbf{Poincaré 不等式：}
\[
\|u-u_\Omega\|_{L^p}\le C\,\|\nabla u\|_{L^p}\qquad
\text{特别地}\quad \|u\|_{L^p}\le C\|\nabla u\|_{L^p}\ \ (u\in W_0^{1,p})
\]
后一式说明在 $H_0^1$ 上 $\|\nabla u\|_{L^2}$ \textbf{本身就是等价范数}——
它是下面「强制性」的来源。

\subsubsection{弱解与 Lax--Milgram}

以 Poisson 方程为例：$-\Delta u=f$，$u|_{\partial\Omega}=0$。

\textbf{弱形式：} 两边乘 $v\in H_0^1$、积分、对左边分部积分（P9 的移导数），求 $u\in H_0^1$ 使
\[
\underbrace{\int_\Omega\nabla u\cdot\nabla v\,dx}_{B(u,v)}=\underbrace{\int_\Omega fv\,dx}_{F(v)}
\qquad\forall v\in H_0^1
\]

\textbf{Lax--Milgram 定理：} $H$ 是 Hilbert 空间，$B:H\times H\to\mathbb R$ 双线性。若
\begin{itemize}
  \item \textbf{有界：} $|B(u,v)|\le C\|u\|\|v\|$，
  \item \textbf{强制（coercive）：} $B(u,u)\ge\alpha\|u\|^2$（$\alpha>0$），
\end{itemize}
且 $F$ 是连续线性泛函，则存在\textbf{唯一} $u\in H$ 使 $B(u,v)=F(v)$ 对一切 $v$ 成立。

\textbf{验证只需两步（这就是它的价值）：}
\begin{itemize}
  \item 有界：$|B(u,v)|\le\|\nabla u\|\|\nabla v\|$（直接 Cauchy--Schwarz，P3 的第一件套）；
  \item 强制：$B(u,u)=\|\nabla u\|^2\ge\alpha\|u\|_{H^1}^2$——\textbf{正是 Poincaré 不等式}。
\end{itemize}

\textbf{方法论上的意义：} 它把「解一个无穷维方程」变成「验证两个不等式」。
弱解存在唯一性一旦拿到，剩下的就是\textbf{正则性}问题（把弱解升级成经典解）。

\subsubsection{正则性提升与能量估计}

\textbf{能量估计（先验估计）：} 在弱形式中取 $v=u$，立刻得
\[
\|\nabla u\|_{L^2}^2=\int f u\le\|f\|_{L^2}\|u\|_{L^2}\le C\|f\|_{L^2}\|\nabla u\|_{L^2}
\ \Longrightarrow\ \|u\|_{H^1}\le C\|f\|_{L^2}
\]
（中间一步用 Cauchy--Schwarz，最后一步用 Poincaré。）\textbf{这不是「顺便得到」的结论——
PDE 理论中几乎所有存在性都是先拿到先验估计，再用紧性取极限。}

\textbf{内部正则性：} 若 $f\in H^k$ 且 $\partial\Omega$ 光滑，则 $u\in H^{k+2}$，
$k$ 足够大时由 Sobolev 嵌入即得 $u\in C^{2,\alpha}$——\textbf{弱解自动变成经典解}。

\textbf{这条链要记住：}
\[
\text{弱形式}\to\text{Lax--Milgram 存在唯一}\to\text{能量估计}\to\text{正则性提升}\to\text{经典解}
\]

\textbf{连到 NS：} 同一条链在 Navier--Stokes 上只能走到第二步——
Leray--Hopf 弱解的存在性（1934）靠的正是 P9 的能量法流水线（乘 $u$、移导数、Gronwall 收尾）。
\textbf{但正则性与唯一性至今未解决}（千禧年问题）。所以 PDE 里「弱解」不是权宜之计，
而是\textbf{非线性和临界情形下唯一能拿到的东西}。

\textbf{与既有章节的接口：}
\begin{itemize}
  \item $L^p$ 与收敛定理见 A1；Hilbert 空间与投影见 A2；
  \item 移导数与能量法五步见 P9；紧嵌入为什么关键见 Q5；Rellich--Kondrachov 的拓扑背景见 Q5 末；
  \item 分布（广义函数）是弱导数的更一般版本，见 A2.5。
\end{itemize}
\section{B. 代数与几何}

\subsection{B1. 从矩阵到群论}

\subsubsection{群的基本概念}

\textbf{群 $(G, \cdot)$：} 集合 + 运算，满足封闭性、结合律、单位元、逆元。

\textbf{关键群：}
\begin{center}
\begin{tabular}{llll}
\toprule
\textbf{群} & \textbf{定义} & \textbf{维度} & \textbf{应用} \\
\midrule
$SO(2)$ & 2D 旋转 & 1 & 2D 旋转 \\
$SO(3)$ & 3D 旋转 & 3 & 刚体旋转 \\
$SE(3)$ & 3D 刚体运动 & 6 & 机器人学、动画 \\
$U(1)$ & 复数单位圆 & 1 & 电磁规范对称 \\
$SU(2)$ & $2\times2$ 酉矩阵，$\det=1$ & 3 & 自旋、四元数 \\
$SU(3)$ & $3\times3$ 酉矩阵，$\det=1$ & 8 & 量子色动力学 \\
$O(3)$ & 正交矩阵 & 3 & 包含反射 \\
$GL(n)$ & 可逆矩阵 & $n^2$ & 一般线性变换 \\
$SL(n)$ & $\det=1$ 的矩阵 & $n^2-1$ & 保体积变换 \\
\bottomrule
\end{tabular}
\end{center}

\subsubsection{李群与李代数}
\label{subsec:lie-intro}

\textbf{李群：} 既是群又是光滑流形（群运算光滑）。

\textbf{李代数 $\mathfrak{g}$：} 李群在单位元处的切空间，配备李括号 $[\cdot, \cdot]$。

\textbf{指数映射：} $\exp: \mathfrak{g} \to G$

$\mathfrak{so}(3)$ 的具体形式：
\[
\mathfrak{so}(3) = \{A \in \mathbb{R}^{3\times 3} : A^T = -A\}
\]

反对称矩阵参数化：
\[
[\omega]_\times = \begin{pmatrix} 0 & -\omega_3 & \omega_2 \\ \omega_3 & 0 & -\omega_1 \\ -\omega_2 & \omega_1 & 0 \end{pmatrix}
\]

Rodrigues 公式：
\[
R = \exp([\omega]_\times) = I + \frac{\sin\theta}{\theta}[\omega]_\times + \frac{1-\cos\theta}{\theta^2}[\omega]_\times^2
\]
其中 $\theta = \|\omega\|$。

\textbf{$SU(2)$ 与四元数：} 四元数 $q = a + bi + cj + dk$，单位四元数构成 $SU(2)$（同构）。$SU(2) \to SO(3)$ 是 2:1 的覆盖映射（旋转 $2\pi$ 后四元数变号，需要 $4\pi$ 才回到自身——自旋 $1/2$ 粒子的数学根源）。

\textbf{经典李群的完整分类见第~\ref{sec:lie-group}~节；李代数结构常数、Killing 形式、Dynkin 图与表示论见第~\ref{sec:lie-algebra}~节及其后。}

\subsubsection{表示论基础}

\textbf{表示：} 群 $G$ 到 $GL(V)$ 的同态 $\rho: G \to GL(V)$。

\textbf{不可约表示：} 没有非平凡不变子空间的表示。

\textbf{在物理中：}
\begin{itemize}
  \item 角动量的量子化来自 $SO(3)$ 的不可约表示（维数 $2l+1$）
  \item 粒子分类来自庞加莱群的不可约表示
  \item 球谐函数 $Y_l^m$ 就是 $SO(3)$ 的第 $l$ 个不可约表示的基函数
\end{itemize}

\textbf{在 CG 中：}
\begin{itemize}
  \item 球谐光照（Spherical Harmonics Lighting）
  \item 等变神经网络（Equivariant Neural Networks）
\end{itemize}

%--------------------------------------------------------------
\subsection{B2. 微分几何}

\subsubsection{曲线与曲面}

\textbf{曲线} $\gamma: [a,b] \to \mathbb{R}^3$：
\begin{itemize}
  \item 弧长参数化：$s = \int_0^t \|\gamma'(\tau)\|\,d\tau$
  \item 曲率：$\kappa = \|\mathbf{T}'(s)\|$
  \item 挠率：$\tau$
  \item Frenet 标架：$(\mathbf{T}, \mathbf{N}, \mathbf{B})$
\end{itemize}

\textbf{曲面} $\mathbf{r}(u,v)$：
\begin{itemize}
  \item 第一基本形式：$I = E\,du^2 + 2F\,du\,dv + G\,dv^2$（度量）
  \item 第二基本形式：$II = L\,du^2 + 2M\,du\,dv + N\,dv^2$（弯曲）
  \item 高斯曲率：$K = \dfrac{LN - M^2}{EG - F^2}$（\textbf{内蕴量}，不依赖嵌入）
  \item 平均曲率：$H = \dfrac{EN - 2FM + GL}{2(EG - F^2)}$
\end{itemize}

\textbf{高斯绝妙定理（Theorema Egregium）：} 高斯曲率是内蕴的，只依赖第一基本形式。

在 CG 中：网格曲率估计、平均曲率流（网格平滑）、高斯曲率（展开/参数化）。

\subsubsection{流形与切空间}

\begin{itemize}
  \item \textbf{流形 $M$：} 局部像 $\mathbb{R}^n$ 的拓扑空间
  \item \textbf{切空间 $T_pM$：} 点 $p$ 处所有“方向导数”构成的向量空间
  \item \textbf{切丛 $TM$：} 所有切空间的并集
  \item \textbf{余切空间 $T_p^*M$：} 切空间的对偶（1-形式生活的地方）
\end{itemize}

\subsubsection{黎曼几何}

\textbf{黎曼度量 $g$：} 每个切空间上的内积，光滑地依赖于点。
\[
g = g_{ij}\,dx^i \otimes dx^j
\]

\textbf{Christoffel 符号：}
\[
\Gamma^k_{ij} = \frac{1}{2}g^{kl}\left(\frac{\partial g_{il}}{\partial x^j} + \frac{\partial g_{jl}}{\partial x^i} - \frac{\partial g_{ij}}{\partial x^l}\right)
\]

\textbf{黎曼曲率张量：}
\[
R^l_{ijk} = \partial_j \Gamma^l_{ik} - \partial_k \Gamma^l_{ij} + \Gamma^m_{ik}\Gamma^l_{mj} - \Gamma^m_{ij}\Gamma^l_{mk}
\]

\textbf{测地线方程：}
\[
\frac{d^2 x^k}{dt^2} + \Gamma^k_{ij}\frac{dx^i}{dt}\frac{dx^j}{dt} = 0
\]

在 CG 中：网格上的测地线计算、测地线距离（形状匹配）、黎曼优化。

\subsubsection{纤维丛（Fiber Bundle）}

\textbf{直觉：} 一个空间，每一点上面“粘着”一个纤维（向量空间或其他结构）。

\textbf{例子：}
\begin{itemize}
  \item 切丛 $TM$：每一点的纤维是切空间
  \item 余切丛 $T^*M$：每一点的纤维是余切空间
  \item 主丛：纤维是群 $G$（规范场论的数学语言）
\end{itemize}

\textbf{在物理中：} 规范场（电磁场、Yang--Mills 场）是主丛上的联络。
\section{C. 概率与测度}

\subsection{C1. 概率的测度论基础}

\textbf{概率空间} $(\Omega, \mathcal{F}, P)$：
\begin{itemize}
  \item $\Omega$：样本空间
  \item $\mathcal{F}$：$\sigma$-代数（事件集合）
  \item $P$：概率测度，$P(\Omega) = 1$
\end{itemize}

\textbf{随机变量：} 可测函数 $X: (\Omega, \mathcal{F}) \to (\mathbb{R}, \mathcal{B}(\mathbb{R}))$

\textbf{分布函数：} $F_X(x) = P(X \leq x)$

\textbf{密度函数（如果存在）：} $f_X(x) = F_X'(x)$，$P(X \in A) = \int_A f_X\,dx$

\subsection{C2. 重要分布族}

\begin{center}
\begin{tabular}{L{3cm}L{5cm}L{2cm}L{3.5cm}}
\toprule
\textbf{分布} & \textbf{密度/质量函数} & \textbf{参数} & \textbf{应用} \\
\midrule
均匀 $U(a,b)$ & $\frac{1}{b-a}$ & $a, b$ & 随机采样基础 \\
正态 $N(\mu, \sigma^2)$ & $\frac{1}{\sqrt{2\pi}\sigma}e^{-\frac{(x-\mu)^2}{2\sigma^2}}$ & $\mu, \sigma$ & 中心极限定理、噪声 \\
指数 $\text{Exp}(\lambda)$ & $\lambda e^{-\lambda x}$ & $\lambda$ & 光子自由程 \\
泊松 $\text{Pois}(\lambda)$ & $\frac{\lambda^k e^{-\lambda}}{k!}$ & $\lambda$ & 光子计数、散粒噪声 \\
冯·米塞斯--费歇尔 & $C(\kappa)e^{\kappa \mu \cdot x}$（球面上） & $\kappa, \mu$ & 方向统计 \\
GGX & $\frac{\alpha^2}{\pi(\cos^2\theta(\alpha^2-1)+1)^2}$ & $\alpha$ & 微表面 BRDF \\
Beta & $\frac{x^{\alpha-1}(1-x)^{\beta-1}}{B(\alpha,\beta)}$ & $\alpha, \beta$ & 共轭先验 \\
\bottomrule
\end{tabular}
\end{center}

\subsection{C3. 多维分布与独立性}

$X, Y$ 独立 $\iff$ $f_{X,Y}(x,y) = f_X(x)\,f_Y(y)$ $\iff$ $F_{X,Y}(x,y) = F_X(x)\,F_Y(y)$

\textbf{换元与独立性：} 设 $(U,V) = g(X,Y)$，则：
\[
f_{U,V}(u,v) = f_{X,Y}(g^{-1}(u,v)) \cdot |J|
\]
其中 $J = \det\frac{\partial(x,y)}{\partial(u,v)}$ 是雅可比行列式。

\textbf{关键洞察：} 独立性不是变换不变的。极坐标变换 $(x,y) \to (r,\theta)$ 中，如果 $(X,Y)$ 是标准正态，则 $R$ 和 $\Theta$ 独立；但一般情况下换元后独立性可能被破坏或创造。

在 CG 中：球面采样需要 $\theta = \arccos(1-2u_1)$，$\phi = 2\pi u_2$；雅可比行列式在路径追踪的变量替换中至关重要。

\subsection{C4. 条件期望与鞅}

\textbf{条件期望：} $E[X|\mathcal{G}]$ 是关于子 $\sigma$-代数 $\mathcal{G}$ 可测的随机变量，满足：
\[
\int_A E[X|\mathcal{G}]\,dP = \int_A X\,dP, \quad \forall A \in \mathcal{G}
\]

\textbf{鞅：} $E[X_{n+1}|\mathcal{F}_n] = X_n$

在 CG 中：俄罗斯轮盘（Russian Roulette）和分裂（Splitting）的无偏性证明、路径追踪中估计量的方差分析。

\subsection{C5. 随机过程}

\textbf{布朗运动 $W_t$：}
\begin{itemize}
  \item $W_0 = 0$
  \item 独立增量
  \item $W_t - W_s \sim N(0, t-s)$
  \item 路径连续但处处不可微
\end{itemize}

\textbf{随机微分方程（SDE）：}
\[
dX_t = \mu(X_t, t)\,dt + \sigma(X_t, t)\,dW_t
\]

在 CG 中：布朗运动（粒子模拟、烟雾）、朗之万方程（分子动力学）。更完整的 Itô 微积分见 C10。

%--------------------------------------------------------------
\subsection{C6. 收敛模式与极限定理}
\label{subsec:convergence-modes}

这是 C1--C5 之后最该补的一块：前面都在讲「分布是什么」，这里讲「一列随机变量怎么趋近另一个」。

\subsubsection{四种收敛模式}

\begin{center}
\begin{tabular}{L{2.6cm}L{5.6cm}L{6.4cm}}
\toprule
\textbf{模式} & \textbf{定义} & \textbf{一句话} \\
\midrule
几乎必然（a.s.） & $\mathbb P(X_n\to X)=1$ & 轨道意义上逐点收敛 \\
依概率（$\xrightarrow{\ \mathbb P\ }$） & $\forall\varepsilon>0:\ \mathbb P(|X_n-X|>\varepsilon)\to0$ &
大偏差事件的概率消失 \\
$L^p$（$\xrightarrow{\ L^p\ }$） & $\mathbb E|X_n-X|^p\to0$ & 平均意义（$p=1$ 即 $L^1$） \\
依分布（$\Rightarrow$） & $F_n(x)\to F(x)$ 在 $F$ 的连续点 & 只看分布形状，不管耦合方式 \\
\bottomrule
\end{tabular}
\end{center}

\textbf{层级关系}（唯一需要背的一张图）：
\[
\text{a.s.}\ \Longrightarrow\ \text{依概率}\ \Longrightarrow\ \text{依分布},
\qquad
L^p\ \Longrightarrow\ \text{依概率}
\]
\textbf{四个箭头都不可逆。} 反向需要一个额外条件：均匀可积性（C6.3）或常数列（若 $X$ 是常数，
则依概率收敛与依分布收敛等价）。

\subsubsection{反例博物馆：为什么箭头不可逆}

\begin{center}
\begin{tabular}{L{3.4cm}L{5.2cm}L{5.6cm}}
\toprule
\textbf{要否定的箭头} & \textbf{反例} & \textbf{机理} \\
\midrule
依概率 $\not\Rightarrow$ a.s. & 区间 $[0,1]$ 上的「打字机」序列
$X_n=\mathbf 1_{[k/2^m,(k+1)/2^m]}$（把 $[0,1]$ 按 $2^m$ 等分后逐个扫过） &
每点被命中无穷多次，故 a.s. 不收敛；但 $\mathbb P(X_n=1)=2^{-m}\to0$ \\
\addlinespace
a.s. $\not\Rightarrow L^1$ & $X_n=n\cdot\mathbf 1_{(0,1/n)}$ &
a.s. 收敛到 $0$，但 $\mathbb E X_n=1\not\to0$：质量「挤进零测点」 \\
\addlinespace
依分布 $\not\Rightarrow$ 依概率 & 取 $X_n$ 与 $X$ 独立同分布（如都 $\sim N(0,1)$） &
$F_n=F$ 故依分布收敛；但 $X_n-X$ 的方差恒为 $2$ \\
\bottomrule
\end{tabular}
\end{center}

\textbf{注意第二个反例}——它和 A1 里 DCT 那个「盖子」是同一件事的两面：a.s. 收敛之所以推不出 $L^1$ 收敛，
就是因为极限过程的「质量」可以逃到无穷远或挤进零测集。$X_n=n\mathbf 1_{(0,1/n)}$ 正是教科书里的标准反例。

\subsubsection{一致可积性：让「依概率」升到 $L^1$}

\textbf{定义：} 族 $\{X_n\}$ 称为\textbf{均匀可积}（uniformly integrable, UI），若
\[
\lim_{M\to\infty}\sup_n \mathbb E\big[|X_n|\,\mathbf 1_{\{|X_n|>M\}}\big]=0
\]
即「大值部分对期望的贡献可以被一致地压小」。

\textbf{核心定理（Vitali 收敛定理）：} 若 $\mathbb E|X_n|<\infty$，$X_n\xrightarrow{\ \mathbb P\ }X$，且 $\{X_n\}$ 均匀可积，则
$X_n\xrightarrow{\ L^1\ }X$ 且 $\mathbb E X_n\to\mathbb E X$。

\textbf{为什么这等于给 DCT 换了一副面孔：} DCT 用「被一个可积函数盖住」来保证 $L^1$ 收敛；
均匀可积是它的抽象化——$|X_n|\le g\in L^1$ 本身就是均匀可积性的一个充分条件。
\textbf{一句话：DCT 的盖子 $\iff$ 均匀可积性。} 这正是 A1 中「三大收敛定理」那条线索的收口。

\subsubsection{三大定理之外的收敛定理}

\begin{center}
\begin{tabular}{L{3.0cm}L{5.0cm}L{6.2cm}}
\toprule
\textbf{定理} & \textbf{条件} & \textbf{结论} \\
\midrule
Egorov & 有限测度，$f_n\to f$ a.e. & 对任意 $\varepsilon>0$，去掉一个测度 $<\varepsilon$ 的集合后\textbf{一致}收敛
（a.e. 收敛 $\approx$ 一致收敛，代价是丢掉一小块）\\
\addlinespace
Scheffé & $f_n,f\ge0$ 可积，$f_n\to f$ a.e.，且 $\int f_n\to\int f$ & $\int|f_n-f|\to0$（即 $L^1$ 收敛） \\
\addlinespace
Vitali 收敛 & $\mu$ 有限，$f_n\to f$ a.e.，$\{f_n\}$ 均匀可积 & $\int|f_n-f|\to0$ \\
\addlinespace
单调类定理（Dynkin $\pi$--$\lambda$） & $\mathcal P$ 是 $\pi$-系，$\mathcal L$ 是 $\lambda$-系，$\mathcal P\subset\mathcal L$ &
$\sigma(\mathcal P)\subset\mathcal L$ \\
\bottomrule
\end{tabular}
\end{center}

\textbf{单调类定理为什么重要：} 它把「证明两个测度相等」从「对一切可测集验证」压缩成「只在一个生成 $\pi$-系上验证」。
典型用法：$\mu,\nu$ 是 $(\Omega,\sigma(\mathcal P))$ 上的概率测度，若 $\mu|_{\mathcal P}=\nu|_{\mathcal P}$ 且 $\mathcal P$ 是 $\pi$-系，则 $\mu=\nu$。
这是测度论里最省力的证明手段之一。

\subsubsection{Borel--Cantelli 与两种大数定律}

\textbf{Borel--Cantelli 引理：}
\begin{itemize}
  \item 第一引理（不需要独立性）：$\sum_n\mathbb P(A_n)<\infty\ \Longrightarrow\ \mathbb P(\limsup_n A_n)=0$，
        即「几乎必然只有有限多个 $A_n$ 发生」。
  \item 第二引理（需要独立性）：$A_n$ 相互独立且 $\sum_n\mathbb P(A_n)=\infty\ \Longrightarrow\ \mathbb P(\limsup_n A_n)=1$。
\end{itemize}
（$\limsup_n A_n=\bigcap_N\bigcup_{n\ge N}A_n$ 表示「无穷多个 $A_n$ 发生」，即 $A_n$ i.o. 事件。）

\textbf{大数定律：} 设 $X_n$ 独立同分布，$\mathbb E X_1=\mu$（有限）。
\begin{itemize}
  \item \textbf{弱大数定律}（Khinchin）：$\bar X_n\xrightarrow{\ \mathbb P\ }\mu$，不需要方差有限。
  \item \textbf{强大数定律}（Kolmogorov）：$\bar X_n\to\mu$ a.s.（需要 $\mathbb E|X_1|<\infty$ 且独立同分布）。
\end{itemize}

\textbf{在 CG 中的意义：} 蒙特卡洛估计量的一致性正是（弱）大数定律——$\bar X_n\to I$ 依概率。
没有它，$L^1$ 收敛 $\Rightarrow$ 依概率收敛这条链路断了，「多采样就能逼近」就没有理论保证（见 C9）。

%--------------------------------------------------------------
\subsection{C7. 特征函数与中心极限定理}

\subsubsection{定义与基本性质}

\[
\varphi_X(t)=\mathbb E\big[e^{itX}\big]=\int_{\mathbb R}e^{itx}\,dF_X(x)
\]

特征函数就是分布的（反向）Fourier 变换——所以 P5 的 Euler 公式与 A2 的 Fourier 分析在这里合流。

\begin{itemize}
  \item $\varphi_{a+bX}(t)=e^{iat}\varphi_X(bt)$；
  \item $X,Y$ 独立 $\Longrightarrow$ $\varphi_{X+Y}=\varphi_X\varphi_Y$（\textbf{卷积变乘法}，这是它最有用的一条）；
  \item $|\varphi_X(t)|\le1$，$\varphi_X(0)=1$，$\varphi_X$ 在 $\mathbb R$ 上一致连续；
  \item $\mathbb E X^k=(-i)^k\varphi_X^{(k)}(0)$（矩由原点处的导数读出）；
  \item \textbf{一一对应：} $\varphi_X$ 唯一决定 $F_X$（反演公式），且两分布弱收敛 $\iff$ $\varphi_n\to\varphi$（局部一致）。
\end{itemize}

\subsubsection{算例：正态分布}

对 $X\sim N(0,1)$，由 P8 的高斯积分：
\[
\varphi_X(t)=\frac{1}{\sqrt{2\pi}}\int e^{itx}e^{-x^2/2}dx
=\frac{1}{\sqrt{2\pi}}\int e^{-(x-it)^2/2}dx\cdot e^{-t^2/2}=e^{-t^2/2}
\]
（配方后把被积函数补成完整的高斯积分 $=\sqrt{2\pi}$，余下的 $e^{-t^2/2}$ 就是答案。）

\subsubsection{中心极限定理（CLT）}

\textbf{定理（Lévy--Lindeberg）：} $X_n$ 独立同分布，$\mathbb E X_1=\mu$，$\operatorname{Var}X_1=\sigma^2\in(0,\infty)$，则
\[
\sqrt n\,\frac{\bar X_n-\mu}{\sigma}\ \xrightarrow{\ d\ }\ N(0,1)
\]

\textbf{证明路线（三步，必须会复述）：}
\begin{enumerate}
  \item 令 $Y_i=(X_i-\mu)/\sigma$，标准化后 $\varphi_Y(0)=1$、$\varphi_Y'(0)=0$、$\varphi_Y''(0)=-1$；
  \item 把 $\varphi_Y$ 在 $0$ 处 Taylor 展开到二阶：$\varphi_Y(t)=1-\dfrac{t^2}{2}+o(t^2)$
        （这一步用到 $e^{iu}=1+iu-\frac{u^2}{2}+O(u^3)$，即 P7 的 Maclaurin 表）；
  \item 独立性使和的特征函数相乘：
        \[
        \varphi_{S_n/\sqrt n}(t)=\Big[\varphi_Y\Big(\frac{t}{\sqrt n}\Big)\Big]^n
        =\Big[1-\frac{t^2}{2n}+o\Big(\frac1n\Big)\Big]^n\ \longrightarrow\ e^{-t^2/2}.
        \]
        右边正是 $N(0,1)$ 的特征函数，由 C7.1 的连续性定理即得依分布收敛。
\end{enumerate}

\textbf{直觉：} 每加一个独立变量，在特征函数上就是乘一次「略小于 $1$ 的因子」；
乘 $n$ 次的累积效果恰好复现 Gaussian 的 $e^{-t^2/2}$——这就是「为什么到处是指数」。
（$n\to\infty$ 那步本质上是 P4 的 $\big(1+\frac xn\big)^n\to e^x$。）

\textbf{在 CG 中的意义：} 蒙特卡洛估计量的误差估计
$\big|\bar X_n-I\big|\lesssim \sigma/\sqrt n$ 与置信区间全部来自 CLT——
\textbf{这就是路径追踪「采样数翻四倍、噪声减半」的出处}。

%--------------------------------------------------------------
\subsection{C8. 马尔可夫链与 MCMC}

\subsubsection{基本对象}

\textbf{马尔可夫链：} $\mathbb P(X_{n+1}\in A\mid X_0,\dots,X_n)=\mathbb P(X_{n+1}\in A\mid X_n)$。
离散状态时由\textbf{转移矩阵} $P=(p_{ij})$ 描述。

\begin{itemize}
  \item \textbf{Chapman--Kolmogorov：} $P^{(m+n)}=P^{(m)}P^{(n)}$（矩阵幂 $=$ 多步转移）。
  \item \textbf{平稳分布：} $\pi$ 满足 $\pi P=\pi$，$\pi_i\ge0$，$\sum_i\pi_i=1$。
  \item \textbf{细致平衡（可逆性）：} $\pi_ip_{ij}=\pi_jp_{ji}$ $\Longrightarrow$ $\pi P=\pi$
        （对 $j$ 求和即得）。\textbf{这是 MCMC 的构造原则}——不去解 $\pi P=\pi$，而是直接造一个满足细致平衡的 $P$。
\end{itemize}

\textbf{遍历定理：} 若链不可约、非周期且正常返，则转移矩阵 $P^{(n)}$ 的每一行都收敛到平稳分布 $\pi$，
且时间平均几乎必然收敛到空间平均：
\[
\frac1n\sum_{k=1}^n f(X_k)\ \longrightarrow\ \sum_i\pi_i f(i)\qquad(\text{a.s.})
\]
（遍历定理是强大数定律在相依序列上的版本——C6 的独立性在这里被换成了马尔可夫性。）

\subsubsection{Metropolis--Hastings 与 MLT}

\textbf{Metropolis--Hastings：} 想从目标分布 $\pi$（可能只知道到差一个常数因子的密度）采样时：

\begin{enumerate}
  \item 从提议分布 $q(\cdot\mid x)$ 抽 $y$；
  \item 以概率 $\alpha(x,y)=\min\Big\{1,\ \dfrac{\pi(y)q(x\mid y)}{\pi(x)q(y\mid x)}\Big\}$ 接受 $y$，否则留在 $x$。
\end{enumerate}

\textbf{为什么它是无偏的：} 接受率的分母设计出来就是为了让 $\pi$ 满足细致平衡——这解释了「为什么未知归一化常数不影响结果」。

\textbf{与渲染的接口：} 路径追踪本质上是\textbf{在路径空间上做积分}（H3）。当路径空间的分布高度集中
（焦散、间接光、窄光源）时，逐样本独立采样效率极低；\textbf{MLT（Metropolis Light Transport）}
就是在这条路径空间上跑 MCMC，让「好路径」附近被采得更密。这是 MCMC 在 CG 里最直接的落地。

\subsubsection{随机游走}

一维简单随机游走（$\pm1$ 等概率）在 $\mathbb Z$ 上是\textbf{常返}的（$d\le2$ 常返，$d\ge3$ 瞬态——
但 $d=2$ 是临界，返回时间期望无穷）。布朗运动（C5）是随机游走的连续极限。
C8.1 的平稳分布 $\pi$ 在 Monte Carlo 里就是「长期采样的直方图」——这正是 MCMC 能估计积分的原因。

%--------------------------------------------------------------
\subsection{C9. 蒙特卡洛与方差缩减}

\subsubsection{估计量、偏差与误差量级}

要算 $I=\int_\Omega f\,dx$，取 $X_i\sim p$ 独立同分布，用
\[
\hat I_N=\frac1N\sum_{i=1}^N\frac{f(X_i)}{p(X_i)}
\]
（要 $\operatorname{supp}p\supset\operatorname{supp}f$ 以保证无偏。）

\begin{itemize}
  \item \textbf{无偏：} $\mathbb E[\hat I_N]=I$（由线性性，不需要任何极限定理）。
  \item \textbf{一致性：} $\hat I_N\to I$ 依概率（C6 弱大数定律）；
        若 $\mathbb E[(f/p)^2]<\infty$ 则 a.s. 收敛（强大数定律）。
  \item \textbf{误差量级：} 由 CLT，$\hat I_N-I$ 的分布近似 $N\!\big(0,\ \sigma^2/N\big)$，
        其中 $\sigma^2=\operatorname{Var}[f(X)/p(X)]$。
        \textbf{所以误差 $\sim\sigma/\sqrt N$：维数无关（这是 MC 战胜求积公式的唯一理由），但收敛极慢}——
        降方差只有两条路：增大 $N$（贵），或减小 $\sigma$（聪明）。
\end{itemize}

\textbf{方差公式}（与 H4 一致）：
\[
\operatorname{Var}[\hat I_N]=\frac1N\Big(\int\frac{f^2}{p}\,dx-I^2\Big)
\]

\subsubsection{方差缩减技术}

\begin{center}
\begin{tabular}{L{3.0cm}L{6.0cm}L{5.2cm}}
\toprule
\textbf{技术} & \textbf{做法} & \textbf{为什么有效} \\
\midrule
重要性采样 & 取 $p\propto|f|$；理论上 $\sigma^2=0$（$f\ge0$ 时） &
把「采样力气」花在 $f$ 大的地方；$p\propto|f|$ 使被积函数变常数 \\
\addlinespace
MIS（多策略） & $w_j=\dfrac{n_jp_j}{\sum_kn_kp_k}$（平衡启发式） &
多套 $p_j$ 各自不完美，加权后互补；避免单一策略的尾部灾难 \\
\addlinespace
分层抽样 & 把 $\Omega$ 切成若干层，各层配额采样 &
层内方差被消掉，只剩层间方差；$\sigma^2_{\text{分层}}\le\sigma^2_{\text{朴素}}$ \\
\addlinespace
对偶变量 & 用 $U$ 与 $1-U$ 配对，取平均 &
两者的负相关抵消线性项误差 \\
\addlinespace
控制变量 & $\hat I_\alpha=\hat I-\alpha(\hat C-\mathbb E C)$ &
减去一个已知期望的相关量；$\alpha^*=\operatorname{Cov}(\hat I,\hat C)/\operatorname{Var}(\hat C)$ \\
\addlinespace
俄式轮盘 & 以概率 $q$ 终止并除以 $(1-q)$ 补偿 &
由条件期望／鞅（C4）保证\textbf{无偏}，同时砍掉低贡献路径 \\
\bottomrule
\end{tabular}
\end{center}

\textbf{拟蒙特卡洛（QMC）：} 用低差异序列替代随机样本，误差量级由
\[
O(N^{-1/2})\qquad\text{改善为}\qquad O\big((\log N)^d/N\big)
\]
（Koksma--Hlawka 不等式，见 H6）。\textbf{注意维数 $d$ 出现在对数里}——
这是 QMC 在高维（路径积分）里收益递减的原因，也是它通常只用于低维积分（如面光源、环境光）的原因。

\textbf{偏差 $-$ 方差权衡：} 上述除 QMC 外都是无偏的；一旦引入有偏估计（如模糊近似、预滤波），
误差里就多出一项系统性偏差，它\textbf{不会随 $N$ 变小}——这是渲染里「模糊比噪声更致命」的数学根源。

%--------------------------------------------------------------
\subsection{C10. Itô 微积分入门}

\subsubsection{Wiener 过程与二次变差}

\textbf{Wiener 过程（布朗运动）$W_t$：} $W_0=0$；增量独立；$W_t-W_s\sim N(0,t-s)$；轨道连续。
（这四条正是 C5 的四条性质，此处升级为可积分的对象。）

\textbf{关键新事实：二次变差不为零。} 对分割 $0=t_0<\cdots<t_n=t$，
\[
\sum_i(W_{t_{i+1}}-W_{t_i})^2\ \xrightarrow{\ L^2\ }\ t
\]
把它写成微分形式就是 $dW^2=dt$，即 $(dW)^2$ 与 $dt$ \textbf{同阶}。
（对比普通函数的二次变差为零：$dt^2$ 是 $dt$ 的高阶无穷小。）

\textbf{后果：} $W_t$ 的轨道几乎必然处处不可微（C5 已提），
所以 $\int_0^t H\,dW$ \textbf{不能}定义为 Riemann--Stieltjes 积分——必须重新定义。

\subsubsection{Itô 积分与 Itô 引理}

\textbf{Itô 积分：} 对适应过程 $H$（只看 $t$ 之前的信息），沿分割取\textbf{左端点}值求和再取 $L^2$ 极限：
\[
\int_0^t H_s\,dW_s=\lim_{n\to\infty}\sum_i H_{t_i}\big(W_{t_{i+1}}-W_{t_i}\big)
\]
（取左端点使 $H_{t_i}$ 与增量 $W_{t_{i+1}}-W_{t_i}$ 独立，于是 $\mathbb E\int H\,dW=0$——这就是\textbf{鞅}，
与 C4 接上。）

\textbf{Itô 引理（链式法则的修正版）：} 若 $f\in C^{2}$，
\[
df(W_t)=f'(W_t)\,dW_t+\frac12f''(W_t)\,dt
\]
多出来的 $\frac12f''dt$ \textbf{完全来自二次变差 $dW^2=dt$}。
若照搬普通链式法则，这一项就会丢掉——这是初学者 90\% 错误的来源。

\textbf{算例：} 令 $X_t=\exp(\sigma W_t-\frac{\sigma^2}{2}t)$，用 $f=e^x$：
\[
dX_t=X_t\Big(\sigma\,dW_t-\frac{\sigma^2}{2}dt\Big)+\frac12X_t\sigma^2dt=X_t\sigma\,dW_t
\]
故 $X_t$ 是鞅，且 $\mathbb E X_t=X_0$。这个「补 $\frac{\sigma^2}{2}$ 的修正项」在金融与统计物理里反复出现。

\subsubsection{SDE 与 Fokker--Planck}

一般 SDE：
\[
dX_t=\mu(X_t,t)\,dt+\sigma(X_t,t)\,dW_t
\]

\textbf{例（几何布朗运动）：} $dX_t=\mu X_t\,dt+\sigma X_t\,dW_t$，解为
$X_t=X_0\exp\big((\mu-\frac{\sigma^2}{2})t+\sigma W_t\big)$——注意指数上的 $-\frac{\sigma^2}{2}$ 修正，
它使 $\mathbb E X_t=e^{\mu t}X_0$ \textbf{与} 典型样本轨道 $e^{(\mu-\sigma^2/2)t}$ 的增长率不同。
\textbf{这个差距不是笔误，而是「期望 vs 典型」的分野}，也是 C6 中 a.s. 收敛与 $L^1$ 收敛方向不同的另一种面孔。

\textbf{Fokker--Planck（前向 Kolmogorov）：} 转移密度 $p(x,t)$ 满足
\[
\frac{\partial p}{\partial t}=-\frac{\partial}{\partial x}\big(\mu p\big)+\frac12\frac{\partial^2}{\partial x^2}\big(\sigma^2p\big)
\]
右边是「对流 $+$ 扩散」——这正是 I 章流体的对流扩散方程的形式。
\textbf{随机微分方程（SDE）与抛物型 PDE 是同一件事的两种语言}：SDE 描述单条轨道，Fokker--Planck 描述密度的演化。

\textbf{在 CG 中的接口：} M5 的朗之万方程 $m\ddot x=-\gamma\dot x+\sqrt{2\gamma k_BT}\,\xi(t)+F(x)$
就是 SDE 的牛顿版（$\xi$ 是白噪声，即 $dW/dt$）；
其 Fokker--Planck 方程给出位置与速度的联合密度演化——分子动力学与被动标量（烟雾）模拟的理论骨架。
\section{D. 经典力学的数学结构}

\subsection{D1. 拉格朗日力学}

\textbf{作用量：} $S[q] = \displaystyle\int_{t_1}^{t_2} L(q, \dot{q}, t)\,dt$

\textbf{拉格朗日量：} $L = T - V$（动能减势能）

\textbf{欧拉--拉格朗日方程：}
\[
\frac{d}{dt}\frac{\partial L}{\partial \dot{q}_i} - \frac{\partial L}{\partial q_i} = 0
\]

\textbf{推导（变分法）：} 令 $q(t) \to q(t) + \epsilon \eta(t)$，$\eta(t_1) = \eta(t_2) = 0$，要求 $\delta S = 0$：
\[
\delta S = \int_{t_1}^{t_2} \left(\frac{\partial L}{\partial q}\eta + \frac{\partial L}{\partial \dot{q}}\dot{\eta}\right)dt = \int_{t_1}^{t_2}\left(\frac{\partial L}{\partial q} - \frac{d}{dt}\frac{\partial L}{\partial \dot{q}}\right)\eta\,dt = 0
\]
由 $\eta$ 的任意性得欧拉--拉格朗日方程。

\subsection{D2. Noether 定理}

每一个连续对称性对应一个守恒量：

\begin{center}
\begin{tabular}{ll}
\toprule
\textbf{对称性} & \textbf{守恒量} \\
\midrule
时间平移不变性 & 能量守恒 \\
空间平移不变性 & 动量守恒 \\
空间旋转不变性 & 角动量守恒 \\
规范不变性 $U(1)$ & 电荷守恒 \\
\bottomrule
\end{tabular}
\end{center}

\textbf{数学表述：} 如果拉格朗日量在变换 $q \to q + \epsilon\,\delta q$ 下不变（至多差一个全导数），则：
\[
\frac{d}{dt}\left(\frac{\partial L}{\partial \dot{q}} \cdot \delta q\right) = 0
\]

\subsection{D3. 哈密顿力学}

\textbf{广义动量：} $p_i = \dfrac{\partial L}{\partial \dot{q}_i}$

\textbf{哈密顿量：} $H(q, p) = \displaystyle\sum_i p_i \dot{q}_i - L$

\textbf{哈密顿方程：}
\[
\dot{q}_i = \frac{\partial H}{\partial p_i}, \quad \dot{p}_i = -\frac{\partial H}{\partial q_i}
\]

\textbf{泊松括号：}
\[
\{f, g\} = \sum_i \left(\frac{\partial f}{\partial q_i}\frac{\partial g}{\partial p_i} - \frac{\partial f}{\partial p_i}\frac{\partial g}{\partial q_i}\right)
\]

运动方程：$\dot{f} = \{f, H\}$

\textbf{辛结构：} 相空间上有辛形式 $\omega = \sum dp_i \wedge dq_i$，哈密顿流保持辛形式（刘维尔定理：相空间体积守恒）。

\subsection{D4. 哈密顿--雅可比理论与正则变换}

\textbf{哈密顿--雅可比方程：}
\[
\frac{\partial S}{\partial t} + H\left(q, \frac{\partial S}{\partial q}, t\right) = 0
\]

这是经典力学与量子力学的桥梁：将 $S$ 替换为 $\frac{\hbar}{i}\ln\psi$，取 $\hbar \to 0$ 极限，就得到薛定谔方程的经典极限。

\textbf{正则变换：} 保持辛结构的坐标变换 $(q,p) \to (Q,P)$。

\textbf{生成函数：} 四种类型 $F_1(q,Q)$，$F_2(q,P)$，$F_3(p,Q)$，$F_4(p,P)$。
\section{E. 场论与微分形式}

\subsection{E1. 向量分析的微分形式统一}

\begin{center}
\begin{tabular}{lll}
\toprule
\textbf{向量分析} & \textbf{微分形式} & \textbf{运算} \\
\midrule
标量场 $f$ & 0-形式 & $df$ = 梯度 \\
向量场 $\mathbf{F}$ & 1-形式 $\omega$ & $d\omega$ = 旋度 \\
向量场 $\mathbf{F}$ & 2-形式 $\eta$ & $d\eta$ = 散度 \\
$\nabla \cdot (\nabla \times \mathbf{F}) = 0$ & $d^2 = 0$ & \\
\bottomrule
\end{tabular}
\end{center}

\subsection{E2. 麦克斯韦方程的微分形式表述}

\textbf{电磁场张量：} $F = dA$，其中 $A$ 是电磁势 1-形式。

\[
dF = 0 \quad \text{（Faraday 定律 + 无磁单极）}
\]
\[
d{*}F = J \quad \text{（Ampère 定律 + Gauss 定律）}
\]

其中 $*$ 是 Hodge 星算子，$J$ 是电流 3-形式。

\textbf{规范不变性：} $A \to A + d\lambda$ 不改变 $F = dA$（因为 $d^2 = 0$）。这就是 $U(1)$ 规范对称性。

\subsection{E3. 闵可夫斯基时空}

\textbf{度规：} $\eta_{\mu\nu} = \text{diag}(-1, 1, 1, 1)$

\textbf{间隔：} $ds^2 = -c^2 dt^2 + dx^2 + dy^2 + dz^2$

\textbf{四维向量：}
\begin{itemize}
  \item 位置：$x^\mu = (ct, x, y, z)$
  \item 四维动量：$p^\mu = (E/c, \mathbf{p})$
  \item 四维电流：$J^\mu = (c\rho, \mathbf{J})$
\end{itemize}

\textbf{质能关系：} $p^\mu p_\mu = -m^2 c^2$，即 $E^2 = p^2 c^2 + m^2 c^4$

\subsection{E4. 规范场论初步}

\begin{center}
\begin{tabular}{lll}
\toprule
\textbf{规范群} & \textbf{规范场} & \textbf{物理} \\
\midrule
$U(1)$ & 电磁场 $A_\mu$ & 电磁力 \\
$SU(2)$ & $W$ 玻色子 & 弱力 \\
$SU(3)$ & 胶子 $G_\mu^a$ & 强力 \\
\bottomrule
\end{tabular}
\end{center}

\textbf{协变导数：} $D_\mu = \partial_\mu - igA_\mu$

\textbf{场强：} $F_{\mu\nu} = \partial_\mu A_\nu - \partial_\nu A_\mu - ig[A_\mu, A_\nu]$

\textbf{Yang--Mills 作用量：} $S = -\dfrac{1}{4}\displaystyle\int F_{\mu\nu}^a F^{a\mu\nu}\,d^4x$
\section{F. 量子力学的数学结构}

\subsection{F1. 希尔伯特空间与态}

\textbf{量子态：} 希尔伯特空间 $\mathcal{H}$ 中的射线（等价类 $|\psi\rangle \sim e^{i\theta}|\psi\rangle$）。

\textbf{Dirac 符号：}
\begin{itemize}
  \item 右矢 $|\psi\rangle \in \mathcal{H}$
  \item 左矢 $\langle\psi| \in \mathcal{H}^*$
  \item 内积 $\langle\phi|\psi\rangle$
  \item 外积 $|\psi\rangle\langle\phi|$（秩 1 算子）
\end{itemize}

\textbf{叠加原理：} $|\psi\rangle = \alpha|0\rangle + \beta|1\rangle$，$|\alpha|^2 + |\beta|^2 = 1$

\subsection{F2. 可观测量与算子}

\textbf{可观测量 = 自伴算子：} $\hat{A} = \hat{A}^\dagger$

\textbf{测量公设：} 测量 $\hat{A}$ 得到特征值 $a_n$ 的概率为 $|\langle a_n|\psi\rangle|^2$

\textbf{期望值：} $\langle \hat{A} \rangle = \langle\psi|\hat{A}|\psi\rangle$

\textbf{不确定性原理：}
\[
\Delta A \cdot \Delta B \geq \frac{1}{2}|\langle[\hat{A}, \hat{B}]\rangle|
\]
其中 $[\hat{A}, \hat{B}] = \hat{A}\hat{B} - \hat{B}\hat{A}$ 是对易子。

\textbf{基本对易关系：} $[\hat{x}, \hat{p}] = i\hbar$

\subsection{F3. 薛定谔方程}

\[
i\hbar\frac{\partial}{\partial t}|\psi(t)\rangle = \hat{H}|\psi(t)\rangle
\]

\textbf{形式解：} $|\psi(t)\rangle = e^{-i\hat{H}t/\hbar}|\psi(0)\rangle$

时间演化算子 $U(t) = e^{-i\hat{H}t/\hbar}$ 是\textbf{酉算子}（保持内积，即保持概率）。

\subsection{F4. 角动量与自旋}

\textbf{角动量算子：} $\hat{\mathbf{L}} = \hat{\mathbf{r}} \times \hat{\mathbf{p}}$

\textbf{对易关系：} $[\hat{L}_i, \hat{L}_j] = i\hbar\varepsilon_{ijk}\hat{L}_k$（$\mathfrak{so}(3)$ 的李代数）

\textbf{量子化：}
\begin{itemize}
  \item $\hat{L}^2|l,m\rangle = \hbar^2 l(l+1)|l,m\rangle$
  \item $\hat{L}_z|l,m\rangle = \hbar m|l,m\rangle$
  \item $l = 0, 1, 2, \ldots$，$m = -l, -l+1, \ldots, l$
\end{itemize}

球谐函数 $Y_l^m(\theta, \phi)$ 是 $\hat{L}^2$ 和 $\hat{L}_z$ 的共同本征函数。

\textbf{自旋：} 内禀角动量，不对应空间旋转，但满足相同的代数。电子自旋 $s = 1/2$。

\subsection{F5. 路径积分（Feynman 表述）}
\label{subsec:pi-brief}

\[
\langle x_f, t_f | x_i, t_i \rangle = \int \mathcal{D}[x(t)]\, e^{iS[x]/\hbar}
\]

其中 $S[x] = \displaystyle\int L\,dt$ 是经典作用量。

\textbf{经典极限：} $\hbar \to 0$ 时，相位快速振荡，只有作用量取极值的路径（经典路径）贡献显著（稳相近似）。

\textbf{完整推导（时间切片法）、自由粒子与谐振子的显式计算、Wick 转动，以及与路径追踪的逐项类比，见第~\ref{sec:path-integral}~节（K 章）。}

\subsection{F6. 量子力学与 CG 的交叉}

\begin{center}
\begin{tabular}{ll}
\toprule
\textbf{量子概念} & \textbf{CG 对应} \\
\midrule
波函数 $\psi$ & 光场/辐射度分布 \\
$|\psi|^2$ = 概率密度 & 辐射度 = 能量密度 \\
薛定谔方程 & 波动方程/亥姆霍兹方程 \\
路径积分 & 路径追踪 \\
算子谱 & 频率分解/球谐分解 \\
不确定性原理 & 时频分辨率的 trade-off \\
\bottomrule
\end{tabular}
\end{center}
\section{G. 数值方法与计算数学}

\subsection{G1. 常微分方程数值解}

\subsubsection{经典 RK4}

\begin{align*}
k_1 &= f(t_n, y_n) \\
k_2 &= f(t_n + h/2,\; y_n + hk_1/2) \\
k_3 &= f(t_n + h/2,\; y_n + hk_2/2) \\
k_4 &= f(t_n + h,\; y_n + hk_3) \\
y_{n+1} &= y_n + \frac{h}{6}(k_1 + 2k_2 + 2k_3 + k_4)
\end{align*}

\subsubsection{辛积分器}

\textbf{Verlet 积分：}
\[
x_{n+1} = 2x_n - x_{n-1} + a_n h^2
\]

\textbf{Leapfrog（蛙跳）：}
\begin{align*}
v_{n+1/2} &= v_{n-1/2} + a_n h \\
x_{n+1} &= x_n + v_{n+1/2} h
\end{align*}

辛积分器能量在 bounded 范围内振荡，不会系统性漂移。布料模拟、刚体模拟、粒子系统都用辛积分器。

\subsection{G2. 偏微分方程数值解}
\label{subsec:pde-numerics}

\subsubsection{有限差分法（FDM）}

\[
\frac{\partial u}{\partial x} \approx \frac{u_{i+1} - u_{i-1}}{2h}, \quad
\frac{\partial^2 u}{\partial x^2} \approx \frac{u_{i+1} - 2u_i + u_{i-1}}{h^2}
\]

\textbf{CFL 条件：} 对波动方程，$\dfrac{c\,\Delta t}{\Delta x} \leq 1$。

\subsubsection{有限元法（FEM）}

\textbf{弱形式：}
\[
\int_\Omega \nabla u \cdot \nabla v\,d\Omega = \int_\Omega fv\,d\Omega
\]

\textbf{刚度矩阵：} $K_{ij} = \displaystyle\int_\Omega \nabla\phi_i \cdot \nabla\phi_j\,d\Omega$

\subsubsection{谱方法}

用全局正交基（傅里叶、切比雪夫、球谐）展开解。对光滑解指数收敛。

\textbf{有限差分/有限元/谱方法的稳定性分析、格式对比与半拉格朗日方法，见第~\ref{sec:de-numerics}~节（第八层：数值方法）。}

\subsection{G3. 数值线性代数}

\begin{center}
\begin{tabular}{lll}
\toprule
\textbf{分解} & \textbf{形式} & \textbf{应用} \\
\midrule
LU & $A = LU$ & 解线性系统 \\
QR & $A = QR$ & 最小二乘、特征值 \\
SVD & $A = U\Sigma V^T$ & PCA、伪逆、低秩近似 \\
Cholesky & $A = LL^T$（正定） & 快速求解 \\
特征分解 & $A = P\Lambda P^{-1}$ & 谱分析、矩阵函数 \\
Schur & $A = QTQ^*$ & 数值稳定的特征值计算 \\
\bottomrule
\end{tabular}
\end{center}

\textbf{迭代法：} 共轭梯度法（CG）、GMRES、多重网格。

\subsection{G4. 逼近论}

\subsubsection{Padé 近似}

\[
f(x) \approx \frac{P_m(x)}{Q_n(x)} = \frac{a_0 + a_1 x + \cdots + a_m x^m}{1 + b_1 x + \cdots + b_n x^n}
\]

可以逼近有极点的函数（泰勒级数做不到），收敛半径通常更大。

\subsubsection{切比雪夫逼近}

切比雪夫节点：$x_k = \cos\dfrac{(2k+1)\pi}{2(n+1)}$（避免 Runge 现象）。

\subsection{G5. 优化}

\begin{itemize}
  \item 梯度下降：$x_{k+1} = x_k - \alpha \nabla f(x_k)$
  \item 牛顿法：$x_{k+1} = x_k - H^{-1}\nabla f$
  \item L-BFGS（拟牛顿，CG 中广泛使用）
  \item 拉格朗日乘子法、KKT 条件
\end{itemize}
\section{H. 计算机图形学中的数学}

\subsection{H1. 离散微分几何}

\textbf{离散拉普拉斯算子：}
\[
(\Delta f)_i = \sum_{j \in N(i)} w_{ij}(f_j - f_i)
\]

权重选择：
\begin{itemize}
  \item 均匀权重：$w_{ij} = 1$
  \item 余切权重：$w_{ij} = \cot\alpha_{ij} + \cot\beta_{ij}$（FEM 离散化，推荐）
\end{itemize}

\textbf{离散高斯曲率（角度亏损）：}
\[
K_i = 2\pi - \sum_{j} \theta_j
\]

\textbf{离散平均曲率：}
\[
H_i = \frac{1}{2A_i}\sum_{j} (\cot\alpha_{ij} + \cot\beta_{ij})(v_i - v_j)
\]

\subsection{H2. 渲染方程与 BRDF}

\textbf{渲染方程：}
\[
L_o(\mathbf{x}, \omega_o) = L_e(\mathbf{x}, \omega_o) + \int_{\Omega} f_r(\mathbf{x}, \omega_i, \omega_o)\, L_i(\mathbf{x}, \omega_i)\, \cos\theta_i\, d\omega_i
\]

\textbf{算子形式：} $L = L_e + \mathcal{T}L$

\textbf{Neumann 级数：} $L = \displaystyle\sum_{k=0}^{\infty} \mathcal{T}^k L_e$

\textbf{BRDF 定义：}
\[
f_r(\omega_i, \omega_o) = \frac{dL_o(\omega_o)}{L_i \cos\theta_i\, d\omega_i}
\]

\textbf{性质：} 互易性、能量守恒、正定性。

\textbf{微表面模型：}
\[
f_r = f_{\text{diffuse}} + \frac{D \cdot F \cdot G}{4\cos\theta_i \cos\theta_o}
\]

\subsection{H3. 路径空间公式化}

\textbf{路径} $\bar{x} = (x_0, x_1, \ldots, x_k)$

\textbf{路径贡献：}
\[
f(\bar{x}) = L_e(x_0 \to x_1) \prod_{i=1}^{k-1} f_r(x_i)\, G(x_i \leftrightarrow x_{i+1})
\]

\textbf{几何项：}
\[
G(x \leftrightarrow y) = \frac{\cos\theta_x \cos\theta_y}{|x - y|^2}
\]

\textbf{路径积分：}
\[
I = \int_{\Omega} f(\bar{x})\, d\mu(\bar{x})
\]

\subsection{H4. 蒙特卡洛积分}

\[
\int_\Omega f(x)\,dx \approx \frac{1}{N}\sum_{i=1}^N \frac{f(X_i)}{p(X_i)}, \quad X_i \sim p
\]

\textbf{方差：}
\[
\text{Var}[\hat{I}] = \frac{1}{N}\left(\int \frac{f^2}{p}\,dx - I^2\right)
\]

\textbf{MIS 权重（平衡启发式）：}
\[
w_j(x) = \frac{n_j\, p_j(x)}{\sum_k n_k\, p_k(x)}
\]

\subsection{H5. 球谐函数在 CG 中的应用}

\begin{itemize}
  \item 环境光照表示（通常 $l \leq 2$，9 个系数）
  \item 预计算辐射传输（PRT）
  \item 球面卷积在球谐域变成逐元素乘法
\end{itemize}

\subsection{H6. 采样理论}

\subsubsection{低差异序列}

\begin{center}
\begin{tabular}{lll}
\toprule
\textbf{序列} & \textbf{构造} & \textbf{特点} \\
\midrule
Halton & 基于不同素数的 radical inverse & 简单，高维有相关 \\
Sobol & 基于方向数 & 高维表现好 \\
Cranley--Patterson 旋转 & 随机偏移 & 去相关 \\
\bottomrule
\end{tabular}
\end{center}

\textbf{Koksma--Hlawka 不等式：}
\[
\left|\frac{1}{N}\sum f(x_i) - \int f\right| \leq V(f) \cdot D_N^*
\]

\subsubsection{蓝噪声}

功率谱中低频分量被抑制，高频分量保留。视觉上最均匀的随机分布。

\subsection{H7. 辐射度量学}

\begin{center}
\begin{tabular}{llll}
\toprule
\textbf{量} & \textbf{符号} & \textbf{单位} & \textbf{定义} \\
\midrule
辐射通量 & $\Phi = dQ/dt$ & W & \\
辐射强度 & $I = d\Phi/d\omega$ & W/sr & 点源 \\
辐射照度 & $E = d\Phi/dA$ & W/m$^2$ & 到达表面 \\
辐射度 & $L = \frac{d^2\Phi}{dA\cos\theta\,d\omega}$ & W/(m$^2\cdot$sr) & \textbf{核心量} \\
\bottomrule
\end{tabular}
\end{center}

\textbf{光度学转换：}
\[
\Phi_v = 683 \int_0^\infty \Phi_e(\lambda)\, V(\lambda)\, d\lambda
\]

\subsection{H8. 物理模拟}

\subsubsection{流体模拟}
\label{subsec:cg-fluid}

\textbf{Navier--Stokes 方程（不可压缩）：}
\[
\frac{\partial \mathbf{u}}{\partial t} + (\mathbf{u}\cdot\nabla)\mathbf{u} = -\frac{1}{\rho}\nabla p + \nu\nabla^2\mathbf{u} + \mathbf{g}, \quad \nabla \cdot \mathbf{u} = 0
\]

\textbf{压力投影：}
\begin{enumerate}
  \item $\mathbf{u}^* = \mathbf{u}^n + \Delta t(\text{对流} + \text{粘性})$
  \item $\nabla^2 p = \frac{1}{\Delta t}\nabla\cdot\mathbf{u}^*$
  \item $\mathbf{u}^{n+1} = \mathbf{u}^* - \Delta t\nabla p$
\end{enumerate}

\textbf{守恒律的完整推导、应力张量本构关系、涡量动力学与守恒性保持，见第~\ref{sec:fluid}~节（I 章）。}

\subsubsection{刚体模拟}

\textbf{运动方程：}
\[
\dot{\mathbf{x}} = \mathbf{v}, \quad \dot{R} = [\boldsymbol{\omega}]_\times R
\]
\[
m\dot{\mathbf{v}} = \mathbf{F}, \quad I\dot{\boldsymbol{\omega}} + \boldsymbol{\omega} \times I\boldsymbol{\omega} = \boldsymbol{\tau}
\]

\subsubsection{布料与软体}

\begin{itemize}
  \item 质点--弹簧模型：$m\ddot{x} = -k(x - x_0) - c\dot{x} + f_{\text{ext}}$
  \item 位置基动力学（PBD/XPBD）：直接操作位置满足约束，无条件稳定
\end{itemize}
\section{I. 流体力学}
\label{sec:fluid}

\subsection{I1. 守恒律}

\textbf{质量守恒：}
\[
\frac{\partial \rho}{\partial t} + \nabla \cdot (\rho \mathbf{u}) = 0
\]

\textbf{动量守恒：}
\[
\frac{\partial(\rho \mathbf{u})}{\partial t} + \nabla \cdot (\rho \mathbf{u} \otimes \mathbf{u}) = \nabla \cdot \boldsymbol{\sigma} + \rho \mathbf{g}
\]

\textbf{物质导数：}
\[
\frac{D}{Dt} = \frac{\partial}{\partial t} + \mathbf{u} \cdot \nabla
\]

\subsection{I2. 应力张量与粘性}

\[
\boldsymbol{\sigma} = -p\mathbf{I} + \boldsymbol{\tau}
\]

\textbf{牛顿流体本构关系：}
\[
\boldsymbol{\tau} = 2\mu\mathbf{D} + \lambda(\nabla\cdot\mathbf{u})\mathbf{I}
\]

其中 $\mathbf{D} = \frac{1}{2}(\nabla\mathbf{u} + (\nabla\mathbf{u})^T)$ 是应变率张量。

\textbf{非牛顿流体：}
\begin{center}
\begin{tabular}{lll}
\toprule
\textbf{类型} & \textbf{特征} & \textbf{例子} \\
\midrule
剪切稀化 & 粘度随剪切率增大而减小 & 油漆、血液 \\
剪切增稠 & 粘度随剪切率增大而增大 & 玉米淀粉水 \\
宾汉流体 & 有屈服应力 & 牙膏、泥浆 \\
粘弹性 & 同时有粘性和弹性 & 聚合物、面团 \\
\bottomrule
\end{tabular}
\end{center}

\subsection{I3. 雷诺数}

\[
Re = \frac{UL}{\nu} = \frac{\text{惯性力}}{\text{粘性力}}
\]

\begin{center}
\begin{tabular}{ll}
\toprule
\textbf{$Re$ 范围} & \textbf{流动特征} \\
\midrule
$Re \ll 1$ & 蠕动流（Stokes 流），粘性主导 \\
$Re \sim 10^2$--$10^3$ & 层流，可能出现分离 \\
$Re > 4000$（管流） & 湍流 \\
$Re \sim 10^6+$ & 完全湍流 \\
\bottomrule
\end{tabular}
\end{center}

\subsection{I4. 涡量动力学}

\textbf{涡量：} $\boldsymbol{\omega} = \nabla \times \mathbf{u}$

\textbf{涡量方程（不可压缩）：}
\[
\frac{D\boldsymbol{\omega}}{Dt} = (\boldsymbol{\omega}\cdot\nabla)\mathbf{u} + \nu\nabla^2\boldsymbol{\omega}
\]

$(\boldsymbol{\omega}\cdot\nabla)\mathbf{u}$：涡量拉伸项（三维特有，湍流能量级联的关键）。

\textbf{Kelvin 环量定理（无粘）：}
\[
\frac{D\Gamma}{Dt} = 0, \quad \Gamma = \oint_C \mathbf{u}\cdot d\mathbf{l}
\]

在 CG 中：涡量约束（Vorticity Confinement）补偿数值耗散对涡量的抹平。

\subsection{I5. 湍流}

\textbf{Kolmogorov 能量谱：} $E(k) \sim k^{-5/3}$（惯性子区）

\textbf{Kolmogorov 尺度：}
\[
\eta = \left(\frac{\nu^3}{\epsilon}\right)^{1/4}
\]

完全解析湍流需要网格数 $\sim Re^{9/4}$，对 $Re = 10^6$ 需要 $\sim 10^{13}$ 个网格点——完全不可行。

\textbf{RANS：} 将速度分解为均值和脉动 $\mathbf{u} = \bar{\mathbf{u}} + \mathbf{u}'$，引入雷诺应力 $-\rho\overline{u_i'u_j'}$，需要模型封闭（$k$-$\varepsilon$、$k$-$\omega$ 等）。

\textbf{LES：} 直接解析大尺度，对小尺度用亚格子模型。

\subsection{I6. 边界层}

\textbf{边界层厚度：} $\delta \sim \dfrac{L}{\sqrt{Re}}$

\textbf{无滑移条件：} $\mathbf{u}|_{\text{wall}} = 0$

\subsection{I7. 可压缩流与激波}

\textbf{马赫数：} $Ma = \dfrac{U}{c}$

$Ma > 1$ 时出现激波。Rankine--Hugoniot 条件描述激波两侧的守恒关系。

\subsection{I8. 自由表面与多相流}

\begin{itemize}
  \item \textbf{Level Set：} 用符号距离函数 $\phi$ 表示界面（$\phi = 0$ 为界面），输运方程 $\frac{\partial \phi}{\partial t} + \mathbf{u}\cdot\nabla\phi = 0$
  \item \textbf{VOF：} 用体积分数 $f \in [0,1]$ 表示，天然守恒
  \item \textbf{SPH：} 拉格朗日描述，天然处理自由表面
\end{itemize}

\textbf{表面张力：} $\Delta p = \sigma \kappa$

\subsection{I9. 数值模拟中的守恒性}

\subsubsection{有限体积法——天然守恒}

相邻单元共享面的通量大小相等方向相反 $\to$ 全局守恒自动满足。

\subsubsection{能量稳定格式}

连续情况：
\[
\frac{d}{dt}\int \frac{1}{2}|\mathbf{u}|^2\,dV = -\nu\int|\nabla\mathbf{u}|^2\,dV \leq 0
\]

离散要求：对流项的离散必须满足“反对称性”或“守恒性”。

\subsubsection{散度自由条件的保持}

\begin{itemize}
  \item 投影法（Chorin/Temam）
  \item MAC/Staggered 网格：离散散度精确为零
  \item 约束输运：对涡量或流函数求解
\end{itemize}

\subsubsection{时间步长约束}

\begin{center}
\begin{tabular}{ll}
\toprule
\textbf{约束来源} & \textbf{条件} \\
\midrule
对流（CFL） & $\Delta t \leq C\frac{\Delta x}{|\mathbf{u}|_{\max}}$ \\
粘性 & $\Delta t \leq C\frac{\Delta x^2}{\nu}$ \\
表面张力 & $\Delta t \leq C\sqrt{\frac{\rho\,\Delta x^3}{\sigma}}$ \\
弹性波 & $\Delta t \leq C\frac{\Delta x}{c_s}$ \\
\bottomrule
\end{tabular}
\end{center}
\section{J. 光线追踪的深层数学}

\subsection{J1. 从几何光学到波动光学}

\textbf{程函方程：}
\[
|\nabla S|^2 = n^2(\mathbf{x})
\]

\textbf{光线方程：}
\[
\frac{d}{ds}\left(n\frac{d\mathbf{x}}{ds}\right) = \nabla n
\]

\textbf{亥姆霍兹方程：}
\[
\nabla^2 \mathbf{E} + k^2 \mathbf{E} = 0, \quad k = \frac{2\pi}{\lambda}
\]

\textbf{自由空间格林函数：}
\[
G(\mathbf{x}, \mathbf{x}') = \frac{e^{ik|\mathbf{x}-\mathbf{x}'|}}{4\pi|\mathbf{x}-\mathbf{x}'|}
\]

\subsection{J2. 衍射}

\subsubsection{夫琅禾费衍射（远场）}

衍射图案是孔径函数的傅里叶变换：
\[
U(\theta_x, \theta_y) \propto \mathcal{F}\{A(x,y)\}\bigg|_{f_x = \theta_x/\lambda,\; f_y = \theta_y/\lambda}
\]

\subsubsection{菲涅尔衍射（近场）}

\[
U(x,y,z) = \frac{e^{ikz}}{i\lambda z}\iint u_0(x',y') \exp\left[\frac{ik}{2z}((x-x')^2 + (y-y')^2)\right]dx'\,dy'
\]

\subsection{J3. 偏振}

\textbf{Fresnel 方程：}
\[
r_s = \frac{n_1\cos\theta_i - n_2\cos\theta_t}{n_1\cos\theta_i + n_2\cos\theta_t}, \quad
r_p = \frac{n_2\cos\theta_i - n_1\cos\theta_t}{n_2\cos\theta_i + n_1\cos\theta_t}
\]

\textbf{Brewster 角：} $\theta_B = \arctan(n_2/n_1)$，此时 $r_p = 0$。

\subsection{J4. 相干性}

\textbf{互相干函数：}
\[
\Gamma(\mathbf{r}_1, \mathbf{r}_2, \tau) = \langle E^*(\mathbf{r}_1, t)\, E(\mathbf{r}_2, t+\tau) \rangle
\]

$|\gamma| = 1$：完全相干（激光）；$|\gamma| = 0$：完全不相干（热光源）。

大多数渲染假设不相干光（强度叠加）。

\subsection{J5. 参与介质}

\textbf{辐射传输方程（RTE）：}
\[
\frac{dL}{ds} = -\sigma_t L + \sigma_a L_e + \sigma_s \int_{S^2} p(\omega', \omega)\, L(\omega')\, d\omega'
\]

\textbf{Beer--Lambert 定律：} $T = e^{-\tau}$，$\tau = \displaystyle\int_0^s \sigma_t\, ds'$

\textbf{Delta tracking（Woodcock tracking）：} 对非均匀介质，用最大消光系数 $\sigma_{\max}$ 做包络。

\subsection{J6. 次表面散射（SSS）}

\textbf{BSSRDF：} $S(x_i, \omega_i;\, x_o, \omega_o)$，入射点和出射点可以不同。

\textbf{偶极近似：} 将次表面散射近似为两个虚拟点源，得到解析的扩散剖面。

\subsection{J7. 路径采样策略}

\begin{center}
\begin{tabular}{lll}
\toprule
\textbf{方法} & \textbf{描述} & \textbf{优势} \\
\midrule
路径追踪 (PT) & 从相机逐步采样 & 简单通用 \\
双向路径追踪 (BDPT) & 光源+相机同时构建 & 处理焦散 \\
光子映射 & 从光源发射光子 & 焦散、参与介质 \\
MLT & MCMC 在路径空间 & 复杂光路 \\
VCM/UPBP & BDPT + 光子映射统一 & 通用 \\
SPPM & 渐进收敛的光子映射 & 焦散 \\
\bottomrule
\end{tabular}
\end{center}

\subsection{J8. 焦散与灾变理论}

几何光学在焦散处预言无限亮度（振幅发散）。波动光学给出有限值（Airy 函数、Pearcey 函数）。

焦散的形态（尖点、褶皱）由灾变理论分类。
\section{K. Feynman 路径积分}
\label{sec:path-integral}

\subsection{K1. 推导（时间切片法）}

\textbf{传播子：}
\[
K(x_f, t_f;\, x_i, t_i) = \langle x_f | e^{-i\hat{H}(t_f-t_i)/\hbar} | x_i \rangle
\]

将时间区间分为 $N$ 段，$\epsilon = (t_f - t_i)/N$，在每个时间片插入完备性关系 $\int |x\rangle\langle x|\,dx = 1$。

对于 $\hat{H} = \frac{\hat{p}^2}{2m} + V(\hat{x})$，取 $N \to \infty$：
\[
\boxed{K(x_f, t_f;\, x_i, t_i) = \int \mathcal{D}[x(t)]\, \exp\left[\frac{i}{\hbar}S[x(t)]\right]}
\]

其中 $S[x] = \displaystyle\int_{t_i}^{t_f} L(x, \dot{x}, t)\,dt$。

\subsection{K2. 物理含义}

\begin{itemize}
  \item 每条路径贡献相位 $e^{iS/\hbar}$
  \item 经典路径是 $\delta S = 0$ 的驻点，附近路径相干叠加
  \item 远离经典路径的路径相位快速振荡，相互抵消
  \item $\hbar \to 0$：只有经典路径贡献 $\to$ 经典力学
\end{itemize}

\subsection{K3. 具体计算}

\subsubsection{自由粒子}

\[
K_{\text{free}} = \sqrt{\frac{m}{2\pi i\hbar(t_f-t_i)}} \exp\left[\frac{im(x_f-x_i)^2}{2\hbar(t_f-t_i)}\right]
\]

\subsubsection{谐振子}

\[
K_{\text{HO}} = \sqrt{\frac{m\omega}{2\pi i\hbar\sin\omega T}} \exp\left[\frac{im\omega}{2\hbar\sin\omega T}\left((x_f^2+x_i^2)\cos\omega T - 2x_fx_i\right)\right]
\]

对于二次作用量（自由粒子、谐振子），路径积分可以精确计算（高斯积分）。

\subsection{K4. Wick 转动与统计力学}

$t \to -i\tau$：
\[
e^{iS/\hbar} \to e^{-S_E/\hbar}
\]

路径积分变为配分函数形式：
\[
Z = \int \mathcal{D}[x(\tau)]\, e^{-S_E[x]/\hbar}
\]

$d$ 维量子场论 $\leftrightarrow$ $(d+1)$ 维经典统计力学。

\subsection{K5. 与路径追踪的类比}

\begin{center}
\begin{tabular}{ll}
\toprule
\textbf{量子力学} & \textbf{渲染} \\
\midrule
路径 $x(t)$ & 光路 $(x_0, x_1, \ldots, x_k)$ \\
作用量 $S[x]$ & 路径贡献 $f(\bar{x})$ \\
$e^{iS/\hbar}$（相位） & 路径 throughput（实数） \\
$\int \mathcal{D}x$ & $\int_\Omega d\mu(\bar{x})$ \\
经典路径（$\delta S = 0$） & 镜面反射路径（Fermat 原理） \\
微扰展开（Feynman 图） & Neumann 级数（弹射次数） \\
$\hbar \to 0$ & $\lambda \to 0$（几何光学极限） \\
\bottomrule
\end{tabular}
\end{center}

\textbf{关键区别：} 量子路径积分的权重是复数相位（干涉），渲染路径积分的权重是正实数（能量）。
\section{L. Feynman 图与量子场论微扰论}

\subsection{L1. 二次量子化}

\textbf{Fock 空间：}
\[
\mathcal{F} = \mathbb{C} \oplus \mathcal{H} \oplus (\mathcal{H}\otimes_s\mathcal{H}) \oplus \cdots
\]

\textbf{产生/湮灭算子：}
\begin{itemize}
  \item 玻色子：$[a(\mathbf{p}), a^\dagger(\mathbf{p}')] = \delta(\mathbf{p}-\mathbf{p}')$
  \item 费米子：$\{a(\mathbf{p}), a^\dagger(\mathbf{p}')\} = \delta(\mathbf{p}-\mathbf{p}')$
\end{itemize}

\subsection{L2. Dyson 级数}

\[
U = 1 + \left(-\frac{i}{\hbar}\right)\int dt_1\, H_{\text{int}}(t_1) + \left(-\frac{i}{\hbar}\right)^2\int dt_1 dt_2\, T[H_{\text{int}}(t_1)H_{\text{int}}(t_2)] + \cdots
\]

每一项对应一个 Feynman 图的阶数。

\subsection{L3. Feynman 规则}

\begin{enumerate}
  \item \textbf{外线：} 入射/出射粒子的波函数
  \item \textbf{内线（传播子）：} $\dfrac{i}{p^2 - m^2 + i\epsilon}$（标量场）
  \item \textbf{顶点：} $-ig$（耦合常数）
  \item \textbf{动量守恒：} 每个顶点处
  \item \textbf{环路积分：} $\displaystyle\int \frac{d^4k}{(2\pi)^4}$
  \item \textbf{对称因子：} 除以图的对称群阶数
\end{enumerate}

\subsection{L4. QED}

\textbf{拉格朗日量：}
\[
\mathcal{L} = \bar{\psi}(i\gamma^\mu D_\mu - m)\psi - \frac{1}{4}F_{\mu\nu}F^{\mu\nu}
\]

精细结构常数：$\alpha = e^2/4\pi \approx 1/137$

\subsection{L5. Wick 定理}

时间序乘积可以展开为所有可能“收缩”之和：
\[
T[\phi(x_1)\phi(x_2)\phi(x_3)\phi(x_4)] = :\phi_1\phi_2\phi_3\phi_4: + \text{所有收缩}
\]

每个完全收缩对应一个 Feynman 图。

\subsection{L6. S 矩阵}

\[
S = T\exp\left(-i\int d^4x\, \mathcal{H}_{\text{int}}(x)\right)
\]

\textbf{散射振幅：} $\langle f|S|i\rangle = \delta_{fi} + i(2\pi)^4\delta^4(p_f - p_i)\mathcal{M}_{fi}$

\subsection{L7. 重整化}

环路积分通常发散。重整化步骤：
\begin{enumerate}
  \item \textbf{正则化：} 引入截断（维度正则化 $d = 4-\epsilon$）
  \item \textbf{重定义参数：} 将发散吸收到物理参数（质量、电荷）的重定义中
  \item \textbf{取极限：} $\epsilon \to 0$，得到有限物理预言
\end{enumerate}

\subsection{L8. 规范固定与鬼场}

规范对称性导致路径积分中对等价构型的无穷重复计数。

\textbf{Faddeev--Popov 方法：}
\[
Z = \int \mathcal{D}A\,\Delta_{FP}[A]\,\delta(G[A])\,e^{iS[A]}
\]

鬼场：反对易的标量场，只出现在内部线。

\subsection{L9. 与渲染的类比}

\begin{center}
\begin{tabular}{ll}
\toprule
\textbf{Feynman 图} & \textbf{渲染} \\
\midrule
树图（无环） & 直接光照（一次弹射） \\
单圈图 & 二次弹射 \\
多圈图 & 多次弹射 \\
耦合常数 $e$ & 反射率 $\rho$ \\
传播子 & 几何衰减 $1/r^2$ \\
顶点 & 散射事件 \\
环路积分 & 对所有中间路径积分 \\
\bottomrule
\end{tabular}
\end{center}
\section{M. 补充：统计力学、热辐射、信息论与离散守恒结构}

\subsection{M1. 统计力学基础}

\textbf{配分函数：}
\[
Z = \sum_i e^{-\beta E_i}, \quad \beta = \frac{1}{k_B T}
\]

\textbf{Boltzmann 分布：} $P_i = \dfrac{e^{-\beta E_i}}{Z}$

\textbf{自由能：} $F = -k_B T \ln Z$

\subsection{M2. 黑体辐射}

\[
B_\lambda(T) = \frac{2hc^2}{\lambda^5} \cdot \frac{1}{e^{hc/(\lambda k_B T)} - 1}
\]

\textbf{Wien 位移定律：} $\lambda_{\max} T = 2.898 \times 10^{-3}$ m$\cdot$K

\subsection{M3. 信息论与熵}

\textbf{Shannon 熵：} $H(X) = -\displaystyle\sum p_i \log p_i$

\textbf{KL 散度：} $D_{KL}(p\|q) = \displaystyle\int p\log\frac{p}{q}\,dx \geq 0$

在 CG 中：最优重要性采样、最大熵正则化。

\subsection{M4. 线性互补问题（LCP）}

找 $z \geq 0$ 使得：
\[
w = Mz + q \geq 0, \quad z^T w = 0
\]

物理意义：接触力 $z$ 和穿透深度 $w$ 不能同时为正。

求解方法：Lemke 算法、投影 Gauss--Seidel (PGS)、内点法。

\subsection{M5. 朗之万方程}

\[
m\ddot{x} = -\gamma\dot{x} + \sqrt{2\gamma k_B T}\,\xi(t) + F(x)
\]

涨落耗散定理：摩擦系数和噪声强度由同一参数 $\gamma$ 决定。

\subsection{M6. 离散外微分（DEC）}

\begin{center}
\begin{tabular}{ll}
\toprule
\textbf{连续} & \textbf{离散} \\
\midrule
0-形式（标量场） & 顶点上的值 \\
1-形式 & 边上的值 \\
2-形式 & 面上的值 \\
外微分 $d$ & 上边界算子（关联矩阵） \\
Hodge 星 $*$ & 对偶网格上的对角矩阵 \\
\bottomrule
\end{tabular}
\end{center}

$d^2 = 0$ 在离散层面精确成立。

\subsection{M7. 变分积分器}

先离散作用量，再对离散作用量取变分：
\[
S_d = \sum_{k=0}^{N-1} L_d(q_k, q_{k+1}, h)
\]

\textbf{离散欧拉--拉格朗日方程：}
\[
D_2 L_d(q_{k-1}, q_k) + D_1 L_d(q_k, q_{k+1}) = 0
\]

性质：自动辛、自动保持动量、长期能量行为好。
\section{N. 核心公式速查表}

\noindent\textbf{这一节是索引，不是第二遍讲解。}每条只给结论，推导在对应篇里——
\begin{itemize}[leftmargin=1.4em,itemsep=0pt]
\item 分析类公式 → \textbf{A}（分析学）、\textbf{C}（概率与测度）；
\item 代数与几何 → \textbf{B}；经典力学 → \textbf{D}；场论 / 量子 → \textbf{E}、\textbf{F}；
\item 数值与图形 → \textbf{G}、\textbf{H}、\textbf{I}、\textbf{J}；
\item 微分方程 9 层 → \textbf{ODE-1}--\textbf{ODE-10}；李群与李代数 → \textbf{Lie-1}--\textbf{Lie-8}。
\end{itemize}
本节内部的「第一层 / 第二层 / 第三层 / 第四层」指\textbf{重要性分层}（承重墙 / 支撑 / 工具 / 延伸），
与 ODE 那条线的「第 n 层」不是一回事，别混。

\subsection{分析}
\[
\int_{\partial M} \omega = \int_M d\omega \qquad \text{（广义 Stokes）}
\]
\[
\|f\|_{L^2}^2 = \sum_n |\hat{f}_n|^2 \qquad \text{（Parseval）}
\]

\subsection{代数}
\[
\exp: \mathfrak{g} \to G, \quad [\mathfrak{g}, \mathfrak{g}] \subset \mathfrak{g}
\]
\[
SU(2) \cong \text{单位四元数}, \quad SU(2) \to SO(3) \text{（2:1 覆盖）}
\]

\subsection{概率}
\[
E[X] = \int_\Omega X\,dP = \int_{\mathbb{R}} x\,dF_X(x)
\]
\[
\text{Var}[\hat{I}_{MC}] = \frac{1}{N}\left(\int \frac{f^2}{p}\,dx - I^2\right)
\]

\subsection{力学}
\[
\delta S = 0 \implies \frac{d}{dt}\frac{\partial L}{\partial \dot{q}} = \frac{\partial L}{\partial q}
\]
\[
\{f, H\} = \dot{f}
\]

\subsection{量子}
\[
i\hbar\frac{\partial}{\partial t}|\psi\rangle = \hat{H}|\psi\rangle
\]
\[
[\hat{x}, \hat{p}] = i\hbar, \quad \Delta x \cdot \Delta p \geq \frac{\hbar}{2}
\]

\subsection{渲染}
\[
L_o = L_e + \int_\Omega f_r\, L_i \cos\theta\, d\omega
\]
\[
w_j = \frac{n_j\, p_j}{\sum_k n_k\, p_k} \qquad \text{（MIS 权重）}
\]

%==============================================================
\section{O. 推荐学习路径}
%==============================================================

\subsection{第一阶段：夯实基础（1--2 个月）}
\begin{enumerate}
  \item 实分析/测度论：Folland《Real Analysis》或 Stein《Real Analysis》
  \item 线性代数深化：Strang + Trefethen《Numerical Linear Algebra》
  \item 概率论严格化：Durrett《Probability: Theory and Examples》
\end{enumerate}

\subsection{第二阶段：核心数学物理（2--3 个月）}
\begin{enumerate}
  \setcounter{enumi}{3}
  \item 微分几何：do Carmo $\to$ Lee《Introduction to Smooth Manifolds》
  \item 经典力学：Arnold《Mathematical Methods of Classical Mechanics》
  \item 泛函分析：Kreyszig《Introductory Functional Analysis with Applications》
\end{enumerate}

\subsection{第三阶段：量子与场论（2--3 个月）}
\begin{enumerate}
  \setcounter{enumi}{6}
  \item 量子力学：Shankar《Principles of Quantum Mechanics》或 Sakurai
  \item 电动力学/场论：Griffiths $\to$ Baez \& Muniain《Gauge Fields, Knots and Gravity》
\end{enumerate}

\subsection{第四阶段：计算与 CG（持续）}
\begin{enumerate}
  \setcounter{enumi}{8}
  \item 数值方法：Trefethen《Spectral Methods in MATLAB》+ 各种 FEM 教材
  \item 渲染：PBRT 源码 + Kajiya 1986 原始论文
  \item 几何处理：Botsch et al.《Polygon Mesh Processing》+ Crane 的离散微分几何课程
  \item 物理模拟：Bridson《Fluid Simulation for Computer Graphics》+ 各种 SIGGRAPH 课程笔记
\end{enumerate}

\clearpage
\subsection{章节与教材对照索引}

读某一节卡住时，按下表直接翻到对应位置，不必通读全书。
（章节代号用笔记内的局部编号：“A1” 指 A 篇第 1 小节，“A1.5” 指其下第 5 小目。）

\begin{center}
\begin{tabular}{L{3.9cm}L{5.2cm}L{5.5cm}}
\toprule
\textbf{本笔记章节} & \textbf{首选教材} & \textbf{对应位置} \\
\midrule
A1 测度与勒贝格积分 & Folland \emph{Real Analysis} & Ch.1--2（\S1.4 外测度、\S2.4 收敛定理） \\
A1.5 $L^p$ 空间 & Folland & Ch.6（\S6.3 插值可略） \\
A2 泛函分析 & Kreyszig & Ch.2--3（赋范/内积）、Ch.7（谱理论） \\
A2.4 谱理论（进阶） & Reed--Simon Vol.1 & Ch.VI--VIII \\
A3 微分形式 & Lee \emph{Smooth Manifolds} & Ch.14（形式）、Ch.16--17（de Rham 与积分） \\
A4 Sobolev 与弱解 & Evans \emph{PDE} + Brezis \emph{Functional Analysis} & Evans Ch.5（Sobolev）、Brezis Ch.8--9 \\
B1 李群与李代数 & Lee \emph{Smooth Manifolds} & Ch.7--8 \\
B1.3 表示论 & Hall \emph{Lie Groups, Lie Algebras} & Ch.3--5 \\
B2 曲线曲面 & do Carmo & Ch.1--2、Ch.4（内蕴几何） \\
B2.3 黎曼几何 & Lee \emph{Riemannian Manifolds} & Ch.3--6 \\
C1--C4 概率的测度论基础 & Durrett & Ch.1--2、Ch.4（条件期望）、Ch.5（鞅） \\
C6--C10 概率深化 & Durrett + Williams \emph{Probability with Martingales} & Durrett Ch.3--4（收敛与 CLT）、Ch.7--8（马尔可夫链、SDE） \\
D1--D3 经典力学 & Arnold & Ch.1--3（含 Noether 定理） \\
E1--E2 场论与微分形式 & Baez \& Muniain & Part I--II \\
F1--F4 量子力学 & Shankar & Ch.1--7、Ch.12（角动量） \\
G1--G5 数值方法 & Trefethen \emph{Spectral Methods} / \emph{Num.\ Lin.\ Alg.} & 谱方法全书；线代 Ch.1--20 \\
H1 离散微分几何 & Botsch et al. \emph{Polygon Mesh Processing} & Ch.2--3 \\
H2--H4 渲染方程 & Kajiya 1986 & \emph{The Rendering Equation}（11 页，必读原文） \\
H4 MIS & Veach 1997 博士论文 & Ch.9（MIS 的出处） \\
I1--I4、I9 流体力学 & Bridson \emph{Fluid Simulation for CG} & 全书（偏实现） \\
K1--K4 路径积分、L1--L8 QFT & Peskin \& Schroeder & Ch.9（路径积分）、Ch.4--7（Feynman 规则） \\
ODE-1--10 微分方程线（ode1--ode4） & 丁同仁《常微分方程教程》 + Evans \emph{PDE} & 丁 Ch.1--4；Evans Ch.2（分类）、Ch.5--6（弱解） \\
Lie-1--8 李群线（lie1--lie3） & Hall & 全书（比 B1 深一档） \\
M6 离散外微分 & Crane 的 DDG 课程笔记 & 前 5 讲 \\
M7 变分积分器 & Hairer--Lubich--Wanner \emph{GNI} & Part I Ch.III（辛积分器） \\
P 计算必修（P1--P9） & Graham--Knuth--Patashnik \emph{Concrete Mathematics} + Hardy--Littlewood--Pólya \emph{Inequalities} & CM Ch.2--7（求和与递推）；Ineq. 全书（不等式） \\
Q 拓扑（点集 $+$ 同调） & Munkres \emph{Topology} + Hatcher & Munkres Ch.2--4、Ch.7；Hatcher Ch.1--2 \\
R 复分析 & Ahlfors \emph{Complex Analysis} + Needham \emph{Visual Complex Analysis} & Ahlfors Ch.4--5（积分与留数）、Ch.6（共形映射） \\
\bottomrule
\end{tabular}
\end{center}

\textbf{精读/略读分级：}
\begin{itemize}
  \item \textbf{必须自己推一遍：} A1.3 三大收敛定理、A1.4 可数可加的必要性、A4 Lax--Milgram 的验证两步、
        B1.2 指数映射与 Rodrigues 公式、C7 特征函数证 CLT、H4 重要性采样方差、R4 留数定理算一个实积分
  \item \textbf{看结论即可：} A2.4 无限维谱定理、C6.3 均匀可积性、E3 闵可夫斯基时空、L7 重整化、M2 黑体辐射、
        Q5 Tychonoff、Q8 Seifert--van Kampen、R5 Riemann 映射定理
  \item \textbf{写代码时再查：} G 全章、H1 离散算子权重、I9 守恒性格式、M4 LCP 求解器、
        Q9 Euler 特征与网格修复、P10 必修优先级总表
\end{itemize}

\clearpage
\subsection{分层导读：按章分配阅读深度}
\label{subsec:reading-depth}

上一节按「进度阶段」排，本节按\textbf{阅读深度}排——同一个问题换一个问法：
\textbf{这些章节，每一章该投入多少？}判据只有一条：\textbf{它与你的目标（PDE $+$ 计算机图形学）之间隔了几层依赖}。

\begin{center}
\begin{tabular}{L{1.7cm}L{12.9cm}}
\toprule
\textbf{层} & \textbf{判据与达标要求} \\
\midrule
承重墙 & 缺了它后面直接走不下去。要求：\textbf{闭卷能推}（不看笔记写出证明要点），能写代码验证。 \\
支撑 & 不理解就讲不清主线的「为什么」。要求：\textbf{能复述结论与用法}，细节允许回查。 \\
工具 & 只需知道「这里有这个」。要求：\textbf{记得索引}，用到时能翻回来。 \\
延伸 & 知道它存在、知道它在哪。要求：知道\textbf{它解决什么问题}，暂不投入。 \\
\bottomrule
\end{tabular}
\end{center}

\subsubsection{第一层 · 承重墙（八项）}

\begin{center}
\begin{tabular}{L{2.5cm}L{12.0cm}}
\toprule
\textbf{章} & \textbf{为什么它在最里层} \\
\midrule
A1 分析基础 & 三大收敛定理回答「极限与积分何时可换」——后面每一个能量估计、每一次取极限都在用它；
Leray--Hopf 存在性直接靠 $L^2$ 完备性（Riesz--Fischer）。 \\
A2 泛函分析 & 没有 Hilbert 空间就没有 Lax--Milgram，没有分布就没有「$\delta$ 源」的严格含义。
这是 PDE 的语法层。 \\
A4 Sobolev 与弱解 & NS 的「解」按定义就是弱解。这一章不是前置知识，而是 L2/L3 的门槛本身。 \\
B2 微分几何 & PDE 的定义域是流形，CG 的网格是离散流形——两条线共用同一套语言，省不掉。 \\
P 计算必修 & 重点在 P3（不等式）与 P9（能量法五步）：P3 是「每次证明都要掏口袋」的工具箱，
P9 是唯一性证明的收尾手法。手不好使，看得懂也推不动。 \\
H 图形学数学 & 目标本身。H1（离散微分几何）、H3（路径空间）、H4（蒙特卡洛）是引擎，其余是外围。 \\
I 流体力学 & NS 目标的正面战场：I1/I2 是方程的来源，I4 涡量动力学是三维湍流的钥匙，
I9 是「离散格式为什么会掉能量」。 \\
Q5 紧性 & 从 Q 章单独提出来：变分法「取极小化序列 $\to$ 提收敛子列 $\to$ 极限是极小元」的第二步
100\% 靠紧性；Rellich--Kondrachov 是椭圆理论的地基。（Q 章其余部分在第二层。） \\
\bottomrule
\end{tabular}
\end{center}

\subsubsection{第二层 · 支撑（十二项）}

\begin{center}
\begin{tabular}{L{2.5cm}L{12.0cm}}
\toprule
\textbf{章} & \textbf{为什么不到承重，又不能不读} \\
\midrule
A3 微分形式 & 它把梯度/旋度/散度三条公式收成一条 $\int_{\partial M}\omega=\int_M d\omega$，
M6 的 DEC 就是它的离散版。可以边用边学，但缺了它 E 章与 DEC 都是黑话。 \\
B1 群论与李群入门 & 旋转代码（四元数、$\exp([\omega]_\times)$ 与 Rodrigues 公式）的直接依据；
深化版在 lie 专题，所以这里只需读懂。 \\
C1--C5 概率测度基础 & 渲染的蒙特卡洛挂在概率空间上；俄式轮盘的无偏性直接就是鞅的定理（C4）。 \\
C6--C9 概率深化 & C6 的四层收敛模式解释「MC 为什么收敛」，C7 给 CLT 的证明；
\textbf{C9 实际上够承重}——它是写渲染器时的误差手册（误差 $\sim\sigma/\sqrt N$，采样翻四倍、噪声减半）。 \\
D 经典力学 & Noether 定理是「对称性 $\to$ 守恒律」的枢纽，从 PDE 到 CG 一路在用；
D3 哈密顿力学是辛积分器的来源。 \\
E 场论与微分形式 & A3 的物理落地。「场」的几何结构（主丛上的联络）看 B2 的纤维丛一节与 E4，
这是理解「场」的最短路径。 \\
F 量子力学 & 它是 A2 谱理论的\textbf{物理直觉来源}——形式化的谱定理在这里有图景。
\textbf{F5 的路径积分与 H3 的路径积分逐项同构}（K5 有对照表），是渲染方程最深的解释。 \\
G 数值方法 & 骨架已有，深化在 de4 与 P9。先把 G1（时间步进）、G2（PDE 离散）、
G3（矩阵分解）的结论记牢。 \\
Q 拓扑学（除 Q5） & Q7 商空间是网格粘合与切割图的语言（球面一刀、环面两刀）；
Q9 同调与 Euler 特征是网格修复与亏格判定的依据。 \\
R 复分析 & 二维 PDE（位势、调和函数）与 CG 共形参数化的共同基础；
但三维以上用不上（Riemann 映射定理无高维对应），所以停在支撑层。 \\
ODE-3（ode1） & 特殊函数与级数（Bessel、Legendre、幂级数解法）是分离变量法的核心，
PDE 里会一次次回来。（第 1--2 层见 ODE-1、ODE-2，第 4 层见 ODE-4。） \\
Lie-1--8 & 对称性方法的完整版；\texttt{lie3} 的「在 CG 中的应用」（旋转、球谐、等变神经网络）
对你的工程线是直接可用的。 \\
\bottomrule
\end{tabular}
\end{center}

\subsubsection{第三层 · 工具（八项）}

\begin{center}
\begin{tabular}{L{2.5cm}L{12.0cm}}
\toprule
\textbf{章} & \textbf{什么时候翻它} \\
\midrule
C7 特征函数 & 要算误差的分布、或想自己独立证一遍 CLT 时才翻。 \\
C8 马尔可夫链与 MCMC & 写 MLT / 采样器时才需要；平时记住「MLT $=$ 路径空间上的 MCMC」即可。 \\
J 光线追踪深层数学 & 做到体积渲染（J5 参与介质）、次表面散射（J6）、焦散（J8）时按标题回查。 \\
M1--M5 补充章 & M4（LCP）做刚体接触求解时必查；M5（朗之万）做粒子噪声时必查；M1/M2/M3 按需。 \\
M7 变分积分器 & 做长时程保能量模拟时查；与 G1 的辛积分器、D3 的哈密顿力学呼应。 \\
P8 欧拉积分 & 它在正态归一化（C2）、PDE 特殊函数、组合估计处冒出来，届时回查。 \\
N 速查表 & 想不起某条公式时先翻这里，翻不到再回原章。 \\
O 学习路径 $+$ §15.5 & 导航本身：卡在某节时先查教材对照表，再决定翻哪本书。 \\
\bottomrule
\end{tabular}
\end{center}

\subsubsection{第四层 · 延伸（七项）}

\begin{center}
\begin{tabular}{L{2.5cm}L{12.0cm}}
\toprule
\textbf{章} & \textbf{为什么现在不投入} \\
\midrule
ODE-1--2 & 考研/高数级 ODE 套路与系统理论——你的底子够用，需要时回查即可。 \\
ODE-4 & 李群方法解 ODE 很漂亮，但用得少；知道「对称性怎么找」就够。 \\
C10 Itô 微积分 & 除非要做\textbf{随机流体/随机 NS}，先记住一句：
「SDE 与抛物型 PDE 是同一件事的两种语言」。 \\
K 路径积分专题 & 与渲染的逐项类比极漂亮（K5），但不是工程必需。 \\
L QFT 微扰论 & 同理；L7 重整化是思想性收益，与你的路线接口弱。 \\
E3--E4 时空与规范场 & 物理兴趣向，与 PDE／CG 的接口很弱。 \\
M6 离散外微分 & 若你确定走\textbf{几何处理}方向，把它提到支撑层；否则先知道它在 M6。 \\
\bottomrule
\end{tabular}
\end{center}

\textbf{怎么用这张表：}
\begin{itemize}
  \item \textbf{时间分配}按层：承重墙占七成，支撑占两成，工具与延伸共占一成（延伸层现在是零）。
  \item \textbf{层内按依赖走}，不要按字母顺序：P 计算必修 $\to$ A1 $\to$ A2 $\to$ A3 $\to$ A4 $\to$ B2 $\to$ Q $\to$ R，
        这条链走完，L2/L3 的地基就齐了；H 与 I 可以与之并行推进（它们是目标，不是前置）。
  \item \textbf{例外比标签重要}：Q5（紧性）、C9（方差缩减）、de1 第三层（特殊函数）、
        F5（路径积分）在各自章里都够承重级——章级标签只给整体印象，不要因为标签而跳过它们。
\end{itemize}

\section{P. 计算必修：从平方和到欧拉积分}
\label{sec:computation}

\textbf{定位：} 前面的 A--O 讲的是「概念地图」，本章补的是「一副会算的手」。概念要配得上一副算得动的手，
否则遇到 $\int_0^\infty e^{-x^2}dx$ 只能干看。本章服务所有层：A 分析（放缩、级数）、A2 泛函（$L^p$ 估计）、
I 流体（能量法）、H 图形（旋转、复数、线代）。

\textbf{标记：} 【必】= 必须肌肉记忆（闭卷 30 秒内写出或推完）；【熟】= 熟练（查表能立刻上手）；
【知】= 知道出处即可。

\textbf{代号说明：} 本章九组原记作 G1--G9。因 \textbf{G. 数值方法} 已占用 G1--G5，
在此统一改记为 \textbf{P1--P9}，与黑板里的旧编号一一对应（P$n$ 即旧 G$n$）。

\begin{center}
\begin{tabular}{L{0.9cm}L{3.7cm}L{5.9cm}L{3.5cm}}
\toprule
\textbf{组} & \textbf{主题} & \textbf{代表内容} & \textbf{主要服务} \\
\midrule
P1 & 求和与数列 & $\sum k$、$\sum k^2$、$\sum k^3$、裂项、错位、递推 & 分析、级数 \\
P2 & 恒等式与变形 & 乘法公式、配方、部分分式、韦达 & 积分、Laplace \\
P3 & 不等式与放缩 & AM--GM、Cauchy--Schwarz、Young、Hölder、Gronwall & \textbf{一切证明} \\
P4 & 数列与极限 & 特征根法、单调有界、等价无穷小、Taylor、洛必达 & 分析入门 \\
P5 & 三角与复数 & 和角、积化和差、Euler 公式、单位根 & Fourier、CG 旋转 \\
P6 & 线性代数手算 & 行列式、消元、特征值、二次型、Gram--Schmidt & A2、CG \\
P7 & 微积分计算 & 换元、分部、Jacobian、Maclaurin 表 & 全场 \\
P8 & 欧拉积分与特殊函数 & 高斯积分、$\Gamma$、Beta、Stirling & 概率、PDE、量子 \\
P9 & 估计与渐近 & 移导数、能量法流水线、大 $O$ & I 流体、NS \\
\bottomrule
\end{tabular}
\end{center}

\textbf{读法：} 先吃透标【必】的条目（硬核心），【熟】留到用时回查。

%--------------------------------------------------------------
\subsection{P1. 求和与数列}

\subsubsection{核心思想：差分法（telescoping）}

想求 $\sum_k f(k)$，就找一个 $F$ 使 $f(k) = F(k+1) - F(k)$，整列求和只剩首尾相消。
\textbf{这是全部求和公式的统一来源}——不是背公式，而是找那个 $F$。

\subsubsection{必备公式}

\begin{center}
\begin{tabular}{L{3.2cm}L{6.2cm}L{5.2cm}}
\toprule
\textbf{公式} & \textbf{差分/证明} & \textbf{直觉} \\
\midrule
$\displaystyle\sum_{k=1}^n k=\frac{n(n+1)}{2}$ & 高斯配对：$(1+n)+(2+n-1)+\cdots$ & 阶梯面积 $=$ 半个矩形 \\
$\displaystyle\sum_{k=1}^n k^2=\frac{n(n+1)(2n+1)}{6}$ &
差分为 $F(k+1)-F(k)$，$F(k)=\frac{k(k-1)(2k-1)}{6}$，望远镜即得 &
抛物线下的面积 $\approx n^3/3$ 加修正项 \\
$\displaystyle\sum_{k=1}^n k^3=\Big[\frac{n(n+1)}{2}\Big]^2$ &
$F(k)=\frac{k^2(k-1)^2}{4}$ & 巧合之美：立方和 $=$ 和的平方 \\
\bottomrule
\end{tabular}
\end{center}

\subsubsection{推导示范（$\sum k^2$，必须会复述）}

对恒等式
\[
(k+1)^3-k^3 = 3k^2+3k+1
\]
从 $k=1$ 加到 $n$：左边望远镜 $=(n+1)^3-1$，右边 $=3\sum_k k^2+3\sum_k k+n$，解出
\[
\sum_{k=1}^n k^2=\frac{(n+1)^3-1-3\cdot\frac{n(n+1)}{2}-n}{3}=\frac{n(n+1)(2n+1)}{6}.
\]
\textbf{套路：升一阶的差，凑出想求的幂。}（求 $\sum k^3$ 就用 $(k+1)^4-k^4$。）

\subsubsection{常用技巧}

\begin{itemize}
  \item \textbf{裂项相消}（【必】）：$\dfrac{1}{k(k+1)}=\dfrac1k-\dfrac1{k+1}\Rightarrow\sum=\dfrac{n}{n+1}$；
        更复杂的用待定系数拆成同类差。
  \item \textbf{等比求和}（【必】）：$\sum_{k=0}^{n-1}q^k=\dfrac{1-q^n}{1-q}$；$|q|<1$ 时无穷版 $=\dfrac{1}{1-q}$。
  \item \textbf{错位相减}（【熟】）：等差 $\times$ 等比（如 $\sum k q^k$）——写 $S$，乘 $q$ 错位相减，化归等比。
  \item \textbf{二项式定理}（【必】）：$(1+x)^n=\sum_k\binom{n}{k}x^k$（$n$ 个小括号各挑 $x$ 或 $1$ 的计数）；
        推论 $\sum_k\binom{n}{k}=2^n$。
  \item \textbf{组合和技巧}（【熟】）：$k\binom{n}{k}=n\binom{n-1}{k-1}$（先挑一个当队长）。
  \item \textbf{$\sum \leftrightarrow \int$}（【熟】）：$\dfrac1n\sum_k f\!\big(\tfrac kn\big)\to\int_0^1 f$；
        反向用于估计：$f$ 递减非负时 $\int_1^{n+1}f\le\sum_{k=1}^n f(k)\le f(1)+\int_1^n f$。
\end{itemize}

%--------------------------------------------------------------
\subsection{P2. 恒等式与代数变形}

\subsubsection{乘法公式族（【必】）}

\begin{center}
\begin{tabular}{ll}
\toprule
\textbf{公式} & \textbf{用途备注} \\
\midrule
$(a\pm b)^2=a^2\pm 2ab+b^2$ & 一切配方的基础 \\
$(a+b)(a-b)=a^2-b^2$ & 有理化、差分分解 \\
$(a\pm b)^3=a^3\pm3a^2b+3ab^2\pm b^3$ & 二项式 $n=3$ 特例 \\
$a^3\pm b^3=(a\pm b)(a^2\mp ab+b^2)$ & 立方和差 \\
$(a+b+c)^2=a^2+b^2+c^2+2ab+2bc+2ca$ & 向量 $\|\mathbf{a}+\mathbf{b}\|^2$ 的代数版 \\
\bottomrule
\end{tabular}
\end{center}

\subsubsection{配方（【必】）}

\[
ax^2+bx+c=a\Big(x+\frac{b}{2a}\Big)^2+\frac{4ac-b^2}{4a}
\]

\textbf{直觉：} 平移原点让抛物线对称——配方就是换元 $x\to x+\frac{b}{2a}$。

\textbf{用途：} 求顶点/最值；推导求根公式；二次型正定性（见 P6）；
把 $\int\frac{dx}{x^2+bx+c}$ 化为 $\arctan$ 或 $\log$。\textbf{类比：一元情形下的「对角化」。}

\subsubsection{因式分解与多项式（【熟】）}

\begin{itemize}
  \item 十字相乘、分组分解、试根法（有理根定理）；多项式长除／综合除法。
  \item \textbf{韦达定理}（【必】）：$x_1+x_2=-\dfrac ba$，$x_1x_2=\dfrac ca$；$n$ 次推广见【熟】。
  \item 对称式与轮换式：知道 $a^2+b^2+c^2=(a+b+c)^2-2(ab+bc+ca)$ 这类转换。
\end{itemize}

\subsubsection{部分分式分解（【必】· 积分／Laplace／传递函数全靠它）}

\[
\frac{1}{(x-a)(x-b)}=\frac{A}{x-a}+\frac{B}{x-b},\qquad A=\frac{1}{a-b}\ (\text{代 }x=a)
\]

\textbf{直觉类比：} 像把 $\frac16=\frac12-\frac13$ 拆开——把复杂分式拆成「基本块」，
每块积分都是 $\log$ 或 $\arctan$。

\textbf{四类基本块}（连原函数一起记）：

\begin{center}
\begin{tabular}{lll}
\toprule
\textbf{块} & \textbf{原函数} & \textbf{条件} \\
\midrule
$\dfrac{1}{x-a}$ & $\ln|x-a|$ & 单实根 \\
$\dfrac{1}{(x-a)^k}$ & $\dfrac{(x-a)^{1-k}}{1-k}$ & $k\ge2$ \\
$\dfrac{1}{x^2+px+q}$ & $\arctan$ 型（配方后） & 复根，判别式 $<0$ \\
$\dfrac{2x+p}{x^2+px+q}$ & $\ln(x^2+px+q)$ & 同上（分子是分母的导数） \\
\bottomrule
\end{tabular}
\end{center}

\textbf{目标：} 任何有理分式 $=$ 多项式 $+$ 四类基本块的组合。

\subsubsection{根式、指数与对数（【必】）}

\begin{itemize}
  \item 有理化：乘共轭 $(\sqrt a+\sqrt b)(\sqrt a-\sqrt b)=a-b$；双重根式 $\sqrt{a\pm 2\sqrt b}$ 配方化简。
  \item 指数对数：$a^xa^y=a^{x+y}$、$\ln(xy)=\ln x+\ln y$、换底 $\log_b x=\dfrac{\ln x}{\ln b}$。
  \item 常用界：$\ln 2\approx0.693<1$、$e\approx2.718>2$。
\end{itemize}

%--------------------------------------------------------------
\subsection{P3. 证明需要的核心不等式}

这一组是分析与能量估计「每次证明都要掏口袋」的工具。每条给一行证明加一句直觉。

\subsubsection{基础四件套（【必】）}

\begin{enumerate}
  \item \textbf{AM--GM}：$\dfrac{a+b}{2}\ge\sqrt{ab}$（$a,b\ge0$）。
        证明：$(\sqrt a-\sqrt b)^2\ge0$ 展开即得。
        \textbf{直觉：周长固定的矩形中正方形面积最大。}
        $n$ 元版（【熟】）：$\dfrac{x_1+\cdots+x_n}{n}\ge(x_1\cdots x_n)^{1/n}$。
  \item \textbf{Cauchy--Schwarz}：$\Big(\sum_i a_ib_i\Big)^2\le\Big(\sum_i a_i^2\Big)\Big(\sum_i b_i^2\Big)$。
        证明：$\sum_i(a_i-tb_i)^2\ge0$ 对一切 $t$ 成立，作为 $t$ 的二次式其判别式 $\le0$。
        \textbf{直觉：投影 $\le$ 全长}；取等 $\iff$ 两向量共线。
        积分版：$\big(\int fg\big)^2\le\int f^2\int g^2$——即 $L^2$ 内积的 Cauchy 不等式（见 A2）。
  \item \textbf{Young 不等式}：$ab\le\dfrac{a^p}{p}+\dfrac{b^q}{q}$（$a,b\ge0$，$\frac1p+\frac1q=1$）。
        证明：$\ln$ 是凹函数，
        $\ln(ab)=\frac1p\ln a^p+\frac1q\ln b^q\le\ln\!\Big(\frac{a^p}{p}+\frac{b^q}{q}\Big)$。
        \textbf{直觉：Cauchy--Schwarz 的「指数家族」版}；$p=q=2$ 时退化为 $2ab\le a^2+b^2$。
        这是 \textbf{PDE 能量估计的头号工具}。
  \item \textbf{三角不等式}：$|a+b|\le|a|+|b|$（向量同理：$\|\mathbf{x}+\mathbf{y}\|\le\|\mathbf{x}\|+\|\mathbf{y}\|$）。
        \textbf{直觉：直走不比绕路远。}
\end{enumerate}

\subsubsection{放缩盒子（【必】／【熟】）}

\begin{center}
\begin{tabular}{ll}
\toprule
\textbf{不等式} & \textbf{出处／用法} \\
\midrule
$e^x\ge 1+x$（【必】） & 凸函数切线在图形下方 \\
$\ln(1+x)\le x$（【必】） & 上一式的另一面（取对数） \\
$\ln x\le x-1$（【必】） & 换元即得 \\
$\dfrac{x}{1+x}\le\ln(1+x)\le x$（【熟】） & 靠近 $0$ 的二阶精确版 \\
$|\sin x|\le|x|$ 与 $|\sin x|\le1$（【必】） & 近 $0$ 用前者，远处用后者 \\
$\big(1+\frac1n\big)^n\uparrow e$、$n^{1/n}\to1$（【必】） & 常用极限，证明见 P4 \\
$\sum_{k>n}2^{-k}=2^{-n}$（【必】） & 几何级数尾，$\varepsilon/2^k$ 技巧的弹药 \\
\bottomrule
\end{tabular}
\end{center}

\subsubsection{进阶（【熟】· 记结论 $+$ 用法）}

\begin{itemize}
  \item \textbf{Jensen}：$\varphi$ 凸 $\implies \varphi(\mathbb E X)\le\mathbb E\,\varphi(X)$。
        \textbf{直觉：凸函数的弦在图形上方。}
  \item \textbf{Hölder}：$\int|fg|\le\|f\|_p\|g\|_q$（$\frac1p+\frac1q=1$）；Cauchy--Schwarz 是其 $p=q=2$ 特例。
  \item \textbf{Gronwall}（【必】记结论）：$u'\le\beta u+f$，$u(0)=u_0$ $\implies$
        \[
        u(t)\le e^{\beta t}u_0+\int_0^t e^{\beta(t-s)}f(s)\,ds
        \]
        \textbf{证明：} 乘积分因子 $e^{-\beta t}$ 后两边积分。
        \textbf{直觉：微分不等式 $\to$ 积分不等式的出厂设置}；NS 唯一性证明的收尾工具。
\end{itemize}

%--------------------------------------------------------------
\subsection{P4. 数列与极限计算}

\subsubsection{数列}

\begin{itemize}
  \item 等差／等比／混合数列估值。
  \item \textbf{线性递推特征根法}（【熟】）：$a_n=pa_{n-1}+qa_{n-2}$——试 $a_n=r^n$ 得特征方程 $r^2=pr+q$；
        两实单根时通解 $Ar_1^n+Br_2^n$；重根时 $(A+Bn)r^n$。
        \textbf{直觉：递推 $=$ 矩阵幂}——$r$ 是转移矩阵的特征值，长期行为由主导特征值决定
        （Fibonacci：$a_n\approx\varphi^n/\sqrt5$）。
  \item \textbf{单调有界定理}（【必】）：单调 $+$ 有界 $\implies$ 收敛；再对递推式两边取极限（不动点）解出极限值。
        \textbf{例：} $a_1=\sqrt2$，$a_{n+1}=\sqrt{2+a_n}$：单调增且 $<2$，
        极限 $L$ 满足 $L=\sqrt{2+L}\Rightarrow L=2$。
\end{itemize}

\subsubsection{极限三件套（【必】）}

\begin{enumerate}
  \item \textbf{等价无穷小}（$x\to0$）：$\sin x\sim x$、$\tan x\sim x$、$1-\cos x\sim\frac{x^2}2$、
        $\ln(1+x)\sim x$、$e^x-1\sim x$、$(1+x)^\alpha-1\sim\alpha x$。
        \textbf{直觉：取 Taylor 首项}；\textbf{只适用于乘除结构，加减结构慎用}。
  \item \textbf{Taylor 展开}：按 P7 的 Maclaurin 表展开到够用的阶，余项估计见 P9。
  \item \textbf{洛必达}（【熟】）：$\frac00$ 或 $\frac\infty\infty$ 型。
        \textbf{警告：先化简／用等价无穷小，洛必达是最后手段}；数列极限先化成函数再用。
\end{enumerate}

\subsubsection{常用极限（【必】）}

\[
\lim_{n\to\infty}\Big(1+\frac1n\Big)^n=e,\qquad
\lim_{n\to\infty}n^{1/n}=1,\qquad
\lim_{n\to\infty}\Big(1+\frac xn\Big)^n=e^x
\]

前两个的证明路线：第一个用二项式展开逐项比较证明单调有界；
第二个取对数后用 $\ln t\le t-1$ 放缩。

%--------------------------------------------------------------
\subsection{P5. 三角与复数}

\subsubsection{三角公式}

\begin{center}
\begin{tabular}{L{7.4cm}L{6.8cm}}
\toprule
\textbf{公式} & \textbf{直觉／来源} \\
\midrule
$\sin(A\pm B),\ \cos(A\pm B)$（【必】） & \textbf{旋转矩阵的乘法}（和角 $=$ 两个旋转叠加） \\
$\sin 2x=2\sin x\cos x$，$\cos2x=1-2\sin^2x$（【必】） & 和角特例 \\
辅助角：$a\sin\theta+b\cos\theta=R\sin(\theta+\varphi)$，$R=\sqrt{a^2+b^2}$（【必】） &
向量 $(a,b)$ 的极坐标；两个同频振动合成 \\
积化和差／和差化积（【熟】） & 和角两式相加减即得 \\
\bottomrule
\end{tabular}
\end{center}

\subsubsection{Euler 公式与复数（【必】）}

\[
e^{i\theta}=\cos\theta+i\sin\theta
\]

\begin{itemize}
  \item \textbf{直觉：指数是旋转的引擎}——单位圆上的匀速圆周运动。
        严格证明：把 $e^{i\theta}$ 定义为幂级数，实部／虚部分别就是 $\cos,\sin$ 的级数。
  \item \textbf{De Moivre}：$(\cos\theta+i\sin\theta)^n=\cos n\theta+i\sin n\theta$（幂的辐角 $n$ 倍）。
  \item \textbf{复数乘法 $=$ 旋转 $+$ 缩放}：模相乘、辐角相加。
        \textbf{CG 的旋转矩阵／四元数全从这里长出来}（与 B1 的 $SU(2)$ 一节合看）。
  \item \textbf{单位根}：$z^n=1$ 的解为 $e^{2\pi ik/n}$，$n$ 个均匀分布在单位圆上——DFT／Fourier 的骨架。
\end{itemize}

\subsubsection{三角正交性（【熟】· Fourier 前置）}

\[
\int_0^{2\pi}\sin(m\theta)\sin(n\theta)\,d\theta=\pi\delta_{mn},\qquad
\int_0^{2\pi}\sin(m\theta)\cos(n\theta)\,d\theta=0
\]

用积化和差一句话即可证明。这就是 A2「正交基与 Fourier 分析」里 $\{e^{inx}\}$ 正交性的实数版。

\subsubsection{常用三角积分（【必】）}

\[
\int\sin^2x\,dx=\frac x2-\frac{\sin2x}{4},\qquad
\int\cos^2x\,dx=\frac x2+\frac{\sin2x}{4},\qquad
\int\sin x\cos x\,dx=\frac{\sin^2x}{2}
\]

%--------------------------------------------------------------
\subsection{P6. 线性代数手算}

\subsubsection{行列式（【必】）}

\begin{itemize}
  \item $2\times2$：$\det=ad-bc$；$3\times3$：对角线法则（Sarrus）；$4\times4$ 起用展开／消元（【熟】）。
  \item 性质：换行变号；某行乘 $\lambda$ 则行列式乘 $\lambda$；行加不变；三角阵 $=$ 对角元之积。
  \item \textbf{Vandermonde}（【熟】）：$\det=\prod_{i<j}(x_j-x_i)$（记结论 $+$ 会 $3\times3$ 手推）。
  \item \textbf{几何直觉：} 行列式 $=$ 面积／体积的伸缩率——
        即 A1 里的 $\mu(T(A))=|\det T|\,\mu(A)$（此处 $\mu$ 是 Lebesgue 测度）。
\end{itemize}

\subsubsection{消元与求逆（【必】）}

\begin{itemize}
  \item 高斯消元 $\to$ 阶梯形 $\to$ 回代；LU 分解（消元过程的记录）。
  \item $2\times2$ 求逆：
        $\begin{pmatrix}a&b\\c&d\end{pmatrix}^{-1}=\dfrac{1}{ad-bc}\begin{pmatrix}d&-b\\-c&a\end{pmatrix}$。
  \item $3\times3$：增广 $[A\mid I]$ 消到 $[I\mid A^{-1}]$（【熟】）。
  \item 可逆 $\iff\det\ne0\iff$ 满秩 $\iff$ 列向量线性无关。
\end{itemize}

\subsubsection{特征值与二次型（【必】／【熟】）}

\begin{itemize}
  \item 特征值：解 $\det(A-\lambda I)=0$ 得 $\lambda$，再解 $(A-\lambda I)v=0$（$2\times2$／$3\times3$ 手算必须熟）。
  \item 二次型配方：$\mathbf{x}^TA\mathbf{x}$ 配成平方和（与 P2 配方同源）。
  \item \textbf{正定性三等价}（【必】记结论）：特征值全正 $\iff$ 顺序主子式全正（Sylvester）
        $\iff$ 存在 Cholesky 分解 $A=LL^T$。
  \item \textbf{直觉：正定 $=$ 「碗」}；特征值 $=$ 各主方向的伸缩率——
        与 A2 谱定理、H1 的 Laplace 谱是同一件事的三个层次。
\end{itemize}

\subsubsection{正交化（【熟】）}

\begin{itemize}
  \item Gram--Schmidt：$v-\dfrac{\langle v,u\rangle}{\langle u,u\rangle}u$，逐次正交。
  \item QR：$A=QR$；小 $2\times2$ 手算一次即可。这就是 A2「正交分解定理」的算法版。
  \item SVD 迷你体验：$2\times2$ 对称阵手算 $\sigma_i=\sqrt{\lambda_i(A^TA)}$。
\end{itemize}

%--------------------------------------------------------------
\subsection{P7. 微积分计算技巧}

\subsubsection{必背 Maclaurin 表（【必】）}

\begin{center}
\begin{tabular}{lll}
\toprule
\textbf{函数} & \textbf{展开} & \textbf{收敛半径} \\
\midrule
$e^x$ & $1+x+\frac{x^2}{2!}+\frac{x^3}{3!}+\cdots$ & $\infty$ \\
$\sin x$ & $x-\frac{x^3}{3!}+\frac{x^5}{5!}-\cdots$ & $\infty$ \\
$\cos x$ & $1-\frac{x^2}{2!}+\frac{x^4}{4!}-\cdots$ & $\infty$ \\
$\ln(1+x)$ & $x-\frac{x^2}{2}+\frac{x^3}{3}-\cdots$ & $1$ \\
$\frac{1}{1-x}$ & $1+x+x^2+\cdots$ & $1$ \\
$(1+x)^\alpha$ & $1+\alpha x+\frac{\alpha(\alpha-1)}{2!}x^2+\cdots$ & $1$（$\alpha\notin\mathbb N$） \\
$\arctan x$ & $x-\frac{x^3}{3}+\frac{x^5}{5}-\cdots$ & $1$ \\
\bottomrule
\end{tabular}
\end{center}

\subsubsection{积分技术（【必】）}

\begin{itemize}
  \item \textbf{换元}：凑微分 $f'(x)g(f(x))$；三角替换 $\sqrt{a^2-x^2}\to a\sin t$、
        $\sqrt{x^2+a^2}\to a\tan t$、$\sqrt{x^2-a^2}\to a\sec t$（【熟】）。
  \item \textbf{分部积分}：$\int u\,dv=uv-\int v\,du$；表格法（多项式 $\times$ 三角／指数）；
        循环积分（$\int e^x\sin x$ 型）。
        \textbf{直觉：乘积法则的积分版；「移导数」像搬家}——把导数从难搞的因子搬到好搞的因子
        （P9 能量法的同一手法）。
  \item \textbf{部分分式积分}：拆成 P2 的四类基本块后逐块积。
  \item \textbf{对称性}：奇函数在对称区间积分为 $0$；偶函数 $=2\times$ 半区间；周期函数平移对齐。
  \item 万能替换 $t=\tan\frac x2$（【熟】· 三角有理式的保险手段）。
\end{itemize}

\subsubsection{常用原函数小表（【必】）}

\[
\int\frac{dx}{1+x^2}=\arctan x,\qquad
\int\frac{x\,dx}{1+x^2}=\frac12\ln(1+x^2),\qquad
\int\tan x\,dx=-\ln|\cos x|
\]
\[
\int\ln x\,dx=x\ln x-x,\qquad
\int xe^{-x}dx=-(x+1)e^{-x},\qquad
\int\frac{dx}{x^2-1}=\frac12\ln\Big|\frac{x-1}{x+1}\Big|
\]

\subsubsection{多重积分（【熟】）}

\begin{itemize}
  \item 换元公式：$\displaystyle\iint f\,dx\,dy=\iint f\,|J|\,du\,dv$，
        $J=\dfrac{\partial(x,y)}{\partial(u,v)}$（Jacobian 行列式）——与 P6 的行列式是同一条。
  \item 极坐标 $J=r$；球坐标 $J=r^2\sin\theta$（【必】记，用一次推一次）。
  \item 交换积分次序的直觉 $\to$ Fubini--Tonelli（A1 收敛定理表的成员）。
\end{itemize}

\subsubsection{微分（【必】）}

链式／乘积／商法则熟练；隐函数求导与对数求导法（【熟】）；
多变量：梯度 $\nabla f$、方向导数、全微分；Hessian 与凸性判定（【熟】）。

%--------------------------------------------------------------
\subsection{P8. 欧拉积分与特殊函数}

\subsubsection{高斯积分（【必】· 全分析最有名的技巧）}

\[
\int_{-\infty}^{\infty}e^{-x^2}\,dx=\sqrt{\pi}
\]

\textbf{完整证明（会复述）：}
\[
I^2=\iint_{\mathbb R^2}e^{-(x^2+y^2)}\,dx\,dy
\ \xrightarrow{\ \text{极坐标}\ }\
\int_0^{2\pi}\!\!d\theta\int_0^\infty e^{-r^2}r\,dr=2\pi\cdot\frac12=\pi
\]
\textbf{直觉：「平方换维度」}——一维算不动，就升到二维，让 $r^2$ 恰好配出 $r\,dr$ 的换元。
（这个积分在 C2 的正态分布密度归一化里会再出现一次。）

\subsubsection{$\Gamma$ 函数（Euler 第二积分）（【必】）}

\[
\Gamma(s)=\int_0^\infty t^{s-1}e^{-t}\,dt\qquad(\operatorname{Re}s>0)
\]

\begin{itemize}
  \item \textbf{递推}（分部积分一行）：$\Gamma(s+1)=s\,\Gamma(s)$ $\implies$ $\Gamma(n+1)=n!$。
        即 \textbf{$\Gamma$ 是阶乘的连续插值}。
  \item \textbf{$\Gamma(1/2)=\sqrt\pi$}（【必】会推）：换元 $t=x^2$，
        \[
        \Gamma\Big(\frac12\Big)=\int_0^\infty t^{-1/2}e^{-t}dt=2\int_0^\infty e^{-x^2}dx=\sqrt\pi.
        \]
  \item 常用：$\int_0^\infty x^n e^{-x}dx=n!$；$\Gamma(3/2)=\frac{\sqrt\pi}{2}$。
  \item \textbf{直觉：} $e^{-t}$ 是无穷长的衰减，$t^{s-1}$ 是幂权重；$\Gamma$ 把离散的阶乘「抹平」到连续参数。
\end{itemize}

\subsubsection{Beta 函数（Euler 第一积分）（【熟】）}

\[
B(x,y)=\int_0^1 t^{x-1}(1-t)^{y-1}\,dt
\]

\begin{itemize}
  \item 对称：$B(x,y)=B(y,x)$；与 $\Gamma$ 的关系（【必】记）：
        $B(x,y)=\dfrac{\Gamma(x)\Gamma(y)}{\Gamma(x+y)}$。
  \item 用途：$\displaystyle\int_0^{\pi/2}\sin^m\theta\cos^n\theta\,d\theta
        =\frac12 B\Big(\frac{m+1}{2},\frac{n+1}{2}\Big)$（【熟】）。
  \item \textbf{直觉：} 单位区间上的「双参数权重」；与 $\Gamma$ 的关系来自把 $e^{-(u+v)}$ 的乘积换元拆解。
\end{itemize}

\subsubsection{Stirling 公式（【熟】）}

\[
n!\ \sim\ \sqrt{2\pi n}\,\Big(\frac ne\Big)^n\qquad(n\to\infty)
\]

\textbf{直觉：} $\ln n!\approx\int_1^n\ln x\,dx=n\ln n-n$，再补 $\sqrt{2\pi n}$ 修正。
\textbf{用途：} 组合估计（$\binom{2n}{n}\approx 4^n/\sqrt{\pi n}$）、大数极限、C5 的概率渐近。

%--------------------------------------------------------------
\subsection{P9. 估计与渐近}

\subsubsection{分部积分移导数（【必】· 能量法的发动机）}

\[
\int_\Omega u'v=-\int_\Omega uv'+\big[uv\big]_{\partial\Omega}
\]

\textbf{手法：} 想把导数从 $u$ 搬到 $v$，就分部一次；零边界条件时边界项消失。

\subsubsection{能量法流水线（【必】· 必须会走全流程）}

以热方程 $u_t=u_{xx}$（零边界）为例：

\begin{enumerate}
  \item 两边乘 $u$ 并积分：$\frac12\frac{d}{dt}\int u^2=\int u\,u_{xx}$；
  \item 分部积分（移导数，边界项消失）：$=-\int u_x^2$；
  \item 得衰减 $\frac{d}{dt}\|u\|^2=-2\|u_x\|^2\le0$；
  \item 若出现非线性项，用\textbf{吸收不等式} $ab\le\varepsilon a^2+C_\varepsilon b^2$ 把它「吸收」进左边
        （选 $\varepsilon$ 让系数合拢）；
  \item 最后用 \textbf{Gronwall} 收尾得显式界。
\end{enumerate}

\textbf{NS 的 Leray--Hopf 能量不等式就是这条流水线的产物}——
第 1 步「两边乘 $u$ 再积分」与第 2 步「移导数」的掉能量现象，对应 A1 里 Fatou 引理那个单边不等式。

\subsubsection{渐近记号（【必】）}

\begin{itemize}
  \item 大 $O$：$f=O(g)\iff|f|\le C|g|$（最终成立）；小 $o$：$f/g\to0$；$\sim$：比值 $\to1$。
  \item Taylor 余项估计：$|R_n|\le C|x|^{n+1}$（带界形式，够用）。
  \item 常用：$e^x=1+x+O(x^2)$、$\ln(1+x)=x+O(x^2)$、$\sin x=x+O(x^3)$。
\end{itemize}

\subsubsection{决策树：证明里该掏哪个不等式？（【熟】）}

\begin{center}
\begin{tabular}{L{6.4cm}L{8.2cm}}
\toprule
\textbf{情形} & \textbf{武器} \\
\midrule
乘积 $\to$ 求和分离 & Young（$ab\le\frac{a^p}{p}+\frac{b^q}{q}$）或 $2ab\le a^2+b^2$ \\
和／积的绝对值分离 & Cauchy--Schwarz（$L^2$）／Hölder（$L^p$） \\
微分不等式 $\to$ 显式界 & Gronwall \\
非线性项要「吃掉」 & 吸收不等式（带 $\varepsilon$ 的 Young） \\
凸性／平均值 & Jensen \\
\bottomrule
\end{tabular}
\end{center}

%--------------------------------------------------------------
\subsection{P10. 必修优先级总表}

目标：每个【必】条目练到\textbf{闭卷 30 秒内写出或推完}；【熟】条目查表能立刻上手。

\begin{center}
\begin{tabular}{L{2.6cm}L{9.8cm}L{2.0cm}}
\toprule
\textbf{组} & \textbf{条目} & \textbf{级别} \\
\midrule
P1 & $\sum k,\sum k^2,\sum k^3$ 闭卷推导 & 【必】 \\
P1 & 裂项、等比、错位相减、二项式定理 & 【必】 \\
P1 & $\sum\leftrightarrow\int$ 估计 & 【熟】 \\
P2 & 乘法公式、配方、韦达 & 【必】 \\
P2 & 部分分式（四类基本块 $+$ 待定系数） & 【必】 \\
P2 & 有理化与常用界 & 【必】 \\
P3 & AM--GM、Cauchy--Schwarz、Young、三角不等式 & 【必】 \\
P3 & 放缩盒子（$e^x$、$\ln$、$|\sin x|$、级数尾） & 【必】 \\
P3 & Jensen、Hölder & 【熟】 \\
P3 & Gronwall（结论 $+$ 用法） & 【必】 \\
P4 & 递推特征根法、单调有界 $+$ 不动点 & 【必】 \\
P4 & 等价无穷小、Taylor、洛必达 & 【必】 \\
P4 & $(1+1/n)^n$、$n^{1/n}$ & 【必】 \\
P5 & 和角、倍角、辅助角、常用三角积分 & 【必】 \\
P5 & Euler 公式、De Moivre、单位根 & 【必】 \\
P6 & $2\times2$／$3\times3$ 行列式与求逆、Vandermonde & 【必】 \\
P6 & 手算特征值、二次型配方、正定性三等价 & 【必】 \\
P6 & Gram--Schmidt、QR／SVD 小例 & 【熟】 \\
P7 & 换元、分部（表格法）、部分分式、对称性 & 【必】 \\
P7 & Maclaurin 表七条闭卷默写 & 【必】 \\
P7 & Jacobian、极坐标／球坐标 & 【熟】 \\
P8 & 高斯积分完整证明会复述 & 【必】 \\
P8 & $\Gamma$ 递推、$\Gamma(1/2)=\sqrt\pi$ 会推 & 【必】 \\
P8 & Beta 与 Stirling（结论 $+$ 用法） & 【熟】 \\
P9 & 移导数、能量法五步、吸收不等式 & 【必】 \\
P9 & 大 $O$、Taylor 余项、不等式决策树 & 【熟】 \\
\bottomrule
\end{tabular}
\end{center}

\section{Q. 拓扑学：从点集到同调}
\label{sec:topology}

\textbf{为什么单列一章：} A2 里我们已经在用「紧性」「完备性」这些词，但那是度量的版本；
弱收敛、弱紧性、算子谱这些概念\textbf{在度量框架里根本表述不出来}，必须先把「邻近」这件事抽象出来。
另一头，H1／M6 的网格拓扑（亏格、切割、Euler 特征）也需要同一套语言。

\textbf{与 \emph{li} 专题的分工：} \textbf{lie2} 里的「拓扑结构」一节（基本群、高阶同伦群、纤维丛、de Rham 上同调）
是\textbf{按李群路线切片}的版本；本章是\textbf{系统版}，二者互为参照，不必重复读。

%--------------------------------------------------------------
\subsection{Q1. 为什么需要拓扑}

\textbf{度量的不便：} 度量空间里「$x$ 离 $y$ 多近」由具体数字给出，可是：
\begin{itemize}
  \item $L^p$ 与 $C^k$ 是\textbf{不同}的度量，但它们给出同一批「开集」——该保留的是后者，不是前者；
  \item 弱收敛 $\langle f_n,\varphi\rangle\to\langle f,\varphi\rangle$ 对一切测试函数成立——
        它\textbf{不是任何范数}诱导的，却是一个完全合理的「收敛」概念；
  \item CG 里的网格粘合、切割、亏格，是纯粹的「邻接关系」问题，与边长无关。
\end{itemize}

所以\textbf{保留「哪些集合算开的」，丢掉具体距离}——这就是拓扑。

\begin{center}
\begin{tabular}{L{2.6cm}L{5.6cm}L{6.2cm}}
\toprule
\textbf{领域} & \textbf{用到的拓扑概念} & \textbf{用在哪} \\
\midrule
分析 & 紧性、可数性、分离公理 & Arzelà--Ascoli、弱*紧、$L^p$ 弱紧 \\
PDE & 紧嵌入、弱拓扑 & Rellich--Kondrachov、变分法的直接方法（见 A4） \\
几何 & 同胚、同伦、上同调 & 流形的定义（B2）、de Rham（A3） \\
CG & 同调、Euler 特征、亏格 & 网格拓扑修复、切割、参数化的拓扑障碍 \\
\bottomrule
\end{tabular}
\end{center}

%--------------------------------------------------------------
\subsection{Q2. 拓扑空间与基}

\textbf{拓扑空间：} 集合 $X$ 加一族子集 $\mathcal T$（称为\textbf{开集}），满足：
\begin{itemize}
  \item $\varnothing\in\mathcal T$，$X\in\mathcal T$；
  \item \textbf{任意}并封闭：$\{U_\alpha\}\subset\mathcal T\Rightarrow\bigcup_\alpha U_\alpha\in\mathcal T$；
  \item \textbf{有限}交封闭：$U_1,\dots,U_n\in\mathcal T\Rightarrow\bigcap_{i=1}^nU_i\in\mathcal T$。
\end{itemize}

\textbf{注意「任意」与「有限」的不对称}——与 A1 的 $\sigma$-代数（可数并）对照记忆：
$\sigma$-代数是「可数并 $+$ 补」，拓扑是「任意并 $+$ 有限交 $+$ 无补」。

\textbf{闭集：} 开集的补。有限并封闭、任意交封闭——注意开闭两族的封闭性方向正好相反。

\begin{center}
\begin{tabular}{lll}
\toprule
\textbf{拓扑} & \textbf{开集是什么} & \textbf{用途} \\
\midrule
度量拓扑 & 开球之并 & 标准分析 \\
离散拓扑 & 所有子集 & 反例 \\
平庸（平凡）拓扑 & 只有 $\varnothing,X$ & 反例（不 Hausdorff） \\
余有限拓扑 & 补集有限的集合 $+\varnothing$ & 反例（紧但不 Hausdorff） \\
子空间拓扑 & $U\cap A$，$U$ 开于 $X$ & $S^2\subset\mathbb R^3$、网格面 \\
积拓扑 & 各分量开集之积生成 & $\mathbb R^n$、$\prod_i X_i$ \\
商拓扑 & 使投影连续的最细拓扑 & 粘合、环面、射影平面 \\
\bottomrule
\end{tabular}
\end{center}

\textbf{基与子基：} $\mathcal B\subset\mathcal T$ 是\textbf{基}，若每个开集都是 $\mathcal B$ 中元素之并。
$\mathbb R^n$ 的开球构成基；子基是「生成」基的中间步骤。
\textbf{细／粗：} 开集多者为细。细拓扑下连续更严格、紧更困难——\textbf{同一集合上的两种拓扑，性质可以完全不同}。

%--------------------------------------------------------------
\subsection{Q3. 连续性与同胚}

\textbf{核心定义（比 $\varepsilon$-$\delta$ 更本质）：} $f:X\to Y$ 连续 $\iff$
对每个开集 $V\subset Y$，原像 $f^{-1}(V)$ 是 $X$ 中的开集。

\textbf{关键事实：} 度量空间中「开集原像开」与「$\varepsilon$-$\delta$ 定义」\textbf{等价}。
所以这是同一概念的更一般版本——把「距离小于 $\varepsilon$」换成了「落在某个开集里」。

\textbf{同胚：} 连续双射且逆也连续，记 $X\cong Y$。\textbf{拓扑学的中心问题就是分类同胚类。}
（甜圈与咖啡杯同胚；$[0,1]$ 与 $[0,1]^2$ 不同胚——这需要不变量，见 Q8、Q9。）

\textbf{三条实用判据：}
\begin{itemize}
  \item 连续映射把\textbf{紧}集送成\textbf{紧}集、把\textbf{连通}集送成\textbf{连通}集
        （这是把分析定理「搬」到拓扑层的手段：极值定理、介值定理都是它的特例）；
  \item 连续双射从紧空间到 Hausdorff 空间 $\Longrightarrow$ 同胚（\textbf{紧 Hausdorff 是「好」空间}）；
  \item \textbf{嵌入}：$f$ 是同胚到其像，且像带子空间拓扑。
\end{itemize}

%--------------------------------------------------------------
\subsection{Q4. 分离公理与可数性}

\begin{center}
\begin{tabular}{L{1.5cm}L{8.2cm}L{5.0cm}}
\toprule
\textbf{公理} & \textbf{条件} & \textbf{意义} \\
\midrule
$T_1$ & 任两点可各自被开集分开 & 单点集是闭集 \\
$T_2$（Hausdorff） & 任两点有\textbf{不相交}的开邻域 & \textbf{极限唯一}、网收敛的起点 \\
$T_3$（正则） & $T_1$ $+$ 点与不交闭集可被开集分开 & Urysohn 引理的前提 \\
$T_4$（正规） & $T_1$ $+$ 两个不交闭集可被开集分开 & 连续函数的构造（单位分解） \\
\bottomrule
\end{tabular}
\end{center}

\textbf{为什么 Hausdorff 是底线：} 没有它，序列（更一般是网）可以收敛到多个点，
「$\lim$」这个记号就没有意义——而整个分析都是建立在取极限上的。

\textbf{可数性（分析里更常用）：}
\begin{itemize}
  \item \textbf{第一可数}：每点有可数局部基（度量空间都有）$\Rightarrow$ 序列足够刻画连续与闭包；
  \item \textbf{第二可数}：整个空间有可数基 $\Rightarrow$ 可分（$\mathbb R^n$ 都有）；
  \item \textbf{可分}：存在可数稠密子集。$\mathbb R$ 可分这一点在 A1 里用过——
        它正是「每个开集是可数个基开集之并」，也就是 Borel $\sigma$-代数只需可数层级就能生成的理由。
\end{itemize}

\textbf{两个结论（陈述，证明见点集拓扑教材）：}
\begin{itemize}
  \item \textbf{Urysohn 引理：} $X$ 正规，$A,B$ 不交闭集 $\Longrightarrow$ 存在连续 $f:X\to[0,1]$ 使 $f|_A=0$、$f|_B=1$；
  \item \textbf{Urysohn 度量化定理：} 第二可数 $+$ 正则 $\Longrightarrow$ 可度量化。
        （即：拓扑空间「看起来像度量空间」的充要条件之一。）
\end{itemize}

%--------------------------------------------------------------
\subsection{Q5. 紧性：分析里最有用的一个拓扑性质}

\textbf{定义：} $K$ 紧 $\iff$ 任一开覆盖都有\textbf{有限}子覆盖。

\textbf{等价刻画（度量空间）：} 每个序列都有收敛子列（列紧）$\iff$ 完备 $+$ 全有界。

\begin{center}
\begin{tabular}{L{4.0cm}L{10.6cm}}
\toprule
\textbf{结论} & \textbf{内容} \\
\midrule
Heine--Borel & $\mathbb R^n$ 中：紧 $\iff$ 有界且闭 \\
连续像 & $f$ 连续，$K$ 紧 $\Longrightarrow$ $f(K)$ 紧 $\Rightarrow$ 极值定理（$\max$ 取到） \\
Heine--Cantor & $f$ 在紧集上连续 $\Longrightarrow$ $f$ 在紧集上一致连续（$\delta$ 与点无关） \\
Tychonoff & 任意一族紧空间的积仍紧（等价于选择公理） \\
紧 $+$ Hausdorff & $\Longrightarrow$ 正规（故可构造连续函数、可单位分解） \\
闭子集 & 紧空间的闭子集紧；Hausdorff 空间中的紧子集闭 \\
\bottomrule
\end{tabular}
\end{center}

\textbf{紧性在 PDE 里为什么是命根子：} 变分法要「找一个使能量最小的 $u$」，
标准套路是：取极小化序列 $u_n$ $\to$ 提取收敛子列 $\to$ 极限就是极小元。
\textbf{「提取收敛子列」这一步百分之百靠紧性。} 相应地：
\begin{itemize}
  \item \textbf{Banach--Alaoglu：} 对偶空间的单位球在弱*拓扑下紧（\textbf{不需要自反性}）；
  \item \textbf{Kakutani：} 自反空间（$1<p<\infty$ 的 $L^p$、Hilbert 空间）的单位球弱紧；
  \item $L^1$ \textbf{不自反}，弱紧性得换条件——由 Dunford--Pettis 定理，等价条件正是\textbf{均匀可积}
        （见 C6.3：与 DCT 的「盖子」是同一件事）；
  \item \textbf{Arzelà--Ascoli：} $C(K)$ 中，等度连续 $+$ 一致有界 $\Longrightarrow$ 存在一致收敛子列；
  \item \textbf{Rellich--Kondrachov（预告，见 A4）：} 有界区域上 $W^{1,p}\hookrightarrow\hookrightarrow L^p$ 是\textbf{紧嵌入}——
        这是椭圆 PDE 存在性理论的地基。
\end{itemize}

%--------------------------------------------------------------
\subsection{Q6. 连通性}

\begin{itemize}
  \item \textbf{连通：} 不能写成两个不交非空开集之并（$\varnothing,X$ 是仅有的既开又闭者）；
  \item \textbf{道路连通：} 任意两点有连续道路相连。道路连通 $\Longrightarrow$ 连通，\textbf{反向不成立}；
  \item \textbf{连通分支：} 极大连通子集；道路连通分支同理（可以严格多于连通分支）。
\end{itemize}

\textbf{介值定理是连通性的定理：} 连续像连通 $+$ $\mathbb R$ 中的连通集是区间 $\Rightarrow$ 取遍中间值。

\textbf{两个标准反例}（值得记住，专治「显然」）：
\begin{itemize}
  \item \textbf{拓扑学家的正弦曲线} $\{(x,\sin\frac1x):x>0\}\cup(\{0\}\times[-1,1])$：连通但\textbf{不}道路连通；
  \item \textbf{有理数集} $\mathbb Q$：完全不连通（每点是其连通分支），但仍可数稠密。
\end{itemize}

%--------------------------------------------------------------
\subsection{Q7. 商空间：把东西粘起来}

\textbf{商拓扑：} $X$ 上给定等价关系 $\sim$，$X/{\sim}$ 上取「使投影 $\pi:X\to X/{\sim}$ 连续」的\textbf{最细}拓扑。
等价地：$U\subset X/{\sim}$ 开 $\iff$ $\pi^{-1}(U)$ 开于 $X$。

\textbf{泛性质（实际操作时最有用的形式）：} 在 $X/{\sim}$ 上定义连续映射，只需在 $X$ 上定义、
并验证它与 $\sim$ 相容（在等价类上取常值）。

\begin{center}
\begin{tabular}{lll}
\toprule
\textbf{商空间} & \textbf{从什么粘} & \textbf{Euler 特征} \\
\midrule
$S^1$（圆） & $[0,1]$ 粘两端 & $0$ \\
$S^2$（球面） & 圆盘粘边界于一点 & $2$ \\
$T^2$（环面） & 正方形对边粘合 & $0$ \\
Möbius 带 & 长方形一边扭转后粘 & $0$ \\
Klein 瓶 & 正方形一边扭转后粘 & $0$ \\
$\mathbb{RP}^2$（射影平面） & 球面粘对径点 & $1$ \\
\bottomrule
\end{tabular}
\end{center}

\textbf{在 CG 里这就是「网格的边界粘合」：} 三角网格 $=$ 单纯复形（Q9），
它的\textbf{切割图（cut graph）}就是把曲面剪成圆盘的那组边——
\textbf{球面可以一刀剪开铺平，环面必须剪两刀}。这个「几刀」就是亏格，就是拓扑障碍，
也是 H1 的网格参数化为什么在环面上会失败的根源。

%--------------------------------------------------------------
\subsection{Q8. 同伦与基本群}

\textbf{同伦：} $f,g:X\to Y$ 同伦，若存在连续 $H:X\times[0,1]\to Y$ 使 $H(\cdot,0)=f$、$H(\cdot,1)=g$。
直观：一根连续地变形成另一根。\textbf{同伦等价}比同胚弱（$\mathbb R^n$ 与一点同伦等价，却不同胚）。

\textbf{基本群：} 固定基点 $x_0$，闭道路的同伦类在「首尾相接」下成群，记 $\pi_1(X,x_0)$。
\textbf{函子性：} 连续映射 $f$ 诱导同态 $f_*:\pi_1(X)\to\pi_1(Y)$，且同伦等价 $\Longrightarrow$ 同构。
\textbf{同胚 $\Rightarrow$ 同伦等价 $\Rightarrow$ $\pi_1$ 同构}——所以 $\pi_1$ 是分类同胚类的\textbf{不变量}。

\begin{center}
\begin{tabular}{ll}
\toprule
\textbf{空间} & $\pi_1$ \\
\midrule
$S^1$ & $\mathbb Z$（绕数） \\
$S^n$，$n\ge2$ & $0$（平凡） \\
$T^2$ & $\mathbb Z^2$ \\
$\mathbb{RP}^2$ & $\mathbb Z/2$ \\
$\mathbb R^n$、任意凸集 & $0$（可缩） \\
\bottomrule
\end{tabular}
\end{center}

\textbf{$\pi_1(S^1)=\mathbb Z$ 的用法（几个「一眼看出」的定理）：}
\begin{itemize}
  \item \textbf{Brouwer 不动点定理（二维）：} $f:D^2\to D^2$ 连续必有不动点；
  \item \textbf{毛球定理：} 偶数维球面 $S^{2k}$ 上的连续切向量场必有零点——
        所以\textbf{球面上不可能处处有非零风向}，而环面上可以；
  \item 代数基本定理（复系数多项式必有根）也有一个纯拓扑证明。
\end{itemize}

\textbf{计算工具（陈述）：} \textbf{Seifert--van Kampen 定理}——把 $X=U\cup V$ 的 $\pi_1$
用 $\pi_1(U)$、$\pi_1(V)$、$\pi_1(U\cap V)$ 的自由积拼出来（环面的 $\mathbb Z^2$ 就是这么算的）。
\textbf{更多同伦群与纤维丛}见 lie2 的「拓扑结构」一节，本章不重复。

%--------------------------------------------------------------
\subsection{Q9. 单纯同调与 Euler 特征}

同伦群（Q8）计算极难（$\pi_n(S^m)$ 至今未全知），\textbf{同调群}则是可算的——
它是「同伦群的线性化」。

\subsubsection{链条与边界算子}

把空间剖分为\textbf{单纯复形}（顶点、边、三角形、四面体……）。记 $C_k$ 为 $k$ 维单形的形式线性组合
（系数取 $\mathbb Z$ 或 $\mathbb Z/2$）。\textbf{边界算子} $\partial_k:C_k\to C_{k-1}$ 把每个 $k$-单形映为其 $k+1$ 个
$(k-1)$-面之和（带符号）。核心恒等式：
\[
\partial_{k-1}\circ\partial_k=0 \qquad(\text{“边界的边界为零”})
\]
它的连续对应物就是 A3 的 $d^2=0$。

\subsubsection{同调群与 Betti 数}

\[
H_k=\frac{\ker\partial_k}{\operatorname{im}\partial_{k+1}}
\qquad(\text{闭链 模去 边界})
\]
$b_k=\operatorname{rank}H_k$ 称为 \textbf{Betti 数}：$b_0=$ 连通分支数，$b_1=$ 「独立环路数」，$b_2=$ 「封闭空腔数」。

\begin{center}
\begin{tabular}{l l l l l}
\toprule
\textbf{空间} & $b_0$ & $b_1$ & $b_2$ & $\chi=\sum_k(-1)^kb_k$ \\
\midrule
$S^2$ & $1$ & $0$ & $1$ & $2$ \\
$T^2$ & $1$ & $2$ & $1$ & $0$ \\
$\mathbb{RP}^2$ & $1$ & $0$（$\mathbb Z/2$）& $0$ & $1$ \\
\bottomrule
\end{tabular}
\end{center}

\subsubsection{Euler 特征}

\[
\chi=V-E+F\qquad(\text{三角剖分：顶点}-\text{边}+\text{面})
\]
\textbf{它是拓扑不变量，与剖分无关}——这正是它有用的原因：换一个更细的网格，$\chi$ 不变。
亏格 $g$ 的闭可定向曲面：$\chi=2-2g$（球面 $g=0$，环面 $g=1$）。

\textbf{离散 Gauss--Bonnet（与 H1 直接对接）：}
\[
\sum_{i}\Big(2\pi-\sum_{j\in N(i)}\theta_{ij}\Big)=2\pi\chi
\]
左边是 H1 的\textbf{角度亏损}（离散高斯曲率）之和，右边只依赖拓扑。
\textbf{一句话：曲率的积分是拓扑的}——这是「离散微分几何」最漂亮的一条桥。

\textbf{在 CG 中的用途：}
\begin{itemize}
  \item 网格亏格 $=$ 手柄数，决定参数化策略与切割图（Q7）；
  \item $\chi\ne2-2g$ 说明网格有孔洞、非流形边或自交——\textbf{网格修复的判据}；
  \item Betti 数用于形状分析、拓扑简化、特征匹配。
\end{itemize}

%--------------------------------------------------------------
\subsection{Q10. de Rham 上同调：分析与拓扑的合流}

同调把空间变成群；\textbf{上同调}则从\textbf{函数（微分形式）}出发反过来读拓扑。

\textbf{构造（对照 A3 与 M6）：} 微分形式 $+$ 外微分 $d$ 构成复形
\[
\Omega^0\xrightarrow{\ d\ }\Omega^1\xrightarrow{\ d\ }\Omega^2\xrightarrow{\ d\ }\cdots,
\qquad d^2=0
\]
于是可以定义
\[
H^k_{dR}(M)=\frac{\ker d_k}{\operatorname{im}d_{k-1}}
\qquad(\text{闭形式 模去 恰当形式})
\]

\textbf{de Rham 定理（陈述）：} 光滑流形上 $H^k_{dR}(M)\cong H^k(M;\mathbb R)$（与单纯上同调同构）。
\textbf{这条定理把「解微分方程（找 $d\omega=\alpha$）」与「数洞」变成同一件事。}

\begin{center}
\begin{tabular}{lll}
\toprule
\textbf{问题} & \textbf{差分几何语言} & \textbf{拓扑语言} \\
\midrule
$\mathbf F$ 有势 $\iff$ $\nabla\times\mathbf F=0$（单连通域） & 闭 1-形式是恰当 1-形式 & $H^1=0$ \\
$\oint_C\mathbf F\cdot d\mathbf l$ 绕洞不为零 & 闭而非恰当 & $H^1\ne0$ \\
磁单极不存在 & $dF=0$ 且 $F$ 恰当 & $H^2(\mathbb R^3)=0$ \\
\bottomrule
\end{tabular}
\end{center}

\textbf{离散版本就是 M6 的 DEC：} 把 $\Omega^k$ 换成「$k$ 维单形上的值」，$d$ 换成关联矩阵，
$d^2=0$ 在离散层面\textbf{精确}成立（不是近似）。\textbf{Hodge 理论}（每个上同调类有唯一调和代表元）
则是「$d^2=0$ 拉普拉斯方程 $\leftrightarrow$ 同调」的进一步合流，与 A2 的谱理论、H1 的离散 Laplace 是同一主线。

%--------------------------------------------------------------
\subsection{Q11. 小结：拓扑概念在两条主线上的落点}

\begin{center}
\begin{tabular}{L{3.0cm}L{5.4cm}L{5.6cm}}
\toprule
\textbf{概念} & \textbf{PDE／分析线上的落点} & \textbf{CG／几何线上的落点} \\
\midrule
紧性（Q5） & Arzelà--Ascoli、弱*紧、直接方法、Rellich--Kondrachov（A4） &
参数域是紧流形；网格有限性 \\
\addlinespace
分离公理（Q4） & 极限唯一、弱拓扑的 Hausdorff 性 & —— \\
\addlinespace
可数性（Q4） & 可分 $\Rightarrow$ Borel $\sigma$-代数可数生成（A1） & 网格可数性、有限元剖分 \\
\addlinespace
商空间（Q7） & 商测度、周期边界条件 & 网格粘合、切割图、参数化障碍 \\
\addlinespace
$\pi_1$（Q8） & 定义域拓扑 $\Rightarrow$ 解的量子化／障碍 & 环路、形变、形状匹配 \\
\addlinespace
同调／Euler 特征（Q9） & 指标定理、临界点计数 & 亏格、网格修复、拓扑简化 \\
\addlinespace
上同调（Q10） & 解的存在性障碍（闭 vs 恰当） & DEC／Hodge 分解（M6） \\
\bottomrule
\end{tabular}
\end{center}

\textbf{与既有章节的接口：}
\begin{itemize}
  \item 微分形式与 $d^2=0$ 见 A3；离散版本见 M6；
  \item 流形与切空间见 B2；李群本人的拓扑（$\pi_1$、覆盖空间）见 lie1 与 lie2；
  \item 紧性与弱收敛的\textbf{分析侧用法}见 A2 与 A4；
  \item 离散 Gauss--Bonnet 与角度亏损见 H1；$\chi$ 与亏格在渲染管线中的用法见 lie3 的「拓扑在 CG 中的应用」一节。
\end{itemize}

\section{R. 复分析：解析函数、留数与共形映射}
\label{sec:complex}

\textbf{为什么单列一章：} 复分析是三条线的交汇处，而这三条线在本笔记里都已经出现却缺了总枢纽：
\begin{itemize}
  \item \textbf{PDE 线：} 调和函数与位势理论是 Laplace 方程的核心（A4 的弱解理论之外的另一半）；
  \item \textbf{变换线：} de3 里的 Laplace 变换讲了「极点 $\leftrightarrow$ $e^{s_0t}$」，但没讲\textbf{怎么反演}——
        那个公式就是留数定理；
  \item \textbf{CG 线：} 共形参数化、调和映射的数学基础全在这里。
\end{itemize}
另外 K4 的 Wick 转动 $t\to-i\tau$ 本质上是解析延拓，L7 重整化的收敛性也依赖复平面上的解析性。

\textbf{前置：} 复数与 Euler 公式见 P5（本章直接用）。

%--------------------------------------------------------------
\subsection{R1. 复可微比实可微强得多}

\textbf{定义（全纯／解析）：} $f:\Omega\to\mathbb C$ 在 $z_0$ 全纯，若
\[
f'(z_0)=\lim_{h\to0}\frac{f(z_0+h)-f(z_0)}{h}
\]
存在，其中 $h$ \textbf{沿复平面的任意方向}趋于 $0$。

\textbf{为什么这比实可微强：} 实可微只要求「从左右两个方向趋近得到同一斜率」（一维），
复可微要求「从无穷多个方向（一整圈）趋近都得到同一个复数」。这个约束强到\textbf{会产生奇迹}（R2）。

写成 $f=u+iv$（$u,v$ 实值）后，沿实轴与虚轴趋近相等，得
\[
\boxed{\ \frac{\partial u}{\partial x}=\frac{\partial v}{\partial y},\qquad
\frac{\partial u}{\partial y}=-\frac{\partial v}{\partial x}\ }
\qquad(\text{Cauchy--Riemann 方程})
\]

\textbf{两个立刻可用的推论：}
\begin{itemize}
  \item 分别对两式求导相加，得 $\Delta u=0$、$\Delta v=0$——即 \textbf{$u,v$ 都是调和函数}。
        全纯函数的实部与虚部自动是 Laplace 方程的解，这解释了 R7 里二维 PDE 为什么特殊。
  \item \textbf{反例：} $f(z)=\bar z$（即 $u=x,v=-y$）在\textbf{任何}点都不满足 C-R；
        $f(z)=|z|^2$（$u=x^2+y^2,v=0$）只在原点满足。
        可见「看起来光滑」根本不够——C-R 是很严的条件。
\end{itemize}

\textbf{调和共轭：} 给定调和 $u$，总能（在单连通域上）找到 $v$ 使 $u+iv$ 全纯，$v$ 称 $u$ 的调和共轭。
\textbf{这就是 A4 里「$H^1=0$ 时闭形式恰当」的复分析面孔}（对照 Q10 的表）。

%--------------------------------------------------------------
\subsection{R2. Cauchy 积分定理与公式}

\subsubsection{围道积分}

$\oint_\gamma f(z)\,dz$ 按参数化 $z(t)$ 定义为 $\int_a^b f(z(t))z'(t)\,dt$，与实数积分同构。

\subsubsection{Cauchy 积分定理}

\textbf{定理：} $f$ 在\textbf{单连通}域 $\Omega$ 上全纯，$\gamma\subset\Omega$ 是闭合可求长曲线，则
\[
\oint_\gamma f(z)\,dz=0
\]

\textbf{这是全章的地基，也是它与拓扑的接点。} 用 A3／Q10 的语言：
$f(z)dz$ 是闭 $1$-形式（C-R 方程正是「闭」的条件），
而「单连通」保证 $H^1(\Omega)=0$，于是闭形式自动恰当，积分必然为零。
\textbf{反例记住这一个：} $f(z)=1/z$ 在 $\mathbb C\setminus\{0\}$ 上全纯（该域\textbf{不}单连通），而
\[
\oint_{|z|=1}\frac{dz}{z}=2\pi i\ne0
\]
这个 $2\pi i$ 就是「洞」的度量——它将在 R4 变成留数的定义常数。

\subsubsection{Cauchy 积分公式}

\[
f(a)=\frac{1}{2\pi i}\oint_\gamma\frac{f(z)}{z-a}\,dz
\qquad(a\ \text{在}\ \gamma\ \text{内部})
\]
\textbf{读法：全纯函数在区域内部的值，完全由它在边界上的值决定。}
这与实数情形完全不同——实函数的内部值不受边界约束，而全纯函数是「刚性的」。

反复求导（对 $a$ 微分即可，因为被积函数对 $a$ 光滑）得
\[
f^{(n)}(a)=\frac{n!}{2\pi i}\oint_\gamma\frac{f(z)}{(z-a)^{n+1}}\,dz
\]

\textbf{最重要的推论：} 全纯 $\Longrightarrow$ 无穷次复可微 $\Longrightarrow$ 局部可展开为幂级数（解析）。
\textbf{「可微一次」等于「可微无穷次、且等于自己的 Taylor 级数」}——
这是实分析里绝对没有的性质，也是全章所有「奇迹」的共同来源。

\subsubsection{三个经典推论}

\begin{center}
\begin{tabular}{L{3.0cm}L{5.4cm}L{5.6cm}}
\toprule
\textbf{定理} & \textbf{内容} & \textbf{用法} \\
\midrule
Liouville & 有界整函数（全平面全纯）必为常数 &
由 $|f'|\le M/R\to0$ 得 $f'\equiv0$ \\
代数基本定理 & 任意 $n\ge1$ 次复系数多项式必有根 &
若 $p$ 无处为根，则 $1/p$ 有界整函数，矛盾 \\
最大模原理 & $|f|$ 在区域内部取不到极大（除非 $f$ 为常数） &
与调和函数的平均值性质同源（R7） \\
\bottomrule
\end{tabular}
\end{center}

%--------------------------------------------------------------
\subsection{R3. 幂级数、Laurent 级数与奇点}

\textbf{Taylor 级数：} $f(z)=\sum_{n\ge0}a_n(z-a)^n$，收敛半径 $R=$ \textbf{圆心到最近奇点的距离}
（这条「到奇点的距离决定收敛半径」是判收敛最快的手段）。

\textbf{Laurent 级数：} 在圆环 $r<|z-a|<R$ 上
\[
f(z)=\sum_{n=-\infty}^{\infty}c_n(z-a)^n
\qquad(\text{允许负幂项})
\]
于是可以按负幂项对孤立奇点分类：

\begin{center}
\begin{tabular}{L{2.6cm}L{5.2cm}L{6.2cm}}
\toprule
\textbf{类型} & \textbf{负幂项} & \textbf{判据／例子} \\
\midrule
可去奇点 & 全无（$c_n=0,\ n<0$） & $\lim_{z\to a}f(z)$ 存在且有限；$\frac{\sin z}{z}$ 在 $0$ \\
$m$ 阶极点 & 有限，最低次为 $-m$ & $\lim_{z\to a}(z-a)^mf(z)$ 是非零有限数；$\frac{1}{z^2}$ 在 $0$ 是二阶 \\
本质奇点 & 无穷多项 & $\lim$ 不存在（值域在任意邻域内稠密）；$e^{1/z}$ 在 $0$ \\
\bottomrule
\end{tabular}
\end{center}

%--------------------------------------------------------------
\subsection{R4. 留数定理}

\subsubsection{定义与定理}

\textbf{留数：} $f$ 在孤立奇点 $a$ 处的 Laurent 展开中 $c_{-1}$ 那一项，记 $\operatorname{Res}(f,a)$——\textbf{只有这一项在积分后存活}（因为 $\oint(z-a)^n dz=0$ 当 $n\ne-1$，当 $n=-1$ 时等于 $2\pi i$）。

\textbf{留数定理：} $f$ 在 $\gamma$ 内部除有限个孤立奇点 $a_1,\dots,a_m$ 外全纯，则
\[
\boxed{\ \oint_\gamma f(z)\,dz=2\pi i\sum_{k=1}^m\operatorname{Res}(f,a_k)\ }
\]

\subsubsection{怎么算留数}

\begin{center}
\begin{tabular}{L{3.4cm}L{9.4cm}}
\toprule
\textbf{情形} & \textbf{公式} \\
\midrule
单极点 & $\operatorname{Res}(f,a)=\lim_{z\to a}(z-a)f(z)$ \\
$f=p/q$，$q$ 在 $a$ 有单零点 & $\operatorname{Res}(f,a)=\dfrac{p(a)}{q'(a)}$（最常用，免求极限） \\
$m$ 阶极点 & $\operatorname{Res}(f,a)=\dfrac{1}{(m-1)!}\lim_{z\to a}\dfrac{d^{m-1}}{dz^{m-1}}\big[(z-a)^mf(z)\big]$ \\
\bottomrule
\end{tabular}
\end{center}

\subsubsection{三类实积分的标准套路}

\begin{center}
\begin{tabular}{L{4.6cm}L{4.6cm}L{4.8cm}}
\toprule
\textbf{积分类型} & \textbf{替换／围道} & \textbf{要点} \\
\midrule
$\displaystyle\int_0^{2\pi}R(\cos\theta,\sin\theta)\,d\theta$ &
$z=e^{i\theta}$，$d\theta=\dfrac{dz}{iz}$，$\cos\theta=\dfrac{z+z^{-1}}{2}$ &
化为单位圆上的留数计算 \\
\addlinespace
$\displaystyle\int_{-\infty}^{\infty}R(x)\,dx$ &
上半平面闭合，加半圆 & 需 $\deg q\ge\deg p+2$，半圆上积分 $\to0$ \\
\addlinespace
$\displaystyle\int_{-\infty}^{\infty}R(x)e^{iax}\,dx$ &
上半平面闭合 & $a>0$ 时由 \textbf{Jordan 引理} 保证半圆贡献 $\to0$ \\
\bottomrule
\end{tabular}
\end{center}

\textbf{另外两个计数工具：}
\begin{itemize}
  \item \textbf{辐角原理：} $\dfrac{1}{2\pi i}\oint_\gamma\dfrac{f'}{f}\,dz=$ $\gamma$ 内部零点数 $-$ 极点数（计重数）；
  \item \textbf{Rouché 定理：} $|g|<|f|$ 于 $\partial D$ 上 $\Longrightarrow$ $f$ 与 $f+g$ 在 $D$ 内零点数相同。
\end{itemize}

\textbf{与 de3 的接口（补上「反演」这一环）：} Laplace 反演公式
\[
f(t)=\frac{1}{2\pi i}\int_{c-i\infty}^{c+i\infty}e^{st}F(s)\,ds
\]
是一个沿竖直直线的「积分」，把它向右闭合（取 $t>0$）后用留数定理，即得
\[
f(t)=\sum_k \operatorname{Res}\big(e^{st}F(s),\ s_k\big)
\]
\textbf{而每个极点 $s_k$ 贡献一个 $e^{s_k t}$——这正是 de3 里那张「极点 $\leftrightarrow$ 指数解」对照表的来源。}
重极点贡献 $t^ke^{s_kt}$（共振），也是同一公式的自动推论。

%--------------------------------------------------------------
\subsection{R5. 共形映射}

\textbf{保角性：} $f'(z_0)\ne0$ 时，$f$ 在 $z_0$ 处\textbf{保角}——
把过 $z_0$ 的两条曲线的夹角按大小与方向原样保持。
（理由：局部 $f(z)\approx f(z_0)+f'(z_0)(z-z_0)$，是「旋转 $+$ 缩放」，而这两者都保角。）

\textbf{与 Laplace 方程的相容性：} 若 $u$ 调和、$f$ 全纯，则 $u\circ f$ 也调和
（因为调和性 $\Delta u=0$ 在保角变换下不变）。
\textbf{一句话：共形映射是 Laplace 方程的对称群}——这就是二维 PDE 与流体力学用复分析的根据。

\textbf{Riemann 映射定理：} 任意单连通真子域 $\Omega\subsetneq\mathbb C$ 共形等价于单位圆盘
（即存在共形双射 $f:\Omega\to\mathbb D$）。这是二维特有的奇迹；\textbf{三维及以上没有对应结论}。

\textbf{Möbius 变换：} $z\mapsto\dfrac{az+b}{cz+d}$（$ad-bc\ne0$）把广义圆映成广义圆；
单位圆盘的自同构为 $z\mapsto e^{i\theta}\dfrac{z-a}{1-\bar a z}$（$|a|<1$）。

\textbf{在 CG 中的落地——共形参数化：}
\begin{itemize}
  \item 共形参数化把曲面摊平到平面域时\textbf{消灭角度畸变}（因为保角），代价是面积畸变无法同时消除
        （由 R5 的「保角 $+$ 保面积 $\Rightarrow$ 等距」可知：两者不可兼得）；
  \item 常用算法（圆模式圆填充、离散 Ricci 流）本质是在离散层面解一个 Ricci 流方程——
        与 H1 的余切 Laplace 算子共用同一套离散结构；
  \item 与 B2 的黎曼度量、H1 的离散 Laplace 算子接上：\textbf{「共形」在离散层面就是「保持余切权重比」}。
\end{itemize}

%--------------------------------------------------------------
\subsection{R6. 解析延拓与多值函数}

\textbf{唯一性定理：} 两个在连通域上全纯的函数，若在一个有聚点的集合上相等，则处处相等。
（这是全纯函数「刚性」的又一次体现：局部信息决定整体。）

\textbf{解析延拓：} 沿路径把局部幂级数一块块接下去。若路径绕奇点一圈后回来得到\textbf{不同}的值，
函数就是\textbf{多值}的——$z=0$ 称为\textbf{支点}。

\begin{itemize}
  \item $\log z$、$\sqrt z$、$z^\alpha$（$\alpha\notin\mathbb Z$）都多值；
  \item 处理办法：\textbf{割线}（单值分支）或 \textbf{Riemann 面}（把多值函数变成单值函数的新空间）；
  \item 留数定理对多值函数照样能用，只是要小心割线两侧的跳跃。
\end{itemize}

\textbf{$\Gamma$ 函数的解析延拓（接 P8）：} P8 里 $\Gamma(s)$ 只在 $\operatorname{Re}s>0$ 有定义。
用递推关系反解
\[
\Gamma(s)=\frac{\Gamma(s+1)}{s}
\]
可把定义域逐段向左推进，得到除 $s=0,-1,-2,\dots$ 外处处全纯的函数，
这些点是\textbf{单极点}，且 $\operatorname{Res}\big(\Gamma,-n\big)=\dfrac{(-1)^n}{n!}$
（又是一个留数计算）。\textbf{所以「$\Gamma$ 是阶乘的连续插值」这句话，严格含义是「它的解析延拓」。}

\textbf{与 K4 的接口：} Wick 转动 $t\to-i\tau$ 就是把实时间上的传播子解析延拓到虚时间轴——
之所以合法，是因为传播子作为函数在适当区域上全纯（物理上称「解析性」，
它的谱表示就是 R2 积分公式的物理版）。

%--------------------------------------------------------------
\subsection{R7. 与 PDE 的接口：二维为什么特殊}

\textbf{调和函数的基本性质（与 R2 同源）：}

\begin{center}
\begin{tabular}{L{4.0cm}L{10.6cm}}
\toprule
\textbf{性质} & \textbf{内容} \\
\midrule
平均值性质 & $u(a)=\dfrac{1}{2\pi}\displaystyle\int_0^{2\pi}u(a+re^{i\theta})\,d\theta$（圆心值 $=$ 圆周平均） \\
最大（小）值原理 & 非常数调和函数在区域内部取不到极值 \\
基本解 & $n=2$：$\dfrac{1}{2\pi}\log|z|$；$n=3$：$-\dfrac{1}{4\pi|\mathbf x|}$（A2.5 的 $\delta$ 的解） \\
Liouville（调和版） & 全平面上有上界的调和函数是常数 \\
\bottomrule
\end{tabular}
\end{center}

\textbf{二维不可压无粘流 $=$ 全纯函数。} 记\textbf{复势}
\[
w(z)=\varphi(x,y)+i\psi(x,y)
\]
其中 $\varphi$ 是速度势、$\psi$ 是流函数。由 C-R 方程，
\[
u-iv=w'(z)\qquad(u,v\ \text{为速度分量})
\]
且流线就是 $\operatorname{Im}w=\text{const}$。
\textbf{于是「解 2D 无粘流」变成「找一个全纯函数」}——而 R5 的共形映射允许先把复杂形状
共形搬到简单形状（圆）上求解，再搬回来。经典的 Joukowski 变换求机翼绕流就是这么做的。

\textbf{但这就是复分析方法的边界，必须说清：} 上述结论依赖\textbf{无粘、无旋}。
一旦要有粘性、要有涡量动力学与湍流，复势方法失效，必须回到 I 章的 Navier--Stokes
（那里二维的涡量方程 $\frac{D\omega}{Dt}=\nu\Delta\omega$ 是标量输运方程，
比三维干净，但已经不是全纯函数问题）。

%--------------------------------------------------------------
\subsection{R8. 小结：复分析概念在两线上的落点}

\begin{center}
\begin{tabular}{L{3.4cm}L{5.4cm}L{5.6cm}}
\toprule
\textbf{概念} & \textbf{PDE／分析线} & \textbf{CG／几何线} \\
\midrule
C-R 方程与调和性（R1） & Laplace 方程的解、调和共轭 & 调和映射、网格上的调和坐标 \\
Cauchy 定理（R2） & 单连通 $\Rightarrow$ 闭形式恰当（Q10） & 环路积分的路径无关性 \\
Cauchy 公式与刚性（R2） & 边界值决定内部（最大模原理） & —— \\
留数定理（R4） & Laplace／Fourier 反演（de3）、$\Gamma$ 的极点 & 传递函数的极点分析 \\
Laplace 级数与奇点（R3） & 极点阶数 $=$ 解的增长型（$t^ke^{s_0t}$） & 滤波器极点、共振 \\
共形映射（R5） & 调和性在保角下不变 & \textbf{共形参数化}、纹理畸变控制 \\
解析延拓（R6） & Wick 转动（K4）、重整化（L7） & —— \\
复势（R7） & 2D 无粘流 $=$ 全纯函数 & 平面流场、矢量场设计 \\
\bottomrule
\end{tabular}
\end{center}

\textbf{与既有章节的接口：} 复数与 Euler 公式见 P5；Laplace 变换的极点表见 de3；
Wick 转动见 K4；离散 Laplace 与余切权重见 H1；调和映射与网格参数化见 H1 与 lie3。


\clearpage
\part{专题一：微分方程——从考研套路到现代理论}
%==============================================================
\section{第一层：考研/高数级别的 ODE 解题套路}
%==============================================================

\subsection{核心思想}

所有方法本质上都在做两件事之一：
\begin{enumerate}
  \item \textbf{变量分离}：将 $x$ 和 $y$ 分到等式两边
  \item \textbf{凑全微分}：将方程写成 $dF(x,y) = 0$ 的形式
\end{enumerate}

\subsection{一阶方程分类与解法}

\begin{center}
\begin{tabular}{L{2.5cm}L{4cm}L{4cm}L{3cm}}
\toprule
\textbf{类型} & \textbf{标准形式} & \textbf{解法} & \textbf{本质} \\
\midrule
可分离变量 & $y' = f(x)g(y)$ & $\frac{dy}{g(y)} = f(x)\,dx$ & 直接分离 \\
齐次方程 & $y' = \varphi(y/x)$ & 令 $u = y/x$ & 尺度不变性 \\
一阶线性 & $y' + P(x)y = Q(x)$ & 积分因子 $\mu = e^{\int P\,dx}$ & 凑全微分 \\
Bernoulli & $y' + Py = Qy^n$ & 令 $z = y^{1-n}$ & 非线性$\to$线性 \\
全微分 & $M\,dx + N\,dy = 0$ & 找 $F$ 使 $dF = M\,dx + N\,dy$ & 恰当形式 \\
\bottomrule
\end{tabular}
\end{center}

\subsubsection{积分因子的系统求法}

对 $M\,dx + N\,dy = 0$（不满足 $\frac{\partial M}{\partial y} = \frac{\partial N}{\partial x}$），找 $\mu$ 使得：
\[
\frac{\partial(\mu M)}{\partial y} = \frac{\partial(\mu N)}{\partial x}
\]

展开得：
\[
M\frac{\partial \mu}{\partial y} - N\frac{\partial \mu}{\partial x} = \mu\left(\frac{\partial N}{\partial x} - \frac{\partial M}{\partial y}\right)
\]

\textbf{特殊情况：}
\begin{itemize}
  \item 若 $\frac{1}{N}\left(\frac{\partial M}{\partial y} - \frac{\partial N}{\partial x}\right) = f(x)$（只依赖 $x$），则 $\mu = e^{\int f(x)\,dx}$
  \item 若 $\frac{1}{M}\left(\frac{\partial N}{\partial x} - \frac{\partial M}{\partial y}\right) = g(y)$（只依赖 $y$），则 $\mu = e^{\int g(y)\,dy}$
\end{itemize}

\textbf{示例：} 解 $(3x^2 + 6xy^2)\,dx + (6x^2y + 4y^2)\,dy = 0$

验证：$\frac{\partial M}{\partial y} = 12xy$，$\frac{\partial N}{\partial x} = 12xy$。满足恰当条件。

找 $F$：$\frac{\partial F}{\partial x} = 3x^2 + 6xy^2 \Rightarrow F = x^3 + 3x^2y^2 + \varphi(y)$

$\frac{\partial F}{\partial y} = 6x^2y + \varphi'(y) = 6x^2y + 4y^2 \Rightarrow \varphi(y) = \frac{4}{3}y^3$

解：$x^3 + 3x^2y^2 + \frac{4}{3}y^3 = C$

\subsubsection{一阶线性方程的完整示例}

\textbf{例题：} 解 $y' + \frac{2}{x}y = x^2$，$x > 0$

\textbf{方法一（积分因子）：}

$\mu = e^{\int 2/x\,dx} = e^{2\ln x} = x^2$

两边乘以 $x^2$：$x^2 y' + 2xy = x^4$，即 $(x^2 y)' = x^4$

积分：$x^2 y = \frac{x^5}{5} + C$

\[
\boxed{y = \frac{x^3}{5} + \frac{C}{x^2}}
\]

\textbf{方法二（常数变易法）：}

齐次解：$y_h = \frac{C}{x^2}$

设 $y = \frac{C(x)}{x^2}$，代入：$\frac{C'(x)}{x^2} = x^2$，$C(x) = \frac{x^5}{5} + C_0$

结果相同。

\subsection{可降阶的高阶方程}

\begin{center}
\begin{tabular}{L{3cm}L{4cm}L{5cm}}
\toprule
\textbf{类型} & \textbf{形式} & \textbf{降阶方法} \\
\midrule
缺 $y$ & $F(x, y', y'', \ldots) = 0$ & 令 $p = y'$，降一阶 \\
缺 $x$ & $F(y, y', y'', \ldots) = 0$ & 令 $p = y'$，$y'' = p\frac{dp}{dy}$ \\
\bottomrule
\end{tabular}
\end{center}

\textbf{示例（缺 $x$）：} 解 $yy'' + (y')^2 = 0$

令 $p = y'$，$y'' = p\frac{dp}{dy}$：
\[
yp\frac{dp}{dy} + p^2 = 0 \implies p\left(y\frac{dp}{dy} + p\right) = 0
\]

$p = 0$ 给出 $y = C$（平凡解）。

$y\frac{dp}{dy} + p = 0$：$\frac{dp}{p} = -\frac{dy}{y}$，$\ln|p| = -\ln|y| + C$，$p = \frac{C_1}{y}$

即 $y' = \frac{C_1}{y}$，$y\,dy = C_1\,dx$，$\frac{y^2}{2} = C_1 x + C_2$

\[
\boxed{y^2 = 2C_1 x + C_2}
\]

\textbf{对称性解读：} 方程缺 $x$ → 平移对称性 $x \to x + c$ → 可以降阶。

\subsection{二阶线性常系数方程}

\[
y'' + py' + qy = 0
\]

\textbf{特征方程：} $r^2 + pr + q = 0$

\begin{center}
\begin{tabular}{L{2.5cm}L{3cm}L{5cm}}
\toprule
\textbf{判别式} & \textbf{根} & \textbf{通解} \\
\midrule
$\Delta > 0$ & 不等实根 $r_1, r_2$ & $y = C_1 e^{r_1 x} + C_2 e^{r_2 x}$ \\
$\Delta = 0$ & 重根 $r$ & $y = (C_1 + C_2 x)e^{rx}$ \\
$\Delta < 0$ & 复根 $\alpha \pm i\beta$ & $y = e^{\alpha x}(C_1\cos\beta x + C_2\sin\beta x)$ \\
\bottomrule
\end{tabular}
\end{center}

\textbf{为什么重根时出现 $xe^{rx}$？}

\textbf{解释一（极限）：} 设 $r_1 = r + \epsilon$，$r_2 = r - \epsilon$，则：
\[
\frac{e^{(r+\epsilon)x} - e^{(r-\epsilon)x}}{2\epsilon} = e^{rx}\frac{e^{\epsilon x} - e^{-\epsilon x}}{2\epsilon} \xrightarrow{\epsilon \to 0} xe^{rx}
\]

\textbf{解释二（算子/若尔当）：} 方程等价于 $(D - r)^2 y = 0$，其中 $D = d/dx$。$(D-r)^2$ 的零空间由 $e^{rx}$ 和 $xe^{rx}$ 张成——这是若尔当块的体现。

\textbf{示例：} 解 $y'' - 4y' + 4y = 0$

特征方程：$r^2 - 4r + 4 = (r-2)^2 = 0$，$r = 2$（重根）

\[
\boxed{y = (C_1 + C_2 x)e^{2x}}
\]

\subsection{Euler 方程（等维方程）}

\[
x^2 y'' + axy' + by = 0
\]

\textbf{特征：} 系数是 $x$ 的幂次，具有尺度不变性（$x \to \lambda x$ 下方程形式不变）。

\textbf{解法一：} 令 $y = x^r$，代入得指标方程：
\[
r(r-1) + ar + b = 0
\]

\textbf{解法二：} 令 $x = e^t$（$t = \ln x$），则 $x\frac{d}{dx} = \frac{d}{dt}$，$x^2\frac{d^2}{dx^2} = \frac{d^2}{dt^2} - \frac{d}{dt}$，化为常系数方程。

\textbf{示例：} 解 $x^2 y'' - 2xy' + 2y = 0$

令 $y = x^r$：$r(r-1) - 2r + 2 = 0$，$r^2 - 3r + 2 = 0$，$r = 1, 2$

\[
\boxed{y = C_1 x + C_2 x^2}
\]

\subsection{非齐次方程：待定系数法与常数变易法}

\subsubsection{待定系数法（常系数）}

对 $y'' + py' + qy = f(x)$，根据 $f(x)$ 的形式猜测特解形式：

\begin{center}
\begin{tabular}{L{4cm}L{5cm}}
\toprule
\textbf{$f(x)$ 的形式} & \textbf{特解猜测} \\
\midrule
$P_n(x)$（$n$ 次多项式） & $Q_n(x)$（$n$ 次多项式） \\
$e^{\alpha x}$ & $Ae^{\alpha x}$（若 $\alpha$ 非特征根） \\
$e^{\alpha x}P_n(x)$ & $e^{\alpha x}Q_n(x)$ \\
$e^{\alpha x}\cos\beta x$ 或 $e^{\alpha x}\sin\beta x$ & $e^{\alpha x}(A\cos\beta x + B\sin\beta x)$ \\
\bottomrule
\end{tabular}
\end{center}

\textbf{共振情况：} 若 $\alpha$（或 $\alpha \pm i\beta$）是特征方程的 $k$ 重根，则特解需乘以 $x^k$。

\textbf{示例：} 解 $y'' + y = x\cos x$

齐次解：$y_h = C_1\cos x + C_2\sin x$

$f(x) = x\cos x$，$\alpha + i\beta = i$ 是特征根（$r^2 + 1 = 0$），$k = 1$。

猜测：$y^* = x[(Ax + B)\cos x + (Cx + D)\sin x]$

（代入比较系数，过程略）

\subsubsection{常数变易法（通用）}

设 $y = u_1(x)y_1 + u_2(x)y_2$，约束：
\[
u_1' y_1 + u_2' y_2 = 0
\]
\[
u_1' y_1' + u_2' y_2' = f(x)
\]

解出：
\[
u_1' = -\frac{y_2 f}{W}, \quad u_2' = \frac{y_1 f}{W}
\]

其中 $W = y_1 y_2' - y_1' y_2$ 是 Wronskian。

%==============================================================
\section{第二层：常微分方程的系统理论}
%==============================================================

\subsection{存在唯一性定理}

\textbf{Picard--Lindelöf 定理：} 考虑
\[
\frac{dy}{dx} = f(x, y), \quad y(x_0) = y_0
\]

若 $f$ 连续且关于 $y$ 满足 Lipschitz 条件 $|f(x,y_1) - f(x,y_2)| \leq L|y_1 - y_2|$，则初值问题在 $x_0$ 的某邻域内存在唯一解。

\textbf{证明思想（Picard 迭代）：}
\[
y_0(x) = y_0, \quad y_{n+1}(x) = y_0 + \int_{x_0}^x f(t, y_n(t))\,dt
\]

在适当条件下 $y_n \to y$（压缩映射原理）。

\subsection{线性系统}

\[
\mathbf{x}' = A\mathbf{x}
\]

\textbf{通解：} $\mathbf{x}(t) = e^{At}\mathbf{c}$

\textbf{矩阵指数：}
\[
e^{At} = \sum_{k=0}^{\infty} \frac{(At)^k}{k!}
\]

\textbf{计算方法：}
\begin{itemize}
  \item 若 $A$ 可对角化：$A = P\Lambda P^{-1}$，则 $e^{At} = Pe^{\Lambda t}P^{-1}$
  \item 若有若尔当块：$e^{Jt}$ 包含 $t^k e^{\lambda t}$ 项
  \item Cayley--Hamilton：$e^{At}$ 可表示为 $A$ 的 $n-1$ 次多项式
\end{itemize}

\textbf{示例：} 解 $\mathbf{x}' = \begin{pmatrix} 1 & 1 \\ 0 & 1 \end{pmatrix}\mathbf{x}$

$A$ 有若尔当块（特征值 $\lambda = 1$，代数重数 2，几何重数 1）：
\[
e^{At} = e^t \begin{pmatrix} 1 & t \\ 0 & 1 \end{pmatrix}
\]

通解：
\[
\mathbf{x}(t) = e^t \begin{pmatrix} C_1 + C_2 t \\ C_2 \end{pmatrix}
\]

\subsection{稳定性理论}

\textbf{平衡点分类（二维线性系统 $\mathbf{x}' = A\mathbf{x}$）：}

\begin{center}
\begin{tabular}{L{3.5cm}L{3cm}L{3cm}}
\toprule
\textbf{特征值} & \textbf{类型} & \textbf{稳定性} \\
\midrule
两个负实根 & 稳定结点 & 渐近稳定 \\
两个正实根 & 不稳定结点 & 不稳定 \\
一正一负 & 鞍点 & 不稳定 \\
纯虚根 $\pm i\beta$ & 中心 & 稳定（非渐近） \\
复根 $\alpha \pm i\beta$，$\alpha < 0$ & 稳定焦点 & 渐近稳定 \\
复根 $\alpha \pm i\beta$，$\alpha > 0$ & 不稳定焦点 & 不稳定 \\
\bottomrule
\end{tabular}
\end{center}

\textbf{Lyapunov 方法：} 找正定函数 $V(\mathbf{x})$（“能量”），若 $\dot{V} \leq 0$，则平衡点稳定。

\subsection{Sturm--Liouville 问题}

\[
-(p(x)y')' + q(x)y = \lambda w(x)y
\]

加上边界条件（如 $y(a) = y(b) = 0$）。

\textbf{性质：}
\begin{itemize}
  \item 特征值 $\lambda_1 < \lambda_2 < \cdots \to \infty$（离散、实数）
  \item 特征函数关于权重 $w$ 正交：$\int_a^b y_m y_n w\,dx = 0$（$m \neq n$）
  \item 特征函数构成 $L^2_w[a,b]$ 的完备正交基
\end{itemize}

\textbf{经典例子：}

\begin{center}
\begin{tabular}{L{4cm}L{3cm}L{4cm}}
\toprule
\textbf{方程} & \textbf{特征函数} & \textbf{对应} \\
\midrule
$-y'' = \lambda y$，$[0,\pi]$ & $\sin(nx)$ & Fourier 正弦级数 \\
$-y'' = \lambda y$，周期 & $e^{inx}$ & Fourier 级数 \\
Bessel & $J_\nu(\sqrt{\lambda}x)$ & Bessel 函数 \\
Legendre & $P_n(x)$ & Legendre 多项式 \\
$-y'' + x^2 y = \lambda y$ & $H_n(x)e^{-x^2/2}$ & Hermite 函数 \\
\bottomrule
\end{tabular}
\end{center}

%==============================================================
\section{第三层：特殊函数与级数方法}
%==============================================================

\subsection{幂级数解法}

\subsubsection{常点展开}

若 $P(x)$，$Q(x)$ 在 $x_0$ 解析，设 $y = \sum_{n=0}^{\infty} a_n (x-x_0)^n$，代入方程比较系数。

\textbf{示例：} 解 $y'' + y = 0$（在 $x_0 = 0$ 展开）

设 $y = \sum a_n x^n$，$y'' = \sum_{n=2}^{\infty} n(n-1)a_n x^{n-2} = \sum_{n=0}^{\infty}(n+2)(n+1)a_{n+2} x^n$

代入：$(n+2)(n+1)a_{n+2} + a_n = 0$，即 $a_{n+2} = -\frac{a_n}{(n+2)(n+1)}$

偶数项：$a_0, a_2 = -\frac{a_0}{2!}, a_4 = \frac{a_0}{4!}, \ldots$ → $a_0\cos x$

奇数项：$a_1, a_3 = -\frac{a_1}{3!}, a_5 = \frac{a_1}{5!}, \ldots$ → $a_1\sin x$

\[
y = a_0\cos x + a_1\sin x
\]

\subsubsection{Frobenius 方法（正则奇点）}

若 $x_0$ 是正则奇点，设：
\[
y = (x-x_0)^r \sum_{n=0}^{\infty} a_n (x-x_0)^n, \quad a_0 \neq 0
\]

代入得\textbf{指标方程}（确定 $r$）和递推关系。

\textbf{指标方程两根 $r_1, r_2$ 的关系：}
\begin{itemize}
  \item $r_1 - r_2 \notin \mathbb{Z}$：两个线性无关解都是 Frobenius 形式
  \item $r_1 - r_2 \in \mathbb{Z}$：第二个解可能含 $\ln(x-x_0)$ 项
  \item $r_1 = r_2$：第二个解一定含 $\ln$ 项
\end{itemize}

\textbf{示例：} 解 $xy'' + y' + xy = 0$（$x_0 = 0$ 是正则奇点）

设 $y = x^r \sum a_n x^n$。代入，最低次幂项给出指标方程：
\[
r^2 = 0 \implies r = 0 \text{（重根）}
\]

第一个解：$y_1 = J_0(x) = \sum_{k=0}^{\infty} \frac{(-1)^k}{(k!)^2}\left(\frac{x}{2}\right)^{2k}$（Bessel 函数 $J_0$）

第二个解含 $\ln x$：$y_2 = J_0(x)\ln x + \cdots$

\subsection{Bessel 方程}

\[
x^2 y'' + xy' + (x^2 - \nu^2)y = 0
\]

\textbf{来源：} 柱坐标下 Laplace 方程的径向部分。

\textbf{解：}
\[
J_\nu(x) = \sum_{k=0}^{\infty} \frac{(-1)^k}{k!\,\Gamma(k+\nu+1)}\left(\frac{x}{2}\right)^{2k+\nu}
\]

\textbf{正交性：}
\[
\int_0^1 x\, J_\nu(\alpha_m x)\, J_\nu(\alpha_n x)\,dx = 0 \quad (m \neq n)
\]
其中 $\alpha_n$ 是 $J_\nu$ 的第 $n$ 个正零点。

\subsection{Legendre 方程}

\[
(1-x^2)y'' - 2xy' + l(l+1)y = 0
\]

\textbf{来源：} 球坐标下 Laplace 方程的角度部分。

\textbf{Rodrigues 公式：}
\[
P_l(x) = \frac{1}{2^l l!}\frac{d^l}{dx^l}(x^2-1)^l
\]

\textbf{正交性：}
\[
\int_{-1}^1 P_m(x) P_n(x)\,dx = \frac{2}{2n+1}\delta_{mn}
\]

\subsection{特殊函数的统一视角}

所有经典特殊函数都是二阶线性 ODE 的解，且都满足：
\begin{enumerate}
  \item 可写成 Sturm--Liouville 形式
  \item 关于权重函数正交
  \item 构成对应 $L^2_w$ 空间的完备正交基
  \item 满足递推关系
  \item 存在生成函数
  \item 可用超几何函数 $_2F_1$ 表示
\end{enumerate}

\subsection{Green 函数方法（ODE 版本）}

\textbf{问题：} $Ly = f(x)$，$L$ 是线性微分算子，加边界条件。

\textbf{Green 函数 $G(x,\xi)$：} 满足 $LG = \delta(x-\xi)$（加齐次边界条件）。

\textbf{解：} $y(x) = \int_a^b G(x,\xi)\,f(\xi)\,d\xi$

\textbf{示例：} $-y'' = f(x)$，$y(0) = y(1) = 0$

\[
G(x,\xi) = \begin{cases} x(1-\xi) & x < \xi \\ \xi(1-x) & x > \xi \end{cases}
\]

验证：$G$ 在 $x = \xi$ 连续，$G'$ 跳跃 $-1$，$G(0,\xi) = G(1,\xi) = 0$。
%==============================================================
\section{第四层：对称性与李群方法}
%==============================================================

\subsection{单参数李群与无穷小生成元}

\textbf{一参数变换群：} $x^* = X(x,y;\varepsilon)$，$y^* = Y(x,y;\varepsilon)$

\textbf{无穷小生成元：}
\[
\mathbf{v} = \xi(x,y)\frac{\partial}{\partial x} + \eta(x,y)\frac{\partial}{\partial y}
\]

其中 $\xi = \frac{\partial X}{\partial \varepsilon}\big|_{\varepsilon=0}$，$\eta = \frac{\partial Y}{\partial \varepsilon}\big|_{\varepsilon=0}$。

\subsection{对称性条件与延拓}

对 $n$ 阶 ODE $F(x, y, y', \ldots, y^{(n)}) = 0$，对称性条件为：
\[
\mathbf{v}^{(n)}[F]\big|_{F=0} = 0
\]

其中 $\mathbf{v}^{(n)}$ 是 $n$ 次延拓：
\[
\mathbf{v}^{(n)} = \mathbf{v} + \eta^{(1)}\frac{\partial}{\partial y'} + \eta^{(2)}\frac{\partial}{\partial y''} + \cdots
\]

\textbf{一次延拓系数：}
\[
\eta^{(1)} = D_x\eta - y'D_x\xi = \eta_x + (\eta_y - \xi_x)y' - \xi_y(y')^2
\]

\subsection{经典对称性对应的标准解法}

\begin{center}
\begin{tabular}{L{4cm}L{3.5cm}L{4cm}}
\toprule
\textbf{对称性} & \textbf{生成元} & \textbf{对应解法} \\
\midrule
$x$ 平移 & $\partial_x$ & 缺 $x$ → 令 $p = y'$ 降阶 \\
$y$ 平移 & $\partial_y$ & 缺 $y$ → 令 $p = y'$ 降阶 \\
尺度 $(x,y)\to(\lambda x, \lambda y)$ & $x\partial_x + y\partial_y$ & 齐次方程 → $u = y/x$ \\
尺度 $x \to \lambda x$ & $x\partial_x$ & Euler 方程 → $t = \ln x$ \\
\bottomrule
\end{tabular}
\end{center}

\subsection{用对称性降阶的一般步骤}

\begin{enumerate}
  \item 找对称性 $\mathbf{v} = \xi\partial_x + \eta\partial_y$
  \item 找不变量（canonical coordinates）：$r(x,y)$ 使 $\mathbf{v}r = 0$，$s(x,y)$ 使 $\mathbf{v}s = 1$
  \item 在新坐标 $(r, s)$ 下，方程降一阶
\end{enumerate}

\textbf{示例：} 齐次方程 $y' = f(y/x)$

对称性：$\mathbf{v} = x\partial_x + y\partial_y$

不变量：$r = y/x$（验证：$\mathbf{v}(y/x) = x \cdot (-y/x^2) + y \cdot (1/x) = 0$ $\checkmark$）

新变量：$s = \ln x$（验证：$\mathbf{v}(\ln x) = x \cdot (1/x) = 1$ $\checkmark$）

在新坐标下：$\frac{dr}{ds} = f(r) - r$（可分离！）

%==============================================================
\section{第五层：偏微分方程的分类与基本理论}
%==============================================================

\subsection{二阶线性 PDE 的分类}

\[
Au_{xx} + 2Bu_{xy} + Cu_{yy} + \cdots = 0
\]

\textbf{判别式：} $\Delta = B^2 - AC$

\begin{center}
\begin{tabular}{L{2cm}L{2.5cm}L{4cm}L{3cm}}
\toprule
\textbf{$\Delta$} & \textbf{类型} & \textbf{典型方程} & \textbf{物理} \\
\midrule
$> 0$ & 双曲型 & $u_{tt} = c^2 u_{xx}$ & 波传播 \\
$= 0$ & 抛物型 & $u_t = ku_{xx}$ & 扩散 \\
$< 0$ & 椭圆型 & $u_{xx} + u_{yy} = 0$ & 稳态/平衡 \\
\bottomrule
\end{tabular}
\end{center}

\textbf{三类方程的关键区别：}

\begin{center}
\begin{tabular}{L{2.5cm}L{3.5cm}L{3.5cm}L{3.5cm}}
\toprule
\textbf{性质} & \textbf{双曲} & \textbf{抛物} & \textbf{椭圆} \\
\midrule
信息传播 & 有限速度 & 无限速度 & 无“传播” \\
问题类型 & 初值（Cauchy） & 初值+边值 & 纯边值 \\
解的光滑性 & 保持奇性 & 瞬间光滑化 & 解析 \\
依赖区域 & 特征锥内 & 整个过去 & 整个边界 \\
\bottomrule
\end{tabular}
\end{center}

\subsection{三大经典方程}

\subsubsection{波动方程（双曲型）}

\[
\frac{\partial^2 u}{\partial t^2} = c^2 \frac{\partial^2 u}{\partial x^2}
\]

\textbf{d'Alembert 通解：}
\[
u(x,t) = f(x - ct) + g(x + ct)
\]

\textbf{特征线：} $x \pm ct = \text{const}$

\textbf{示例：} 弦振动，$u(x,0) = \sin x$，$u_t(x,0) = 0$

$f(x) + g(x) = \sin x$，$-cf'(x) + cg'(x) = 0$

解得 $f = g = \frac{1}{2}\sin x$，所以：
\[
u(x,t) = \frac{1}{2}[\sin(x-ct) + \sin(x+ct)] = \sin x \cos ct
\]

\subsubsection{热方程（抛物型）}

\[
\frac{\partial u}{\partial t} = k\frac{\partial^2 u}{\partial x^2}
\]

\textbf{基本解（热核）：}
\[
\Phi(x,t) = \frac{1}{\sqrt{4\pi kt}} e^{-x^2/(4kt)}, \quad t > 0
\]

\textbf{性质：}
\begin{itemize}
  \item $\int_{-\infty}^{\infty} \Phi(x,t)\,dx = 1$（质量守恒）
  \item $t \to 0^+$ 时 $\Phi \to \delta(x)$（初始集中）
  \item 瞬间光滑化：即使初始数据不连续，$t > 0$ 后解解析
\end{itemize}

\subsubsection{Laplace 方程（椭圆型）}

\[
\nabla^2 u = 0
\]

\textbf{调和函数的性质：}
\begin{itemize}
  \item 均值性质：$u(\mathbf{x}_0) = \frac{1}{|\partial B_r|}\int_{\partial B_r} u\,dS$
  \item 极值原理：非常数调和函数的极值只能在边界取到
  \item 解析性：调和函数是解析的
\end{itemize}

\subsection{边界条件}

\begin{center}
\begin{tabular}{L{2.5cm}L{4cm}L{5cm}}
\toprule
\textbf{类型} & \textbf{形式} & \textbf{物理意义} \\
\midrule
Dirichlet & $u|_{\partial\Omega} = g$ & 固定值（温度、位移） \\
Neumann & $\frac{\partial u}{\partial n}\big|_{\partial\Omega} = g$ & 固定通量（绝热、自由端） \\
Robin & $u + \alpha\frac{\partial u}{\partial n} = g$ & 对流换热 \\
\bottomrule
\end{tabular}
\end{center}

%==============================================================
\section{第六层：PDE 的求解方法}
%==============================================================

\subsection{分离变量法}

\textbf{核心思想：} 设 $u(x,t) = X(x)T(t)$，将 PDE 分解为 ODE。

\textbf{完整示例：一维热方程}

$u_t = ku_{xx}$，$u(0,t) = u(L,t) = 0$，$u(x,0) = f(x)$

\textbf{Step 1：} 设 $u = XT$，代入：$XT' = kX''T$，即 $\frac{T'}{kT} = \frac{X''}{X} = -\lambda$

\textbf{Step 2：} 解空间部分：$X'' + \lambda X = 0$，$X(0) = X(L) = 0$

特征值：$\lambda_n = \left(\frac{n\pi}{L}\right)^2$，特征函数：$X_n = \sin\frac{n\pi x}{L}$

\textbf{Step 3：} 解时间部分：$T' + k\lambda_n T = 0$，$T_n = e^{-k(n\pi/L)^2 t}$

\textbf{Step 4：} 叠加：
\[
u(x,t) = \sum_{n=1}^{\infty} b_n \sin\frac{n\pi x}{L}\, e^{-k(n\pi/L)^2 t}
\]

\textbf{Step 5：} 由初始条件确定系数：
\[
b_n = \frac{2}{L}\int_0^L f(x)\sin\frac{n\pi x}{L}\,dx
\]

\textbf{本质：} 分离变量 = 在微分算子的特征函数基上展开。

\subsection{特征线法（双曲型）}

\textbf{一阶 PDE：} $a(x,y)u_x + b(x,y)u_y = c(x,y,u)$

\textbf{特征方程：}
\[
\frac{dx}{ds} = a, \quad \frac{dy}{ds} = b, \quad \frac{du}{ds} = c
\]

\textbf{示例：} 输运方程 $u_t + cu_x = 0$，$u(x,0) = f(x)$

特征线：$\frac{dx}{dt} = c$，即 $x - ct = \text{const}$

沿特征线 $u$ 不变，所以 $u(x,t) = f(x - ct)$。

\subsection{积分变换法}

\subsubsection{Fourier 变换}

\textbf{适用：} 无界域上的线性 PDE。

\textbf{示例：} $\mathbb{R}$ 上的热方程 $u_t = ku_{xx}$，$u(x,0) = f(x)$

对 $x$ 做 Fourier 变换：$\hat{u}_t = -k\xi^2 \hat{u}$

解：$\hat{u}(\xi, t) = \hat{f}(\xi)\, e^{-k\xi^2 t}$

反变换（卷积定理）：
\[
u(x,t) = \frac{1}{\sqrt{4\pi kt}}\int_{-\infty}^{\infty} f(y)\, e^{-(x-y)^2/(4kt)}\,dy
\]

\subsubsection{Laplace 变换}

\textbf{适用：} 半无界域、初值问题。

\textbf{核心性质：}
\[
\mathcal{L}\{f'\} = s\hat{f}(s) - f(0)
\]
\[
\mathcal{L}\{f''\} = s^2\hat{f}(s) - sf(0) - f'(0)
\]

将 ODE/PDE 化为代数方程。

\textbf{示例：} 解 $y'' + 4y = \sin 2t$，$y(0) = 0$，$y'(0) = 0$

Laplace 变换：$s^2\hat{y} + 4\hat{y} = \frac{2}{s^2+4}$

\[
\hat{y} = \frac{2}{(s^2+4)^2}
\]

反变换（查表或留数）：
\[
\boxed{y(t) = \frac{1}{4}(\sin 2t - 2t\cos 2t)}
\]

（注意 $t\cos 2t$ 项——这是\textbf{共振}的体现）

\subsection{Green 函数方法（PDE 版本）}

\textbf{Laplace 方程的 Green 函数：}

三维自由空间：
\[
G(\mathbf{x}, \mathbf{x}') = \frac{1}{4\pi|\mathbf{x} - \mathbf{x}'|}
\]

二维自由空间：
\[
G(\mathbf{x}, \mathbf{x}') = -\frac{1}{2\pi}\ln|\mathbf{x} - \mathbf{x}'|
\]

\textbf{解的表示：}
\[
u(\mathbf{x}) = \int_\Omega G(\mathbf{x},\mathbf{x}')\,f(\mathbf{x}')\,d\mathbf{x}' + \oint_{\partial\Omega}\left[G\frac{\partial u}{\partial n} - u\frac{\partial G}{\partial n}\right]dS
\]

\subsection{变分法与弱解}

\textbf{许多 PDE 是某个泛函的 Euler--Lagrange 方程：}

\begin{center}
\begin{tabular}{L{5cm}L{5cm}}
\toprule
\textbf{泛函} & \textbf{Euler--Lagrange 方程} \\
\midrule
$\int |\nabla u|^2\,d\Omega$（Dirichlet 能量） & $\nabla^2 u = 0$ \\
$\int \frac{1}{2}|\nabla u|^2 - fu\,d\Omega$ & $-\nabla^2 u = f$ \\
$\int \frac{1}{2}u_t^2 - \frac{c^2}{2}|\nabla u|^2$ & $u_{tt} = c^2\nabla^2 u$ \\
\bottomrule
\end{tabular}
\end{center}

\textbf{弱形式（$-\nabla^2 u = f$，$u|_{\partial\Omega} = 0$）：}

找 $u \in H_0^1(\Omega)$ 使得对所有 $v \in H_0^1(\Omega)$：
\[
\int_\Omega \nabla u \cdot \nabla v\,d\Omega = \int_\Omega fv\,d\Omega
\]

\textbf{Lax--Milgram 定理} 保证弱解存在唯一（若双线性形式有界且强制）。
%==============================================================
\section{第七层：微分方程的内在联系}
%==============================================================

\subsection{从泰勒展开到 Laplace 变换：统一视角}

\subsubsection{两种“将微分方程化为代数”的方法}

\begin{center}
\begin{tabular}{L{2cm}L{5cm}L{5cm}}
\toprule
 & \textbf{Taylor 展开/幂级数} & \textbf{Laplace 变换} \\
\midrule
信息类型 & 局部（各阶导数在一点） & 全局（指数加权积分） \\
微分算子作用 & $D \to$ 移位（指标 $n \to n-1$） & $D \to$ 乘以 $s$（减去初值） \\
适用域 & 常点/正则奇点附近 & $[0,\infty)$，初值问题 \\
收敛性 & 需要收敛半径 & 需要指数阶增长 \\
复平面 & 收敛圆内解析 & 半平面 $\text{Re}(s) > \sigma$ 解析 \\
\bottomrule
\end{tabular}
\end{center}

\subsubsection{Taylor 展开的算子视角}

Taylor 公式可以写成：
\[
f(x + h) = e^{hD}f(x) = \sum_{n=0}^{\infty} \frac{h^n}{n!}D^n f(x)
\]

其中 $D = d/dx$。\textbf{平移算子就是微分算子的指数。}

\subsubsection{Laplace 变换的算子视角}

$\mathcal{L}$ 将 $D$ 变为“乘以 $s$”（减去初值项）：
\[
\mathcal{L}\{Df\} = s\mathcal{L}\{f\} - f(0)
\]

在“零初值”的理想情况下，$\mathcal{L}$ 将 $D$ 对角化为乘以 $s$。

\textbf{类比：}
\begin{itemize}
  \item Taylor/幂级数：$D$ 在基 $\{x^n/n!\}$ 下是\textbf{移位算子}
  \item Laplace 变换：$D$ 在“指数基” $\{e^{st}\}$ 下是\textbf{乘法算子}（乘以 $s$）
  \item Fourier 变换：$D$ 在“平面波基” $\{e^{i\xi x}\}$ 下是\textbf{乘法算子}（乘以 $i\xi$）
\end{itemize}

\textbf{统一：} 所有积分变换都是在找微分算子的“特征函数基”，使其对角化。

\subsubsection{复平面上的联系}

\textbf{Laplace 变换：} $\hat{f}(s) = \int_0^\infty f(t)e^{-st}\,dt$，$s = \sigma + i\omega$

\textbf{Fourier 变换是 Laplace 变换在 $\sigma = 0$ 时的特例：}
\[
\hat{f}(i\omega) = \int_0^\infty f(t)e^{-i\omega t}\,dt
\]

\textbf{极点与解的行为：}
\begin{itemize}
  \item $\hat{f}(s)$ 在 $s = s_0$ 有极点 → $f(t)$ 含 $e^{s_0 t}$ 项
  \item 极点的实部决定增长/衰减，虚部决定振荡频率
  \item 重极点 → $t^k e^{s_0 t}$（共振！）
\end{itemize}

\textbf{这与特征方程的根完全对应：}
\begin{itemize}
  \item 特征根 $r$ $\leftrightarrow$ Laplace 域的极点 $s = r$
  \item 重根 $\leftrightarrow$ 重极点 $\leftrightarrow$ $t^k e^{rt}$
  \item 复根 $\alpha \pm i\beta$ $\leftrightarrow$ 共轭极点 $\leftrightarrow$ $e^{\alpha t}\cos\beta t$
\end{itemize}

\subsection{幂级数与无穷维矩阵}
\label{subsec:infdim-matrix}

\subsubsection{微分算子的矩阵表示}

\textbf{选择基：} $e_n = \frac{x^n}{n!}$，$n = 0, 1, 2, \ldots$

\textbf{微分算子 $D = d/dx$ 在此基下的作用：}
\[
D e_n = D\frac{x^n}{n!} = \frac{x^{n-1}}{(n-1)!} = e_{n-1} \quad (n \geq 1), \quad De_0 = 0
\]

\textbf{矩阵表示（无穷维）：}
\[
D = \begin{pmatrix} 0 & 1 & 0 & 0 & \cdots \\ 0 & 0 & 1 & 0 & \cdots \\ 0 & 0 & 0 & 1 & \cdots \\ 0 & 0 & 0 & 0 & \cdots \\ \vdots & & & & \ddots \end{pmatrix}
\]

这是\textbf{无穷维移位算子}（将第 $n$ 个分量移到第 $n-1$ 个位置）。

\subsubsection{矩阵指数与 Taylor 公式}

\[
e^{tD} = I + tD + \frac{t^2}{2!}D^2 + \frac{t^3}{3!}D^3 + \cdots
\]

作用在 $f(x) = \sum a_n e_n$ 上：
\[
e^{tD}f(x) = f(x + t)
\]

\textbf{这就是 Taylor 公式！} 平移 = 微分算子的指数。

\subsubsection{ODE 作为无穷维矩阵方程}

\textbf{一阶方程 $y' = f(x)y$：}

设 $y = \sum a_n e_n$，$f(x) = \sum b_n e_n$。

$y' = \sum a_{n+1} e_n$（移位），$f(x)y$ 是卷积（Cauchy 乘积）。

方程变为：
\[
a_{n+1} = \sum_{k=0}^n b_k a_{n-k} \quad \text{（递推关系）}
\]

这可以写成\textbf{无穷维下三角矩阵}乘以向量 $\mathbf{a} = (a_0, a_1, a_2, \ldots)^T$ 的形式。

\textbf{二阶方程 $y'' + P(x)y' + Q(x)y = 0$：}

设 $y = \sum a_n x^n$（用标准幂基），代入得：
\[
(n+2)(n+1)a_{n+2} + \sum_{k=0}^n [\text{coefficients}] a_k = 0
\]

这是\textbf{无穷维带状矩阵}的零空间问题。

\subsubsection{特征值问题的无穷维矩阵视角}

\textbf{Sturm--Liouville 问题} $-y'' = \lambda y$，$y(0) = y(\pi) = 0$：

在基 $\{\sin nx\}$ 下，$-D^2$ 的矩阵是\textbf{对角的}：
\[
-D^2 \sin nx = n^2 \sin nx
\]

特征值 $\lambda_n = n^2$，特征向量就是基向量本身。

\textbf{变系数情况} $-(p(x)y')' + q(x)y = \lambda w(x)y$：

在某个基下，算子变成\textbf{无穷维矩阵}，特征值问题变成无穷维矩阵的特征值问题。

\textbf{截断近似：} 取前 $N$ 个基函数，得到 $N \times N$ 矩阵的特征值问题。这就是\textbf{Galerkin 方法/谱方法}的本质。

\subsubsection{若尔当标准型的无穷维推广}

\textbf{有限维：} 重特征值 → 若尔当块 → $t^k e^{\lambda t}$

\textbf{无穷维：} 微分算子的“广义特征函数”

\textbf{示例：} $(D - r)^2 y = 0$

在基 $\{e^{rx}, xe^{rx}\}$ 下，$D - r$ 的矩阵是 $\begin{pmatrix} 0 & 1 \\ 0 & 0 \end{pmatrix}$（$2\times 2$ 若尔当块）。

推广到无穷维：$(D - r)^k y = 0$ 的解空间由 $\{e^{rx}, xe^{rx}, \ldots, x^{k-1}e^{rx}\}$ 张成。

\subsubsection{生成函数与转移矩阵}

\textbf{递推关系} $a_{n+1} = c \cdot a_n$ 的生成函数：
\[
A(z) = \sum_{n=0}^{\infty} a_n z^n = \frac{a_0}{1 - cz}
\]

\textbf{更一般的递推} $a_{n+2} = pa_{n+1} + qa_n$：
\[
A(z) = \frac{a_0 + (a_1 - pa_0)z}{1 - pz - qz^2}
\]

\textbf{分母的根} $1 - pz - qz^2 = 0$ 与特征方程 $r^2 - pr - q = 0$ 直接相关（$z = 1/r$）。

\textbf{统一：} 生成函数是离散版本的 Laplace 变换（$z$ 变换），递推关系是离散版本的微分方程。

\subsection{多元微分方程求解的系统方法}

\subsubsection{方法总览}

\begin{center}
\begin{tabular}{L{3cm}L{4cm}L{4cm}}
\toprule
\textbf{方法} & \textbf{适用场景} & \textbf{核心思想} \\
\midrule
分离变量 & 有界域，规则几何 & 在特征函数基上展开 \\
特征线法 & 双曲型，一阶 & 沿特征线化为 ODE \\
Green 函数 & 线性，任意源 & 点源响应叠加 \\
积分变换 & 无界/半无界域 & 将微分化为代数 \\
变分/弱形式 & 复杂几何，数值 & 能量最小化 \\
行波/自相似 & 有对称性 & 利用对称性降维 \\
摄动法 & 有小参数 & 渐近展开 \\
\bottomrule
\end{tabular}
\end{center}

\subsubsection{分离变量在不同坐标系}

\begin{center}
\begin{tabular}{L{3cm}L{3.5cm}L{4cm}}
\toprule
\textbf{坐标系} & \textbf{分离形式} & \textbf{特殊函数} \\
\midrule
直角 $(x,y,z)$ & $X(x)Y(y)Z(z)$ & 三角/指数 \\
柱 $(r,\theta,z)$ & $R(r)\Theta(\theta)Z(z)$ & Bessel \\
球 $(r,\theta,\phi)$ & $R(r)Y(\theta,\phi)$ & 球谐 + 球 Bessel \\
\bottomrule
\end{tabular}
\end{center}

\textbf{球坐标下 Laplace 方程的完整分离：}

设 $u = R(r)\Theta(\theta)\Phi(\phi)$，代入 $\nabla^2 u = 0$：

\textbf{$\phi$ 方程：} $\Phi'' + m^2\Phi = 0$ → $\Phi = e^{im\phi}$

\textbf{$\theta$ 方程：} 连带 Legendre 方程 → $P_l^m(\cos\theta)$

\textbf{$r$ 方程：} $r^2 R'' + 2rR' - l(l+1)R = 0$ → $R = Ar^l + Br^{-l-1}$

\textbf{通解：}
\[
u(r,\theta,\phi) = \sum_{l=0}^{\infty}\sum_{m=-l}^{l} \left(A_{lm}r^l + B_{lm}r^{-l-1}\right)Y_l^m(\theta,\phi)
\]

\subsubsection{Green 函数的构造方法}

\textbf{方法一：镜像法}（简单几何）

半空间 $z > 0$，Dirichlet 条件 $u|_{z=0} = 0$：
\[
G(\mathbf{x}, \mathbf{x}') = \frac{1}{4\pi|\mathbf{x}-\mathbf{x}'|} - \frac{1}{4\pi|\mathbf{x}-\mathbf{x}'^*|}
\]
其中 $\mathbf{x}'^*$ 是 $\mathbf{x}'$ 关于 $z=0$ 的镜像。

\textbf{方法二：特征函数展开}

\[
G(\mathbf{x}, \mathbf{x}') = \sum_n \frac{\phi_n(\mathbf{x})\phi_n(\mathbf{x}')}{\lambda_n}
\]

其中 $\phi_n$ 是 $-\nabla^2$ 的特征函数，$\lambda_n$ 是特征值。

\textbf{方法三：Fourier/Laplace 变换}

对无界域，直接变换求解。

\subsection{从 ODE 到 PDE：维度提升的关键变化}

\begin{center}
\begin{tabular}{L{3cm}L{4cm}L{4cm}}
\toprule
\textbf{一维（ODE）} & \textbf{多维（PDE）} & \textbf{新现象} \\
\midrule
两个初始条件 & 整个边界上的条件 & 适定性更复杂 \\
特征点 & 特征线/特征面 & 信息有方向传播 \\
Fourier 级数 & Fourier 变换/球谐 & 基函数更丰富 \\
Green 函数分段 & Green 函数奇异 & $\delta$ 函数、主值 \\
解的奇性孤立 & 奇性沿特征传播 & 激波、焦散 \\
\bottomrule
\end{tabular}
\end{center}

\subsection{微分方程、积分方程与算子方程的统一}

\textbf{微分方程：} $Lu = f$（$L$ 是微分算子）

\textbf{积分方程：} $u = g + \mathcal{K}u$（$\mathcal{K}$ 是积分算子）

\textbf{联系：}
\begin{itemize}
  \item 微分方程 + 边界条件 → 等价于积分方程（通过 Green 函数）
  \item $Lu = f$ → $u = L^{-1}f = \int G(x,\xi)f(\xi)\,d\xi$
  \item 渲染方程 $L = L_e + \mathcal{T}L$ 就是第二类 Fredholm 积分方程
\end{itemize}

\textbf{Neumann 级数（积分方程）：}
\[
u = \sum_{k=0}^{\infty} \mathcal{K}^k g
\]

\textbf{类比：}
\begin{itemize}
  \item ODE 的 Picard 迭代 $\leftrightarrow$ 积分方程的 Neumann 级数 $\leftrightarrow$ 路径追踪的弹射展开
  \item 收敛条件：$\|\mathcal{K}\| < 1$（算子范数）
\end{itemize}
%==============================================================
\section{第八层：数值方法}
\label{sec:de-numerics}
%==============================================================

\subsection{有限差分法}

\subsubsection{基本格式}

\textbf{一阶导数：}
\[
u'(x_i) \approx \frac{u_{i+1} - u_{i-1}}{2h} \quad \text{（中心差分，$O(h^2)$）}
\]

\textbf{二阶导数：}
\[
u''(x_i) \approx \frac{u_{i+1} - 2u_i + u_{i-1}}{h^2} \quad \text{（$O(h^2)$）}
\]

\textbf{二维 Laplace（五点格式）：}
\[
\nabla^2 u_{i,j} \approx \frac{u_{i+1,j} + u_{i-1,j} + u_{i,j+1} + u_{i,j-1} - 4u_{i,j}}{h^2}
\]

\subsubsection{稳定性}

\textbf{热方程显式格式：} $u_j^{n+1} = u_j^n + r(u_{j+1}^n - 2u_j^n + u_{j-1}^n)$，$r = k\Delta t/h^2$

\textbf{稳定性条件：} $r \leq 1/2$

\textbf{波动方程：} $\frac{c\Delta t}{h} \leq 1$（CFL 条件）

\subsection{有限元法}

\textbf{核心步骤：}
\begin{enumerate}
  \item 将域 $\Omega$ 剖分为单元
  \item 在每个单元上用多项式基函数近似
  \item 代入弱形式，组装全局刚度矩阵
  \item 施加边界条件
  \item 求解线性系统 $K\mathbf{u} = \mathbf{f}$
\end{enumerate}

\textbf{刚度矩阵元素：}
\[
K_{ij} = \int_\Omega \nabla\phi_i \cdot \nabla\phi_j\,d\Omega
\]

\subsection{谱方法}

\textbf{思想：} 用全局正交基展开 $u(x) = \sum \hat{u}_n \phi_n(x)$

\textbf{Galerkin 方法：} 将残差投影到基函数上为零。

\textbf{优势：} 对光滑解指数收敛。

\textbf{劣势：} 对不连续解有 Gibbs 现象。

\textbf{速查版本见~\ref{subsec:pde-numerics}~节。}

\subsection{半拉格朗日方法}

对 $u_t + \mathbf{v}\cdot\nabla u = 0$：
\begin{enumerate}
  \item 从网格点 $\mathbf{x}$ 出发，沿速度场回溯：$\mathbf{x}^* = \mathbf{x} - \Delta t\,\mathbf{v}(\mathbf{x})$
  \item 在 $\mathbf{x}^*$ 处插值：$u^{n+1}(\mathbf{x}) = u^n(\mathbf{x}^*)$
\end{enumerate}

\textbf{优势：} 无条件稳定。\textbf{劣势：} 数值耗散。

%==============================================================
\section{第九层：与物理/CG 的交叉应用}
%==============================================================

\subsection{微分方程在渲染中的出现}

\begin{center}
\begin{tabular}{L{4cm}L{5cm}}
\toprule
\textbf{方程} & \textbf{渲染中的角色} \\
\midrule
Helmholtz 方程 & 波动光学、衍射 \\
辐射传输方程（RTE） & 参与介质中的光传输 \\
渲染方程（积分方程） & 全局光照 \\
扩散方程 & 次表面散射近似 \\
Poisson 方程 & 压力投影、法线重建 \\
\bottomrule
\end{tabular}
\end{center}

\subsection{微分方程在物理模拟中的出现}

\begin{center}
\begin{tabular}{L{4cm}L{5cm}}
\toprule
\textbf{方程} & \textbf{模拟对象} \\
\midrule
Navier--Stokes & 流体（水、烟、火） \\
弹性方程 & 固体变形 \\
波动方程 & 布料、绳索、声波 \\
热方程 & 温度场、扩散 \\
Hamilton 方程 & 刚体、粒子系统 \\
反应--扩散方程 & 图案生成（Turing） \\
\bottomrule
\end{tabular}
\end{center}

\subsection{微分方程在几何处理中的出现}

\begin{center}
\begin{tabular}{L{4cm}L{5cm}}
\toprule
\textbf{方程} & \textbf{几何处理任务} \\
\midrule
$\Delta u = 0$（Laplace） & 调和映射、网格平滑 \\
$\Delta^2 u = 0$（双调和） & 薄板弯曲、曲面重建 \\
平均曲率流 & 网格平滑、曲面演化 \\
测地线方程 & 最短路径、形状匹配 \\
$-\Delta u = \lambda u$ & 形状分析、谱聚类 \\
\bottomrule
\end{tabular}
\end{center}

%==============================================================
\section{总结：微分方程的统一视角}
%==============================================================

\subsection{五条核心线索}

\begin{enumerate}
  \item \textbf{对称性 → 降阶/简化}（李群方法、分离变量、行波解、自相似解）
  \item \textbf{线性算子的谱 → 正交基 → 展开}（Fourier、球谐、Bessel、Legendre）
  \item \textbf{Green 函数 → 点源响应 → 叠加}（线性问题的通用解法）
  \item \textbf{变分原理 → 弱形式 → 数值离散}（有限元的理论基础）
  \item \textbf{特征值问题 → 振动模式 → 形状/频率分析}（贯穿所有领域）
\end{enumerate}

\subsection{知识网络}

\[
\boxed{
\begin{aligned}
&\text{对称性（李群）} \xrightarrow{\text{降阶/降维}} \text{简化方程} \\
&\text{分离变量} \xrightarrow{\text{特征值问题}} \text{特殊函数（正交基）} \\
&\text{Green 函数} \xrightarrow{\text{点源叠加}} \text{卷积/积分表示} \\
&\text{变分原理} \xrightarrow{\text{弱形式}} \text{有限元/数值解} \\
&\text{Taylor/幂级数} \xrightarrow{\text{无穷维矩阵}} \text{算子对角化/积分变换}
\end{aligned}
}
\]

\subsection{从局部到全局的三条路径}

\[
\boxed{
\begin{aligned}
&\text{Taylor 展开（局部）} \xrightarrow{\text{解析延拓}} \text{复平面上的全局信息} \\
&\text{Laplace 变换（半轴）} \xrightarrow{\sigma = 0} \text{Fourier 变换（全轴）} \\
&\text{幂级数（离散指标）} \xrightarrow{\text{生成函数}} \text{复变函数（连续变量）}
\end{aligned}
}
\]

\subsection{微分方程 $\leftrightarrow$ 代数 $\leftrightarrow$ 几何的三角关系}

\[
\boxed{
\begin{aligned}
&\text{微分方程} \xleftrightarrow{\text{特征值/谱}} \text{线性代数（矩阵）} \\
&\text{微分方程} \xleftrightarrow{\text{对称性/不变量}} \text{几何（流形/群）} \\
&\text{线性代数} \xleftrightarrow{\text{表示论}} \text{几何（群作用）}
\end{aligned}
}
\]

\subsection{最终统一}

所有方法归结为一个核心问题：

\[
\boxed{Lu = f}
\]

其中 $L$ 是（微分/积分/差分）算子。求解策略：

\begin{enumerate}
  \item \textbf{找 $L$ 的特征基} → 对角化 → 逐分量求解（分离变量、谱方法）
  \item \textbf{找 $L^{-1}$（Green 函数）} → 卷积 → 叠加（Green 函数法、边界元法）
  \item \textbf{利用对称性} → 降维 → 简化（李群方法、行波解）
  \item \textbf{变分/能量最小化} → 弱形式 → 离散（有限元）
  \item \textbf{迭代/逼近} → 逐步收敛（Picard 迭代、Neumann 级数、路径追踪）
\end{enumerate}

这五条路径覆盖了从考研解题到前沿研究的所有微分方程方法。

\clearpage
\part{专题二：李群、李代数与拓扑}
%==============================================================
\section{李群的基本概念}
\label{sec:lie-group}
%==============================================================

\subsection{定义}

\textbf{李群} $G$ 是一个同时具有群结构和光滑流形结构的对象，使得群乘法 $G \times G \to G$ 和取逆 $G \to G$ 都是光滑映射。

\textbf{维度：} 李群作为流形的维度。

\subsection{经典李群}

\begin{center}
\begin{tabular}{L{2cm}L{5cm}L{1.5cm}L{4cm}}
\toprule
\textbf{群} & \textbf{定义} & \textbf{维度} & \textbf{应用} \\
\midrule
$GL(n,\mathbb{R})$ & $n\times n$ 可逆实矩阵 & $n^2$ & 一般线性变换 \\
$SL(n,\mathbb{R})$ & $\det A = 1$ & $n^2 - 1$ & 保体积变换 \\
$O(n)$ & $A^TA = I$ & $\frac{n(n-1)}{2}$ & 保距（含反射） \\
$SO(n)$ & $A^TA = I$，$\det A = 1$ & $\frac{n(n-1)}{2}$ & 旋转 \\
$U(n)$ & $A^\dagger A = I$ & $n^2$ & 量子力学 \\
$SU(n)$ & $A^\dagger A = I$，$\det A = 1$ & $n^2 - 1$ & 规范场论 \\
$Sp(n)$ & $A^\dagger A = I$，保辛形式 & $n(2n+1)$ & 哈密顿力学 \\
$SE(n)$ & $\begin{pmatrix} R & t \\ 0 & 1 \end{pmatrix}$，$R \in SO(n)$ & $\frac{n(n+1)}{2}$ & 刚体运动 \\
\bottomrule
\end{tabular}
\end{center}

\subsection{关键低维例子}

\subsubsection{$SO(2)$：平面旋转}

\[
SO(2) = \left\{\begin{pmatrix} \cos\theta & -\sin\theta \\ \sin\theta & \cos\theta \end{pmatrix} : \theta \in \mathbb{R}\right\}
\]

\begin{itemize}
  \item 维度：1（一个参数 $\theta$）
  \item 拓扑：$S^1$（圆）
  \item 阿贝尔群（交换）
  \item 同构于 $U(1)$（复数单位圆）
\end{itemize}

\subsubsection{$SO(3)$：三维旋转}

\[
SO(3) = \{R \in \mathbb{R}^{3\times 3} : R^TR = I,\; \det R = 1\}
\]

\begin{itemize}
  \item 维度：3（三个参数，如欧拉角或轴角）
  \item 拓扑：$\mathbb{RP}^3$（三维实射影空间）
  \item 非阿贝尔群
  \item 基本群：$\pi_1(SO(3)) = \mathbb{Z}_2$（旋转 $2\pi$ 不等于恒等！）
\end{itemize}

\subsubsection{$SU(2)$：自旋群}

\[
SU(2) = \{A \in \mathbb{C}^{2\times 2} : A^\dagger A = I,\; \det A = 1\}
\]

元素可以写成：
\[
A = \begin{pmatrix} \alpha & -\bar{\beta} \\ \beta & \bar{\alpha} \end{pmatrix}, \quad |\alpha|^2 + |\beta|^2 = 1
\]

\begin{itemize}
  \item 维度：3
  \item 拓扑：$S^3$（三维球面）
  \item 单连通：$\pi_1(SU(2)) = 0$
  \item 同构于单位四元数群 $\text{Sp}(1)$
\end{itemize}

\textbf{关键关系：} $SU(2) \xrightarrow{2:1} SO(3)$（二重覆盖）

\textbf{物理意义：}
\begin{itemize}
  \item 电子自旋 $1/2$：旋转 $2\pi$ 后波函数变号，旋转 $4\pi$ 才回到自身
  \item 这就是为什么 $SU(2)$（而非 $SO(3)$）是自旋的正确对称群
\end{itemize}

\subsubsection{$SE(3)$：三维刚体运动}

\[
SE(3) = \left\{\begin{pmatrix} R & \mathbf{t} \\ \mathbf{0}^T & 1 \end{pmatrix} : R \in SO(3),\; \mathbf{t} \in \mathbb{R}^3\right\}
\]

\begin{itemize}
  \item 维度：6（3 旋转 + 3 平移）
  \item 拓扑：$SO(3) \times \mathbb{R}^3$（作为流形）
  \item 半直积结构：$SE(3) = SO(3) \ltimes \mathbb{R}^3$
\end{itemize}

\textbf{在 CG 中：} 刚体变换、骨骼动画、相机运动。

\subsection{李群的同态与覆盖}

\textbf{覆盖映射：} $\phi: \tilde{G} \to G$ 是局部同胚的满同态。

\textbf{万有覆盖群：} 单连通的覆盖群。

\begin{center}
\begin{tabular}{L{3cm}L{3cm}L{3cm}L{4cm}}
\toprule
\textbf{群} & \textbf{万有覆盖} & \textbf{覆盖映射} & \textbf{核} \\
\midrule
$SO(3)$ & $SU(2) \cong S^3$ & 2:1 & $\{\pm I\} \cong \mathbb{Z}_2$ \\
$SO(2) \cong S^1$ & $\mathbb{R}$ & $\theta \mapsto e^{i\theta}$ & $\mathbb{Z}$ \\
$SO(1,3)^+$（Lorentz） & $SL(2,\mathbb{C})$ & 2:1 & $\{\pm I\}$ \\
\bottomrule
\end{tabular}
\end{center}

%==============================================================
\section{李代数}
\label{sec:lie-algebra}
%==============================================================

\subsection{定义}

\textbf{李代数} $\mathfrak{g}$ 是向量空间，配备双线性运算 $[\cdot,\cdot]: \mathfrak{g} \times \mathfrak{g} \to \mathfrak{g}$（李括号），满足：
\begin{enumerate}
  \item 反对称：$[X, Y] = -[Y, X]$
  \item Jacobi 恒等式：$[X, [Y, Z]] + [Y, [Z, X]] + [Z, [X, Y]] = 0$
\end{enumerate}

\textbf{李群与李代数的关系：} 李代数 $\mathfrak{g} = T_e G$（李群在单位元处的切空间），李括号来自向量场的对易子。

\subsection{结构常数}

选基 $\{T_a\}$，$a = 1, \ldots, \dim\mathfrak{g}$：
\[
[T_a, T_b] = f_{ab}^{\;\;c}\, T_c
\]

$f_{ab}^{\;\;c}$ 是\textbf{结构常数}，完全决定了李代数的结构。

\textbf{Jacobi 恒等式的约束：}
\[
f_{ab}^{\;\;d}f_{dc}^{\;\;e} + f_{bc}^{\;\;d}f_{da}^{\;\;e} + f_{ca}^{\;\;d}f_{db}^{\;\;e} = 0
\]

\subsection{指数映射}

\[
\exp: \mathfrak{g} \to G, \quad X \mapsto e^X = \sum_{k=0}^{\infty} \frac{X^k}{k!}
\]

\textbf{性质：}
\begin{itemize}
  \item $\exp(0) = e$（单位元）
  \item $\exp(tX)$ 是过单位元的一参数子群
  \item 在单位元附近是局部微分同胚
  \item 一般不是满射（如 $SL(2,\mathbb{R})$）
\end{itemize}

\textbf{Baker--Campbell--Hausdorff 公式：}
\[
\log(e^X e^Y) = X + Y + \frac{1}{2}[X,Y] + \frac{1}{12}[X,[X,Y]] - \frac{1}{12}[Y,[X,Y]] + \cdots
\]

\textbf{意义：} 群乘法在李代数层面由李括号完全决定（局部）。

\subsection{经典李代数}

\subsubsection{$\mathfrak{so}(3)$}

\[
\mathfrak{so}(3) = \{A \in \mathbb{R}^{3\times 3} : A^T = -A\}
\]

基：
\[
J_1 = \begin{pmatrix} 0&0&0\\0&0&-1\\0&1&0 \end{pmatrix}, \quad
J_2 = \begin{pmatrix} 0&0&1\\0&0&0\\-1&0&0 \end{pmatrix}, \quad
J_3 = \begin{pmatrix} 0&-1&0\\1&0&0\\0&0&0 \end{pmatrix}
\]

对易关系：$[J_i, J_j] = \varepsilon_{ijk}J_k$

\textbf{与角动量的关系：} 量子力学中 $[\hat{L}_i, \hat{L}_j] = i\hbar\varepsilon_{ijk}\hat{L}_k$，正是 $\mathfrak{so}(3)$ 的结构（乘以 $i\hbar$）。

\subsubsection{$\mathfrak{su}(2)$}

\[
\mathfrak{su}(2) = \{A \in \mathbb{C}^{2\times 2} : A^\dagger = -A,\; \text{tr}\,A = 0\}
\]

基（Pauli 矩阵除以 $2$）：
\[
\sigma_1 = \begin{pmatrix}0&1\\1&0\end{pmatrix}, \quad
\sigma_2 = \begin{pmatrix}0&-i\\i&0\end{pmatrix}, \quad
\sigma_3 = \begin{pmatrix}1&0\\0&-1\end{pmatrix}
\]

$T_a = \frac{1}{2}\sigma_a$，对易关系：$[T_a, T_b] = i\varepsilon_{abc}T_c$

\textbf{$\mathfrak{su}(2) \cong \mathfrak{so}(3)$}（作为李代数同构），但 $SU(2) \not\cong SO(3)$（作为群不同构，拓扑不同）。

\subsubsection{$\mathfrak{so}(1,3)$：Lorentz 代数}

\[
\mathfrak{so}(1,3) = \{A : A^T\eta + \eta A = 0\}, \quad \eta = \text{diag}(-1,1,1,1)
\]

6 个生成元：3 个旋转 $J_i$ + 3 个 boost $K_i$

对易关系：
\begin{align*}
[J_i, J_j] &= i\varepsilon_{ijk}J_k \\
[J_i, K_j] &= i\varepsilon_{ijk}K_k \\
[K_i, K_j] &= -i\varepsilon_{ijk}J_k
\end{align*}

\textbf{复化后分解：} $\mathfrak{so}(1,3)_\mathbb{C} \cong \mathfrak{su}(2)_\mathbb{C} \oplus \mathfrak{su}(2)_\mathbb{C}$

令 $A_i = \frac{1}{2}(J_i + iK_i)$，$B_i = \frac{1}{2}(J_i - iK_i)$，则 $[A_i, A_j] = i\varepsilon_{ijk}A_k$，$[B_i, B_j] = i\varepsilon_{ijk}B_k$，$[A_i, B_j] = 0$。

\textbf{物理意义：} 这解释了为什么粒子按 $(j_1, j_2)$ 分类（两个“自旋”量子数）：
\begin{itemize}
  \item $(0, 0)$：标量（Higgs）
  \item $(\frac{1}{2}, 0)$：左手 Weyl 旋量
  \item $(0, \frac{1}{2})$：右手 Weyl 旋量
  \item $(\frac{1}{2}, \frac{1}{2})$：向量（光子）
  \item $(1, 0)$ 或 $(0, 1)$：反对称张量
\end{itemize}

\subsubsection{$\mathfrak{su}(n)$}

\[
\mathfrak{su}(n) = \{A \in \mathbb{C}^{n\times n} : A^\dagger = -A,\; \text{tr}\,A = 0\}
\]

维度：$n^2 - 1$

\begin{itemize}
  \item $\mathfrak{su}(2)$：3 维，弱相互作用
  \item $\mathfrak{su}(3)$：8 维，强相互作用（量子色动力学）
\end{itemize}

\subsection{Killing 形式与半单性}

\textbf{Killing 形式：}
\[
B(X, Y) = \text{tr}(\text{ad}_X \circ \text{ad}_Y)
\]

其中 $(\text{ad}_X)(Y) = [X, Y]$。

\textbf{Cartan 判据：} $\mathfrak{g}$ 半单 $\iff$ $B$ 非退化。

\textbf{紧致性：} $\mathfrak{g}$ 是紧致李代数 $\iff$ $B$ 负定。

\subsection{伴随表示}

\textbf{伴随表示：} $\text{Ad}: G \to GL(\mathfrak{g})$，$\text{Ad}_g(X) = gXg^{-1}$

\textbf{无穷小版本：} $\text{ad}: \mathfrak{g} \to \mathfrak{gl}(\mathfrak{g})$，$\text{ad}_X(Y) = [X, Y]$

\textbf{在物理中：} 规范场（胶子、$W$ 玻色子）在伴随表示下变换。
%==============================================================
\section{李代数的分类}
%==============================================================

\subsection{单李代数与半单李代数}

\textbf{单李代数：} 没有非平凡理想（$[I, \mathfrak{g}] \subset I$ 的子空间 $I$）。

\textbf{半单李代数：} 单李代数的直和。

\textbf{分类定理（Cartan--Killing）：} 复半单李代数完全由 Dynkin 图分类。

\subsection{经典李代数系列}

\begin{center}
\begin{tabular}{L{2cm}L{3cm}L{2cm}L{5cm}}
\toprule
\textbf{类型} & \textbf{李代数} & \textbf{维度} & \textbf{物理} \\
\midrule
$A_n$ & $\mathfrak{sl}(n+1, \mathbb{C})$ & $n(n+2)$ & $SU(n+1)$：色、味对称 \\
$B_n$ & $\mathfrak{so}(2n+1, \mathbb{C})$ & $n(2n+1)$ & $SO(2n+1)$：旋转 \\
$C_n$ & $\mathfrak{sp}(2n, \mathbb{C})$ & $n(2n+1)$ & $Sp(n)$：辛对称 \\
$D_n$ & $\mathfrak{so}(2n, \mathbb{C})$ & $n(2n-1)$ & $SO(2n)$：旋转 \\
\bottomrule
\end{tabular}
\end{center}

\subsection{例外李代数}

\begin{center}
\begin{tabular}{L{2cm}L{2cm}L{6cm}}
\toprule
\textbf{类型} & \textbf{维度} & \textbf{出现} \\
\midrule
$G_2$ & 14 & 八元数的自同构群；弦论 \\
$F_4$ & 52 & 例外 Jordan 代数；弦论 \\
$E_6$ & 78 & 大统一理论候选 \\
$E_7$ & 133 & 超引力；弦论 \\
$E_8$ & 248 & 弦论（$E_8 \times E_8$ 杂化弦）；凝聚态（Ising 模型） \\
\bottomrule
\end{tabular}
\end{center}

\subsection{根系与 Dynkin 图}

\subsubsection{Cartan 子代数与根}

\textbf{Cartan 子代数 $\mathfrak{h}$：} 极大交换子代数（所有元素可同时对角化）。

\textbf{根：} $\alpha \in \mathfrak{h}^*$ 使得存在 $E_\alpha \neq 0$ 满足 $[H, E_\alpha] = \alpha(H)E_\alpha$，$\forall H \in \mathfrak{h}$。

\textbf{分解：}
\[
\mathfrak{g} = \mathfrak{h} \oplus \bigoplus_{\alpha \in \Phi} \mathfrak{g}_\alpha
\]

\subsubsection{Dynkin 图}

\textbf{构造规则：}
\begin{itemize}
  \item 每个单根画一个节点
  \item 两节点间的边数由夹角决定：$90° \to 0$ 边，$60° \to 1$ 边，$45° \to 2$ 边，$30° \to 3$ 边
  \item 不等长根之间加箭头（指向短根）
\end{itemize}

\textbf{经典系列的 Dynkin 图：}

\begin{itemize}
  \item $A_n$：$\circ - \circ - \circ - \cdots - \circ$（$n$ 个节点，链状）
  \item $B_n$：$\circ - \circ - \cdots - \circ \Rightarrow \circ$（末端双键）
  \item $C_n$：$\circ - \circ - \cdots - \circ \Leftarrow \circ$（末端双键，方向反）
  \item $D_n$：$\circ - \circ - \cdots - \circ$ 末端分叉为两个节点
\end{itemize}

\textbf{例外系列：} $G_2$、$F_4$、$E_6$、$E_7$、$E_8$ 有特殊的 Dynkin 图。

\subsection{表示论}

\subsubsection{权与最高权}

\textbf{权：} 表示空间中 $H \in \mathfrak{h}$ 的特征值。

\textbf{最高权：} 在某个偏序下最大的权。

\textbf{定理：} 不可约表示由最高权唯一确定。

\subsubsection{$\mathfrak{su}(2)$ 的表示}

不可约表示由非负整数/半整数 $j$ 标记，维度 $2j + 1$：
\begin{itemize}
  \item $j = 0$：标量（1 维）
  \item $j = 1/2$：旋量（2 维）——电子自旋
  \item $j = 1$：向量（3 维）——光子
  \item $j = 3/2$：（4 维）——$\Delta$ 重子
  \item $j = 2$：（5 维）——引力子
\end{itemize}

\subsubsection{$\mathfrak{su}(3)$ 的表示与粒子物理}

\textbf{基本表示 $\mathbf{3}$：} 夸克（$u, d, s$）

\textbf{共轭表示 $\bar{\mathbf{3}}$：} 反夸克

\textbf{分解：}
\[
\mathbf{3} \otimes \bar{\mathbf{3}} = \mathbf{8} \oplus \mathbf{1}
\]

$\mathbf{8}$：8 个胶子（伴随表示）；$\mathbf{1}$：单态。

\[
\mathbf{3} \otimes \mathbf{3} \otimes \mathbf{3} = \mathbf{10} \oplus \mathbf{8} \oplus \mathbf{8} \oplus \mathbf{1}
\]

$\mathbf{10}$：重子十重态（$\Delta$、$\Sigma^*$、$\Xi^*$、$\Omega^-$）——这是盖尔曼预言 $\Omega^-$ 的基础。

%==============================================================
\section{拓扑结构}
%==============================================================

\subsection{基本群}

\textbf{基本群 $\pi_1(X)$：} 闭合路径的同伦类。

\begin{center}
\begin{tabular}{L{3cm}L{3cm}L{5cm}}
\toprule
\textbf{空间} & \textbf{$\pi_1$} & \textbf{物理意义} \\
\midrule
$S^1$ & $\mathbb{Z}$ & 绕数（Aharonov--Bohm 效应） \\
$SO(3) \cong \mathbb{RP}^3$ & $\mathbb{Z}_2$ & 自旋 $1/2$（$2\pi$ 旋转 $\neq$ 恒等） \\
$SU(2) \cong S^3$ & $0$（单连通） & 万有覆盖 \\
$U(1) \cong S^1$ & $\mathbb{Z}$ & 磁通量子化、涡旋 \\
$SO(3,1)^+$（Lorentz） & $\mathbb{Z}_2$ & 旋量 \\
\bottomrule
\end{tabular}
\end{center}

\subsection{高阶同伦群}

\textbf{$\pi_n(X)$：} $S^n \to X$ 的连续映射的同伦类。

\begin{center}
\begin{tabular}{L{3cm}L{3cm}L{5cm}}
\toprule
\textbf{同伦群} & \textbf{值} & \textbf{物理意义} \\
\midrule
$\pi_2(S^2)$ & $\mathbb{Z}$ & 磁单极（'t Hooft--Polyakov） \\
$\pi_3(S^2)$ & $\mathbb{Z}$ & Skyrmion（核子模型） \\
$\pi_3(SU(2))$ & $\mathbb{Z}$ & 瞬子（Yang--Mills） \\
$\pi_3(SU(3))$ & $\mathbb{Z}$ & QCD 真空结构（$\theta$ 真空） \\
$\pi_1(U(1))$ & $\mathbb{Z}$ & 涡旋（超流、超导） \\
\bottomrule
\end{tabular}
\end{center}

\subsection{纤维丛与拓扑}

\textbf{主 $G$-丛：} $P \xrightarrow{\pi} M$，纤维是群 $G$。

\textbf{分类：} 主 $G$-丛的同构类由 $[M, BG]$ 给出（$BG$ 是分类空间）。

\textbf{特征类：}
\begin{itemize}
  \item $U(1)$ 丛：第一 Chern 类 $c_1 \in H^2(M, \mathbb{Z})$ → 磁通量子化
  \item $SU(n)$ 丛：Chern 类 $c_k$ → 瞬子数
  \item $SO(n)$ 丛：Pontryagin 类 $p_k$、Euler 类 $e$
\end{itemize}

\textbf{在物理中：}
\begin{itemize}
  \item 电磁场是 $U(1)$ 主丛上的联络
  \item Yang--Mills 场是 $SU(n)$ 主丛上的联络
  \item 瞬子数 = 第二 Chern 数 = $\frac{1}{8\pi^2}\int \text{tr}(F \wedge F)$
\end{itemize}

\subsection{de Rham 上同调}

\textbf{de Rham 上同调：}
\[
H^k_{dR}(M) = \frac{\ker(d: \Omega^k \to \Omega^{k+1})}{\text{im}(d: \Omega^{k-1} \to \Omega^k)} = \frac{\text{闭形式}}{\text{恰当形式}}
\]

\textbf{物理意义：}
\begin{itemize}
  \item $H^2_{dR}(M) \neq 0$：存在非平凡的电磁场构型（磁单极）
  \item $H^1_{dR}(M) \neq 0$：存在非平凡的矢势（Aharonov--Bohm）
\end{itemize}

\subsection{同调与流形}

\textbf{Euler 示性数：}
\[
\chi(M) = \sum_{k=0}^n (-1)^k b_k
\]

其中 $b_k = \dim H^k(M)$ 是 Betti 数。

\textbf{Gauss--Bonnet 定理：}
\[
\int_M K\,dA = 2\pi\chi(M)
\]

（$K$ 是高斯曲率，$M$ 是闭曲面）

\textbf{在 CG 中：} 网格的拓扑不变量（亏格）影响参数化和网格生成。
%==============================================================
\section{在物理学中的应用}
%==============================================================

\subsection{量子力学中的对称性}

\subsubsection{角动量与 $\mathfrak{so}(3)$}

\textbf{角动量算子满足 $\mathfrak{so}(3)$ 代数：}
\[
[\hat{L}_i, \hat{L}_j] = i\hbar\varepsilon_{ijk}\hat{L}_k
\]

\textbf{不可约表示 → 角动量量子化：}
\begin{itemize}
  \item $\hat{L}^2|l,m\rangle = \hbar^2 l(l+1)|l,m\rangle$
  \item $\hat{L}_z|l,m\rangle = \hbar m|l,m\rangle$
  \item $l = 0, 1, 2, \ldots$；$m = -l, \ldots, l$
\end{itemize}

\textbf{球谐函数 $Y_l^m$ 是 $SO(3)$ 的第 $l$ 个不可约表示的基函数。}

\subsubsection{自旋与 $\mathfrak{su}(2)$}

\textbf{自旋算子满足 $\mathfrak{su}(2)$ 代数：}
\[
[\hat{S}_i, \hat{S}_j] = i\hbar\varepsilon_{ijk}\hat{S}_k
\]

电子自旋 $s = 1/2$：$SU(2)$ 的 2 维表示。

\textbf{为什么是 $SU(2)$ 而非 $SO(3)$？} 因为 $\pi_1(SO(3)) = \mathbb{Z}_2$，$SO(3)$ 的投影表示（允许相位）等价于 $SU(2)$ 的普通表示。

\subsubsection{氢原子与 $SO(4)$}

氢原子有隐藏对称性：除了角动量 $\mathbf{L}$，还有 Runge--Lenz 向量 $\mathbf{A}$。

$\mathbf{L}$ 和 $\mathbf{A}$ 共同生成 $\mathfrak{so}(4) \cong \mathfrak{su}(2) \oplus \mathfrak{su}(2)$。

\textbf{结果：} 能级 $E_n$ 的简并度是 $n^2$（而非仅 $2l+1$），这由 $SO(4)$ 的不可约表示解释。

\subsection{规范场论}

\subsubsection{规范对称性 → 力}

\textbf{核心思想：} 要求拉格朗日量在\textbf{局部}规范变换下不变 → 必须引入规范场（联络）→ 规范场就是力。

\begin{center}
\begin{tabular}{L{2cm}L{3cm}L{3cm}L{4cm}}
\toprule
\textbf{规范群} & \textbf{规范场} & \textbf{力} & \textbf{粒子} \\
\midrule
$U(1)$ & $A_\mu$（1 个） & 电磁力 & 光子 \\
$SU(2)$ & $W_\mu^a$（3 个） & 弱力 & $W^\pm, Z^0$ \\
$SU(3)$ & $G_\mu^a$（8 个） & 强力 & 8 个胶子 \\
\bottomrule
\end{tabular}
\end{center}

\subsubsection{标准模型的群结构}

\[
G_{SM} = SU(3)_C \times SU(2)_L \times U(1)_Y
\]

\begin{itemize}
  \item $SU(3)_C$：色对称（强相互作用）
  \item $SU(2)_L$：弱同位旋（弱相互作用）
  \item $U(1)_Y$：弱超荷
\end{itemize}

\textbf{电弱对称性破缺：} $SU(2)_L \times U(1)_Y \xrightarrow{\text{Higgs}} U(1)_{EM}$

\subsubsection{规范场的数学结构}

\textbf{规范场 = 主丛上的联络：}
\[
A = A_\mu^a T_a\, dx^\mu
\]

\textbf{场强 = 曲率：}
\[
F = dA + A \wedge A = \frac{1}{2}F_{\mu\nu}^a T_a\, dx^\mu \wedge dx^\nu
\]

\textbf{Yang--Mills 作用量：}
\[
S_{YM} = -\frac{1}{2}\int \text{tr}(F \wedge *F) = -\frac{1}{4}\int F_{\mu\nu}^a F^{a\mu\nu}\,d^4x
\]

\subsection{广义相对论与时空对称性}

\textbf{时空对称群：}
\begin{itemize}
  \item 牛顿力学：Galilei 群（10 维）
  \item 狭义相对论：Poincaré 群 $ISO(1,3)$（10 维）
  \item 广义相对论：微分同胚群 $\text{Diff}(M)$（无穷维）
\end{itemize}

\textbf{Poincaré 群 = Lorentz 群 + 平移：}
\[
ISO(1,3) = SO(1,3) \ltimes \mathbb{R}^4
\]

\textbf{Wigner 分类：} 粒子 = Poincaré 群的不可约酉表示，由 $(m, s)$ 标记（质量、自旋）。

\subsection{凝聚态物理}

\begin{center}
\begin{tabular}{L{4cm}L{4cm}L{4cm}}
\toprule
\textbf{对称性} & \textbf{破缺} & \textbf{结果} \\
\midrule
$U(1)$（相位） & 超导 & 磁通量子化、Josephson 效应 \\
$SO(3)$（自旋） & 铁磁 & 磁畴、自旋波 \\
平移对称性 & 晶体 & 声子、能带结构 \\
$SU(2) \times SU(2)$ & 手征对称性破缺 & 核子质量 \\
\bottomrule
\end{tabular}
\end{center}

\textbf{Goldstone 定理：} 连续对称性自发破缺 → 无质量 Goldstone 玻色子。

\textbf{Higgs 机制：} 规范对称性自发破缺 → 规范玻色子获得质量（Goldstone 玻色子被“吃掉”）。

%==============================================================
\section{在微分方程中的应用}
%==============================================================

\subsection{对称性方法求解 ODE}

（详见微分方程专题第四层，此处总结要点）

\begin{center}
\begin{tabular}{L{4cm}L{4cm}L{4cm}}
\toprule
\textbf{对称性} & \textbf{生成元} & \textbf{效果} \\
\midrule
$x$ 平移 & $\partial_x$ & 降阶（缺 $x$） \\
$y$ 平移 & $\partial_y$ & 降阶（缺 $y$） \\
尺度变换 & $x\partial_x + y\partial_y$ & 齐次方程降阶 \\
旋转 & $-y\partial_x + x\partial_y$ & 极坐标简化 \\
\bottomrule
\end{tabular}
\end{center}

\subsection{对称性方法求解 PDE}

\textbf{热方程 $u_t = ku_{xx}$ 的对称群（6 维）：}
\begin{itemize}
  \item $x$ 平移：$\partial_x$
  \item $t$ 平移：$\partial_t$
  \item $u$ 的线性叠加：$u\partial_u$
  \item 尺度：$(2t\partial_t + x\partial_x)$
  \item Galilei boost：$(2kt\partial_x + xu\partial_u)$
  \item 特殊共形：$(4k^2t^2\partial_t + 4ktx\partial_x + (2kt - x^2)u\partial_u)$
\end{itemize}

\textbf{利用对称性构造特殊解：}
\begin{itemize}
  \item 平移不变 → 行波解 $u = f(x - ct)$
  \item 尺度不变 → 自相似解 $u = t^{-\alpha}f(x/\sqrt{t})$（热核！）
  \item 旋转不变 → 径向解 $u = u(r, t)$
\end{itemize}

\subsection{守恒律与矩映射}

\textbf{Noether 定理的李群表述：}

若作用量 $S = \int L\,dt$ 在一参数群 $g_\varepsilon$ 下不变，则：
\[
Q = \frac{\partial L}{\partial \dot{q}} \cdot \delta q
\]
是守恒量。

\textbf{矩映射（Moment Map）：} 将对称群的李代数映射到相空间上的守恒量：
\[
\mu: \mathfrak{g} \to C^\infty(M), \quad X \mapsto \mu_X
\]

满足 $\{\mu_X, f\} = X_M(f)$（$X_M$ 是 $X$ 在相空间 $M$ 上诱导的向量场）。

\textbf{例子：}
\begin{itemize}
  \item 平移对称 → 动量 $p$
  \item 旋转对称 → 角动量 $\mathbf{L} = \mathbf{r} \times \mathbf{p}$
  \item 时间平移 → 能量 $H$
\end{itemize}

%==============================================================
\section{在计算机图形学中的应用}
%==============================================================

\subsection{三维旋转与四元数}

\subsubsection{旋转的多种表示}

\begin{center}
\begin{tabular}{L{3cm}L{3cm}L{3cm}L{3cm}}
\toprule
\textbf{表示} & \textbf{参数数} & \textbf{奇异性} & \textbf{应用} \\
\midrule
旋转矩阵 $R \in SO(3)$ & 9（6 约束） & 无 & 通用变换 \\
欧拉角 & 3 & 有（万向锁） & 动画关键帧 \\
轴角 $(\hat{n}, \theta)$ & 3 & $\theta = 0$ 处 & 物理模拟 \\
四元数 $q \in SU(2)$ & 4（1 约束） & 无 & 插值、组合 \\
\bottomrule
\end{tabular}
\end{center}

\subsubsection{四元数与 $SU(2)$}

单位四元数 $q = \cos\frac{\theta}{2} + \sin\frac{\theta}{2}(\hat{n}_x i + \hat{n}_y j + \hat{n}_z k)$

\textbf{旋转作用：} $\mathbf{v}' = q\mathbf{v}q^{-1}$（将 $\mathbf{v}$ 视为纯虚四元数）

\textbf{优势：}
\begin{itemize}
  \item 无万向锁
  \item 组合旋转 = 四元数乘法（$O(1)$ vs 矩阵乘法 $O(27)$）
  \item 插值自然（Slerp）
  \item 归一化便宜（4 个分量）
\end{itemize}

\subsubsection{Slerp（球面线性插值）}

\[
\text{Slerp}(q_0, q_1, t) = \frac{\sin((1-t)\Omega)}{\sin\Omega}q_0 + \frac{\sin(t\Omega)}{\sin\Omega}q_1
\]

其中 $\cos\Omega = q_0 \cdot q_1$（四元数点积）。

\textbf{几何意义：} 在 $S^3$（单位四元数空间）上沿大圆弧插值。

\subsection{刚体运动与 $SE(3)$}

\subsubsection{刚体变换}

\[
T = \begin{pmatrix} R & \mathbf{t} \\ \mathbf{0}^T & 1 \end{pmatrix} \in SE(3)
\]

\subsubsection{$\mathfrak{se}(3)$ 与螺旋运动}

\[
\mathfrak{se}(3) = \left\{\begin{pmatrix} [\omega]_\times & \mathbf{v} \\ \mathbf{0}^T & 0 \end{pmatrix} : \omega, \mathbf{v} \in \mathbb{R}^3\right\}
\]

\textbf{指数映射（螺旋运动）：}
\[
\exp\begin{pmatrix} [\omega]_\times & \mathbf{v} \\ \mathbf{0}^T & 0 \end{pmatrix} = \begin{pmatrix} e^{[\omega]_\times} & V\mathbf{v} \\ \mathbf{0}^T & 1 \end{pmatrix}
\]

其中 $V = I + \frac{1-\cos\theta}{\theta^2}[\omega]_\times + \frac{\theta - \sin\theta}{\theta^3}[\omega]_\times^2$。

\textbf{在 CG 中：} 骨骼动画的关节变换、相机运动、机器人学。

\subsection{球谐函数与 $SO(3)$ 表示}

\subsubsection{球谐函数 = $SO(3)$ 的不可约表示基}

$Y_l^m(\theta, \phi)$ 是 $SO(3)$ 的第 $l$ 个不可约表示（维度 $2l+1$）的基函数。

\textbf{旋转下的变换：}
\[
Y_l^m(R^{-1}\hat{n}) = \sum_{m'=-l}^{l} D^l_{m'm}(R)\, Y_l^{m'}(\hat{n})
\]

$D^l_{m'm}(R)$ 是 Wigner $D$-矩阵（$(2l+1) \times (2l+1)$ 酉矩阵）。

\subsubsection{在渲染中的应用}

\textbf{环境光照：} 用球谐系数表示：
\[
L(\hat{n}) \approx \sum_{l=0}^{L}\sum_{m=-l}^{l} c_l^m Y_l^m(\hat{n})
\]

通常 $L = 2$（9 个系数）即可表示低频环境光。

\textbf{PRT（预计算辐射传输）：} 将传输算子在球谐基上对角化，实现实时全局光照。

\textbf{球面卷积：} 在球谐域变成逐元素乘法：
\[
\widehat{f * g}_l^m = \hat{f}_l \cdot \hat{g}_l^m
\]

（卷积定理的球面版本）

\subsection{等变神经网络}

\textbf{核心思想：} 网络层在群作用下等变：
\[
f(g \cdot x) = g \cdot f(x)
\]

\textbf{实现：} 用不可约表示的张量积构造等变层。

\textbf{应用：}
\begin{itemize}
  \item 分子性质预测（$SE(3)$ 等变）
  \item 流体模拟（旋转等变）
  \item 3D 点云处理
\end{itemize}

\subsection{微分几何与网格处理}

\subsubsection{曲面上的对称性}

\textbf{等距变换群：} 保持第一基本形式不变的变换。

\textbf{共形变换群：} 保持角度不变的变换（无穷维！）。

\textbf{在 CG 中：}
\begin{itemize}
  \item 纹理映射：共形映射（保持角度）
  \item 网格参数化：等距映射（保持距离）
  \item 形状匹配：在对称群下不变的特征
\end{itemize}

\subsubsection{离散微分几何中的群结构}

\textbf{离散联络：} 网格边上的旋转（$SO(3)$ 元素），描述向量沿边的平行移动。

\textbf{离散曲率：} 绕顶点一圈的旋转角度亏损 = 高斯曲率的离散版本。

\subsection{拓扑在 CG 中的应用}

\subsubsection{网格拓扑}

\textbf{Euler 公式：} $V - E + F = 2 - 2g$（$g$ 是亏格）

\begin{center}
\begin{tabular}{L{3cm}L{2cm}L{5cm}}
\toprule
\textbf{拓扑} & \textbf{亏格} & \textbf{例子} \\
\midrule
球面 & 0 & 封闭物体 \\
环面 & 1 & 甜甜圈、手柄 \\
双环面 & 2 & 眼镜框 \\
\bottomrule
\end{tabular}
\end{center}

\textbf{影响：}
\begin{itemize}
  \item 参数化：球面拓扑可以展开到平面（需要切割），环面可以无切割展开
  \item 向量场：球面上不存在处处非零的连续切向量场（Hairy Ball 定理）
  \item 网格生成：拓扑决定了网格的连接方式
\end{itemize}

\subsubsection{Hairy Ball 定理}

$S^2$ 上不存在处处非零的连续切向量场。

\textbf{在 CG 中：}
\begin{itemize}
  \item 球面上的纹理坐标必有奇点（至少一个极点）
  \item 风向/毛发方向场在球面上必有零点
  \item 这解释了为什么地球仪的地图投影必有畸变
\end{itemize}

\subsubsection{同调与网格分析}

\textbf{Betti 数：}
\begin{itemize}
  \item $b_0$：连通分量数
  \item $b_1$：独立环路数（“洞”的个数）
  \item $b_2$：封闭腔体数
\end{itemize}

\textbf{在 CG 中：} 网格质量检查、拓扑简化、特征提取。

%==============================================================
\section{总结：李群--李代数--拓扑的统一视角}
%==============================================================

\subsection{三条主线的交汇}

\[
\boxed{
\begin{aligned}
&\text{李群（全局对称性）} \xleftrightarrow{\text{指数映射}} \text{李代数（无穷小对称性）} \\
&\text{李群} \xleftrightarrow{\text{基本群/覆盖}} \text{拓扑（全局结构）} \\
&\text{李代数} \xleftrightarrow{\text{表示论}} \text{物理（粒子/量子数）}
\end{aligned}
}
\]

\subsection{核心对应表}

\begin{center}
\begin{tabular}{L{3cm}L{3cm}L{3cm}L{3cm}}
\toprule
\textbf{数学} & \textbf{物理} & \textbf{CG} & \textbf{微分方程} \\
\midrule
李群 & 对称群 & 变换群 & 对称性 \\
李代数 & 守恒量代数 & 无穷小变换 & 降阶生成元 \\
表示 & 粒子分类 & 球谐/傅里叶 & 特征函数 \\
指数映射 & 时间演化 & 旋转插值 & 矩阵指数解 \\
覆盖群 & 自旋 & 四元数 & 多值解 \\
特征类 & 拓扑荷 & 网格拓扑 & 指标定理 \\
\bottomrule
\end{tabular}
\end{center}

\subsection{学习路径建议}

\begin{enumerate}
  \item \textbf{入门：} $SO(2)$、$SO(3)$、四元数、旋转矩阵（CG 直接可用）
  \item \textbf{进阶：} 李代数、指数映射、BCH 公式（物理模拟、机器人学）
  \item \textbf{深入：} 表示论、球谐函数（渲染、量子力学）
  \item \textbf{抽象：} 纤维丛、特征类、同伦群（规范场论、拓扑物态）
\end{enumerate}

\subsection{推荐参考}

\begin{itemize}
  \item Hall, \textit{Lie Groups, Lie Algebras, and Representations}（数学系标准教材）
  \item Stillwell, \textit{Naive Lie Theory}（直观入门）
  \item Georgi, \textit{Lie Algebras in Particle Physics}（物理应用）
  \item Nakahara, \textit{Geometry, Topology and Physics}（拓扑与物理）
  \item Crane, \textit{Discrete Differential Geometry}（CG 中的离散几何）
  \item Sola et al., \textit{A micro Lie theory for state estimation in robotics}（工程应用）
\end{itemize}

% ---------------------------------------------------------------------------
% 专题三（跨篇串讲）已于 2026-11-02 拆为**独立册**：bridges/main.tex -> bridges/main.pdf
%
% 为什么只拆专题三、不拆专题一/二:
%   专题一（微分方程）、专题二（李群）是**新知识**，属于本册主线；
%   专题三是「回头贯通」的连接（见 bridges/README.md 第 3 行），
%   且会随学习进度持续追加，单册便于增量维护与单独打印。
%
% 拆分前提已验证：四篇不含 \ref/\eqref，也没有 equation/align/\tag，
%   公式全为无编号 \[ \]，独立编译不会出现 ?? 与编号漂移。
% ---------------------------------------------------------------------------

\end{document}